\documentclass[11pt,a4paper]{article}
\usepackage{fontspec}
\defaultfontfeatures{Ligatures=TeX}
\setmainfont{lmroman10-regular.otf}[BoldFont=lmroman10-bold.otf,ItalicFont=lmroman10-italic.otf,BoldItalicFont=lmroman10-bolditalic.otf]
\setsansfont{lmsans10-regular.otf}[BoldFont=lmsans10-bold.otf,ItalicFont=lmsans10-oblique.otf,BoldItalicFont=lmsans10-boldoblique.otf]
\setmonofont{lmmono10-regular.otf}[Scale=0.88,Ligatures=TeXOff]
\usepackage{amsmath,amssymb,amsthm,mathtools,mathrsfs}
\usepackage[margin=25mm,headheight=16pt]{geometry}
\usepackage{booktabs,array,longtable,enumitem,xcolor}
\usepackage{fancyhdr}
\usepackage{xurl}
\usepackage[unicode,colorlinks=true,linkcolor=blue!38!black,
 citecolor=blue!38!black,urlcolor=blue!38!black]{hyperref}
\usepackage{bookmark}
\hypersetup{pdftitle={Missing-digit primes in bases eight through eleven},
 pdfauthor={Tseng},pdfsubject={Analytic number theory; missing digits; computer-assisted proofs},
 pdfkeywords={missing-digit primes, Harman sieve, Fourier moments, affine correlation}}
\pdfstringdefDisableCommands{\def\leqslant{<=}\def\geqslant{>=}\def\asymp{asymp}\def\sim{sim}}
\newtheoremstyle{paperplain}{8pt}{8pt}{\normalfont}{}{\bfseries}{.}{.6em}{}
\theoremstyle{paperplain}
\newtheorem{theorem}{Theorem}[section]
\newtheorem{lemma}[theorem]{Lemma}
\newtheorem{proposition}[theorem]{Proposition}
\newtheorem{corollary}[theorem]{Corollary}
\newtheorem{definition}[theorem]{Definition}
\newtheorem{remark}[theorem]{Remark}
\numberwithin{equation}{section}
\newcommand{\e}{\mathrm e}
\newcommand{\ind}{\mathbf1}
\newcommand{\dd}{\,\mathrm d}
\setlist[enumerate]{leftmargin=2.2em,itemsep=4pt,topsep=4pt}
\setlist[itemize]{leftmargin=2.2em,itemsep=3pt,topsep=4pt}
\setlength{\parindent}{1.4em}
\setlength{\parskip}{2.5pt}
\setlength{\emergencystretch}{2em}
\makeatletter
\renewcommand*\l@subsection{\@dottedtocline{2}{1.8em}{3.4em}}
\makeatother
\linespread{1.04}
\allowdisplaybreaks[2]
\pagestyle{fancy}
\fancyhf{}
\fancyhead[L]{\small Missing-digit primes in bases eight through eleven}
\fancyhead[R]{\small Orders of magnitude and asymptotics}
\fancyfoot[C]{\thepage}
\renewcommand{\headrulewidth}{.3pt}
\title{\bfseries Missing-digit primes in bases eight through eleven}
\author{Tseng}
\date{}

\begin{document}
\maketitle
\begin{abstract}
Let $A_{b,d}(x)$ count the positive integers below $x$ whose ordinary
base-$b$ expansion omits $d$, and let $\pi_{b,d}(x)$ count the primes
among them. For every omitted digit and arbitrary real $x\to\infty$,
we prove
\[
 \begin{aligned}
 \pi_{b,d}(x)&\asymp \frac{A_{b,d}(x)}{\log x}
                 \asymp \frac{x^{\log_b(b-1)}}{\log x}
                      &&(b=8,9),\\[3pt]
 \pi_{b,d}(x)&\sim \frac{b h_{b,d}}{(b-1)\varphi(b)}
                       \frac{A_{b,d}(x)}{\log x}
                      &&(b=10,11),
 \end{aligned}
\]
where $h_{b,d}$ is the number of allowed last digits coprime to $b$.
In base eight, a centered divisor sum reduces the lower bound to a
comparison for grouped prime factors and a three-prime remainder
whose loss is at most $5/24$ of the reference term. In base nine, a
weighted estimate for shared rational denominators gives a Type II
range, and retaining the markings in complete prime factorizations
gives a sieve loss below $231/250$. In bases ten and eleven, an affine
correlation bound for reversed carries extends the Type II range
below the short-factor exponent $1/4$ and gives the asymptotic via the
Ford--Maynard sieve. The numerical inequalities used in the proofs
are verified by exact arithmetic.
\end{abstract}
\noindent\textbf{Keywords:} missing-digit primes; digital Fourier moments;
Type I/II estimates; Harman sieve; computer-assisted proofs.\\
\textbf{2020 Mathematics Subject Classification:} 11N05, 11N35, 11N36, 11A63.
\setcounter{tocdepth}{2}
\begingroup
\small
\linespread{.88}\selectfont
\setlength{\parskip}{0pt}
\tableofcontents
\endgroup
\setcounter{tocdepth}{2}
\newpage

\section{Introduction}
\label{sec:introduction}

We count primes that omit one prescribed digit in each of the bases
$8$, $9$, $10$, and $11$. The prime count has the expected order of
magnitude in bases eight and nine and the predicted local-density
asymptotic in bases ten and eleven.

For a base $b\geqslant3$ and a digit $d\in\{0,\ldots,b-1\}$, let
\[
 \mathscr A_{b,d}
 =\{n\geqslant1:\text{the ordinary base-}b
                    \text{ expansion of }n\text{ omits }d\},
\]
and write
\[
 A_{b,d}(x)=\#(\mathscr A_{b,d}\cap[1,x)),\qquad
 \pi_{b,d}(x)=\#\{p<x:p\text{ is prime and }p\in\mathscr A_{b,d}\}.
\]
Ordinary expansions have no leading zeros, and exactly one digit is
excluded at a time. In base eleven, $d=10$ denotes the eleventh digit
symbol. Put
\begin{equation}\label{intro:constants}
 c_b=\frac{\log(b-1)}{\log b},\qquad
 h_{b,d}=\varphi(b)-\mathbf1_{(d,b)=1},\qquad
 \kappa_{b,d}=\frac{b\,h_{b,d}}{(b-1)\varphi(b)}.
\end{equation}
Here $h_{b,d}$ counts the allowed last digits coprime to $b$, and
$\varphi$ is Euler's totient function. Logarithms without an indicated
base are natural.

\begin{theorem}[Orders of magnitude in bases eight and nine]
\label{main:order}
For $b\in\{8,9\}$ and every $d\in\{0,\ldots,b-1\}$, one has
\[
 \pi_{b,d}(x)\asymp\frac{A_{b,d}(x)}{\log x}
 \asymp\frac{x^{c_b}}{\log x}
 \qquad(x\to\infty)
\]
through arbitrary real values.
\end{theorem}

\begin{theorem}[Exact asymptotics in bases ten and eleven]
\label{main:exact}
For $b\in\{10,11\}$ and every $d\in\{0,\ldots,b-1\}$,
\begin{equation}\label{main:asymptotic}
 \pi_{b,d}(x)\sim
 \kappa_{b,d}\frac{A_{b,d}(x)}{\log x}
 \qquad(x\to\infty),
\end{equation}
where $x$ tends to infinity through arbitrary real values.
\end{theorem}

Theorems~\ref{octal:prime-count} and \ref{nine:theorem} give explicit
lower and upper constants for every digit of their respective bases.

The local factor in \eqref{intro:constants} compares the proportion
$h_{b,d}/(b-1)$ of allowed unit last digits with the proportion
$\varphi(b)/b$ of units among all integers. The use of $A_{b,d}(x)$
also matters. Although $A_{b,d}(x)\asymp x^{c_b}$, this counting
function has persistent fluctuations. For example, choose an allowed
nonzero digit $a$ and put $C=ab+d$. Every integer in
$[Cb^k,(C+1)b^k)$ begins with the forbidden prefix $a,d$, so
\[
 A_{b,d}(Cb^k)=A_{b,d}((C+1)b^k).
\]
The corresponding powers of the endpoints have the fixed ratio
$((C+1)/C)^{c_b}>1$. Thus the main term in
\eqref{main:asymptotic} cannot in general be replaced by a fixed
constant times $x^{c_b}/\log x$.

\subsection{Relation to previous work}

Maynard \cite[Theorem~1.1]{Maynard2019} proved the decimal
order-of-magnitude result. Equation~(1.1) of his paper gives the
local factor in the conjectural decimal asymptotic, and the discussion
following his Theorem~1.2 records the one-digit asymptotic for
$b\geqslant12$. Karwatowski \cite[Theorem~1]{Karwatowski2022} obtained
two-sided estimates at powers of $b$ for every $b\geqslant10$.
His base-nine preprint \cite[Theorem~1]{KarwatowskiBaseNine} treats
the omissions $d=0,8$. Theorem~\ref{main:order} gives the octal
estimate for every digit and extends the base-nine estimate to
$d=1,\ldots,7$. Theorem~\ref{main:exact} gives the decimal asymptotic
and the corresponding formula in base eleven.

The proofs use the digital circle method and the analysis of
exceptional frequency pairs developed by Maynard
\cite[Sections~10 and~13--15]{Maynard2019}. Their further analytic and
sieve inputs differ by base, as follows.

\subsection{Proofs of the main results}\label{sec:outline}

\paragraph{Base eight.}
The common distribution estimate in arithmetic progressions extends
to level $X^\theta$ for each $\theta<.6515$ by a quadratic Fourier
approximation. It supplies the Selberg upper bound and the divisor
estimates for a centered M\"obius sum. Complementary divisors cancel on
even prime-factor ranks. On odd ranks, grouped comparisons handle
products with a proper divisor in $[n^{2/5},n^{3/5}]$.
The remaining contribution is bounded on a three-prime simplex.
Sections~\ref{sec:octal-typeI} and \ref{sec:octal-prime-count} give
the reduction to the grouped comparison. Sections~\ref{sec:octal-fourier}
and \ref{sec:octal-balanced-proof} prove that comparison using dyadic
characters at low conductors and short row differences at high conductors.

\paragraph{Base nine.}
The sieve uses Maynard's Buchstab decomposition
\cite[Section~6]{Maynard2019} within Harman's framework \cite{Harman2007}.
Sections~\ref{geo:section} and \ref{geo:typeii-section} control the
multiplicities of shared denominators and supply the Type II range.
Sections~\ref{arith:section} and \ref{sieve:section} transfer this
estimate to marked prime-factor sums. Factoring each rough cofactor
exposes additional admissible subproducts. Section~\ref{nine:sieve-section}
bounds the loss on the resulting ordered simplices with each marking
counted at its full multiplicity.

\paragraph{Bases ten and eleven.}
The starting moment exponents, at orders $1$ and $235/154$, are
Maynard's \cite[Lemmas~10.3--10.4]{Maynard2019}. Karwatowski
\cite[Sections~3 and~5]{Karwatowski2022} gives the corresponding
eigenvalue criteria and reports a Mathematica verification for base
eleven. The additional affine correlation bound and reversed digit
decomposition in Sections~\ref{dec:section} and
\ref{dec:bilinear-section} extend the Type II range below $1/4$.
Section~\ref{dec:prime-section} constructs a common comparison
sequence for Type I and Type II and applies the asymptotic sieve of
Ford and Maynard \cite[Theorem~2.2(a)]{FordMaynard2024}.

Section~\ref{sec:preliminaries} supplies the digital notation and
the passage from fixed prefixes to real endpoints. The appendices
give the finite verification criteria and identify the accompanying
certificates. Decimal thresholds in inequalities are exact rationals.

\section{Digital words, local factors, and Fourier products}
\label{sec:preliminaries}

Digital moment estimates are expressed in terms of suffix words of a fixed
length. Their product structure governs the Fourier analysis, while the
allowed unit last digits determine the arithmetic density.

\subsection{Ordinary expansions and fixed-length words}

Fix $b,d$, write $D=\{0,\ldots,b-1\}\setminus\{d\}$, and set
\[
 \mathcal W_k=\left\{\sum_{j=0}^{k-1}\varepsilon_jb^j:
                 \varepsilon_j\in D\right\},\qquad X=b^k,
 \qquad A=(b-1)^k=X^{c_b}.
\]
The $k$ positions here define a suffix word. If $d\ne0$, this set includes
zero and integers with shorter ordinary expansions; if $d=0$, it contains
only words of actual length $k$. For a fixed allowed positive prefix $v$,
the set $v b^k+\mathcal W_k$ always consists of ordinary expansions.

\begin{lemma}[Digital counting]\label{prelim:count}
For fixed $b,d$, one has $A_{b,d}(x)\asymp_b x^{c_b}$. More precisely,
\[
 A_{b,d}(b^k)=
 \begin{cases}
 \displaystyle\frac{(b-1)((b-1)^k-1)}{b-2},&d=0,\\[3pt]
 (b-1)^k-1,&d\ne0.
 \end{cases}
\]
Every interval $I$ of length $Y\geqslant1$ satisfies
\begin{equation}\label{prelim:interval}
 \#(\mathscr A_{b,d}\cap I)\ll_b Y^{c_b}.
\end{equation}
The same estimate holds for fixed-length words and arbitrary translates.
\end{lemma}
\begin{proof}
If $0$ is forbidden, there are $(b-1)^j$ words of actual length $j$; if a
nonzero digit is forbidden, there are $(b-2)(b-1)^{j-1}$ such words.
Summation gives the endpoint formulas, and comparison with adjacent powers
of $b$ proves the first assertion. Choose
$b^{\ell-1}\leqslant Y<b^\ell$. The interval $I$ meets at most two standard
digital cylinders of length $b^\ell$. A cylinder with a positive prefix
contains at most $(b-1)^\ell$ allowed suffixes. The first cylinder, which
contains zero, has $O_b((b-1)^\ell)$ allowed ordinary expansions by the
endpoint formula just proved. Adding these bounds gives the result.
For fixed-length words the same cylinder bound holds when $\ell\leqslant k$;
when $\ell>k$, their total size already gives the assertion. Translation
of the interval proves the translated version.
\end{proof}

\begin{lemma}[The last-digit local density]\label{prelim:local}
Among the allowed words of length $k\geqslant1$, exactly
$h_{b,d}(b-1)^{k-1}$ are coprime to $b$. For ordinary expansions,
\[
 \#\{n<x:n\in\mathscr A_{b,d},\ (n,b)=1\}
   =\frac{h_{b,d}}{b-1}A_{b,d}(x)+O_b(1).
\]
Relative to the unit density $\varphi(b)/b$ among ordinary integers,
the corresponding factor is $\kappa_{b,d}$.
\end{lemma}
\begin{proof}
There are $h_{b,d}$ allowed unit last digits, and all other positions are
free, proving the fixed-length formula. For ordinary expansions, group
integers by the positive prefix $\lfloor n/b\rfloor$. Each complete allowed
group has $b-1$ members and exactly $h_{b,d}$ units. There is at most one
incomplete group at the endpoint; all single-digit integers contribute
$O_b(1)$. Thus the discrepancy is bounded independently of the number of
digit positions. Dividing by the unit density $\varphi(b)/b$ gives the factor.
\end{proof}

\subsection{Fourier products and moments}

Write $\mathrm e(t)=\exp(2\pi i t)$ and
$\|t\|=\min_{n\in\mathbb Z}|t-n|$, and define
\[
 P_{b,d}(z)=\sum_{a\in D}z^a,\qquad
 f_{b,d}(t)=\frac{|P_{b,d}(\mathrm e(t))|}{b-1},
\]
\[
 F_X(t)=\frac1A\left|\sum_{n\in\mathcal W_k}\mathrm e(nt)\right|
       =\prod_{j=0}^{k-1}f_{b,d}(b^jt).
\]
In particular, $0\leqslant F_X\leqslant1$ and
$f_{b,b-1-d}=f_{b,d}$, since
$P_{b,b-1-d}(z)=z^{b-1}P_{b,d}(z^{-1})$. The terminal-digit density is
given by Lemma~\ref{prelim:local}.

For $p>0$, let
\[
 H_p(X;x)=\sum_{a=0}^{X-1}F_X((a+x)/X)^p,
 \qquad
 (T_ph)(x)=\sum_{j=0}^{b-1}f_{b,d}((j+x)/b)^p h((j+x)/b).
\]
Expansion one digit at a time gives the exact identity
\begin{equation}\label{prelim:transfer}
 T_p^k1(x)=H_p(b^k;x).
\end{equation}
Some later certificates use the normalization $\mathcal L_p=b^{-1}T_p$;
their supersolution constants are correspondingly divided by $b$.

\begin{lemma}[Iteration of a positive supersolution]\label{prelim:supersolution}
If a continuous positive $1$-periodic function $h$ satisfies
$T_ph\leqslant\Lambda h$, then for every $k\geqslant0$ and $x\in\mathbb R$,
\[
 H_p(b^k;x)\leqslant\frac{\max h}{\min h}\Lambda^k.
\]
If $\Lambda\leqslant b^g$, then $H_p(X;x)\ll X^g$ and
$\int_0^1F_X(t)^p\,dt\ll X^{g-1}$.
\end{lemma}
\begin{proof}
Iterate using $1\leqslant h/\min h$ and positivity of the operator.
Integrating \eqref{prelim:transfer} over $0\leqslant x\leqslant1$ and
changing variables gives $X\int_0^1F_X(t)^p\,dt$ on the left.
\end{proof}

The following two even moments need no computational certificate:
\begin{equation}\label{prelim:even}
 H_2(X;x)=X^{1-c_b},\qquad
 H_4(X;x)\leqslant
 X^{\log_b(1+(b-1)^{-3})}.
\end{equation}
The first identity is discrete Parseval. For the second, the one-digit
polynomial $f_{b,d}^4$ has nonnegative Fourier coefficients, so its sum over
$b$ equally spaced sample points is largest at $x=0$. At each nontrivial
$b$th root of unity, $|P_{b,d}|=1$, and this maximum is therefore
$1+(b-1)^{-3}$. Iteration of the positive transfer operator proves the bound.

The continuous inequalities needed for the other moment bounds are obtained
from rational values on finite grids. Appendix~\ref{app:continuous}
gives the interpolation and curvature estimates that justify this passage
from nodes to the whole circle.

\subsection{From fixed prefixes to arbitrary endpoints}

A prime asymptotic for each fixed admissible prefix determines the answer
at arbitrary real endpoints. The order of limits in the following argument
requires uniformity only over finitely many prefixes at a time.

\begin{lemma}[Completion from fixed-prefix cylinders]\label{prelim:endpoints}
Suppose that a constant $\kappa>0$ has the following property: for every
fixed admissible positive prefix $v$,
\begin{equation}\label{prelim:prefix-hypothesis}
 \#\{p\text{ prime}:p\in vb^k+\mathcal W_k\}
 =(\kappa+o_v(1))\frac{(b-1)^k}{k\log b}
 \qquad(k\to\infty).
\end{equation}
Then, with ordinary expansions and arbitrary real $x\to\infty$,
\[
 \pi_{b,d}(x)\sim\kappa\,\frac{A_{b,d}(x)}{\log x}.
\]
\end{lemma}
\begin{proof}
Apply \eqref{prelim:prefix-hypothesis} to the finitely many allowed
nonzero one-digit prefixes and sum over suffix lengths. Geometric
summation gives
\begin{equation}\label{prelim:initial-prime-bound}
 \pi_{b,d}(b^s)\ll_{b,d,\kappa}\frac{(b-1)^s}{s}.
\end{equation}
The finitely many small suffix lengths are absorbed in the constant.

Take $b^K\leqslant x<b^{K+1}$, fix $L_0\geqslant2$, and set
$k=K-L_0+1$, $q=b^k$, $V=\lfloor x/q\rfloor<b^{L_0}$.
Partition $[1,x)$ into $[1,q)$, the complete cylinders
$[vq,(v+1)q)$ for $1\leqslant v<V$, and the terminal partial cylinder.
For fixed $L_0$, the hypothesis is uniform over this finite collection
of prefixes. An inadmissible prefix contributes nothing. In every complete
admissible cylinder, $\log x/(k\log b)\to1$ and there are exactly
$(b-1)^k$ allowed integers, so their total discrepancy is
$o_{L_0}((b-1)^K)$.

A terminal admissible partial cylinder has at most $(b-1)^k$ allowed
integers and at most the prime mass of its complete cylinder.
The initial segment is bounded by
\eqref{prelim:initial-prime-bound} and Lemma~\ref{prelim:count}.
Thus
\[
 \bigl|\pi_{b,d}(x)\log x-\kappa A_{b,d}(x)\bigr|
 \leqslant o_{L_0}((b-1)^K)
       +O_{b,d,\kappa}((b-1)^{K-L_0+1}).
\]
The constant in the second term is independent of $L_0$: for each fixed
depth the complete-cylinder prime bound is eventually at most twice
its main term, and $K\geqslant2L_0$ ensures
$\log x/(k\log b)\leqslant2$. Since
$A_{b,d}(x)\geqslant(b-2)(b-1)^{K-1}$, let first $K\to\infty$ and
then $L_0\to\infty$.
\end{proof}


\part{Orders of magnitude in bases eight and nine}

\section{A weighted estimate for common denominators}\label{geo:section}

For the base-nine Type II estimate, we bound pairs of frequencies
whose rational approximations have a shared denominator. Such pairs
arise from Maynard's geometric analysis of large bilinear kernels
\cite[Lemmas~13.2 and~14.1]{Maynard2019}. We first bound the weighted
second moment of the rational incidences, counting each shared
denominator with its multiplicity. Combining this estimate with
digital row bounds gives the minor-arc input to
Theorem~\ref{geo:typeii}.

Fix a prime-power base $b=p^r$ and a forbidden digit $d$, and put
\[
 \mathcal D=\{0,\ldots,b-1\}\setminus\{d\},\qquad
 c=c_b=\frac{\log(b-1)}{\log b},\qquad Z=b^j.
\]
Let $\mathcal W(Z)$ be the set of all allowed words of length $j$, and write
$P_Z(t)=Z^{-c}\sum_{n\in\mathcal W(Z)}\mathrm e(nt)$ and $F_Z=|P_Z|$.
Thus $0\le F_Z\le1$, and for powers $U,V$ of the base,
\begin{equation}\label{geo:product}
 F_{UV}(t)=F_U(t)F_V(Ut),\qquad
 \int_0^1 F_U(t)^2\,dt=U^{-c}.
\end{equation}
For the moment, suppose that the following estimates hold uniformly over all
powers of the base and all real shifts $x$:
\begin{equation}\label{geo:moments}
 \sum_{a<Z}F_Z\left(\frac{a+x}{Z}\right)\ll Z^{1-\sigma},
 \qquad \int_0^1F_Z(t)\,dt\ll Z^{-\sigma},
 \qquad \sum_{a<Z}F_Z(a/Z)^{4/3}\ll Z^{g_{4/3}},
\end{equation}
where $1/2<\sigma<c$. Proposition~\ref{nine:fourier} supplies these
moment bounds for the base-nine application.

There are two complementary scales. When the rational windows are sparse,
direct row estimates retain their decay in the denominator and window width.
When they are dense, four-cycles control the off-diagonal part of an incidence
matrix, producing the $4/3$ moment in the weighted estimate. Keeping these
scales separate is what allows the final bound to save arbitrary powers of
the logarithm even for short denominators. All implicit constants may depend
on the fixed base, alphabet, moment exponents, and stated fixed parameters,
but never on the word length or the denominators being summed.

\subsection{The weighted energy estimate}

For a fixed denominator, the row mass below is the total frequency weight
near its reduced rational points. The bilinear argument requires the sum
of the squares of these masses, so pairs sharing many denominators must
be counted with their full weight.

For $D\ge1,\Delta\ge1$ and fixed $C>0$, take integer rows
$D\le q<2D$ and columns $a\in\mathbb Z/Z\mathbb Z$, and define
\begin{equation}\label{geo:incidence-matrix}
 M_{q,a}=\mathbf1_{\{\exists h\pmod q:\ (h,q)=1,
       \ \|a/Z-h/q\|\le C\Delta/Z\}},\qquad D\le q<2D.
\end{equation}
The row and column sets may be restricted arbitrarily in what follows.
Write $H=D^2\Delta/Z$.
We shall first treat the range
\begin{equation}\label{geo:det-hyp}
 H\ge1,\qquad 4CD\Delta<Z,\qquad
 48C^2D^3\Delta^2<Z^2.
\end{equation}

\begin{proposition}[Weighted common-denominator second moment]\label{geo:energy}
Assume~\eqref{geo:det-hyp}, let $0\le w_a\le1$, and put
\[
 H_u=\sum_aw_a^u,\qquad S_1=\max_q\sum_aw_aM_{q,a}.
\]
Then, for every $\varepsilon>0$,
\begin{equation}\label{geo:energy-bound}
 \sum_q\left(\sum_aw_aM_{q,a}\right)^2
 \ll_\varepsilon Z^\varepsilon\left\{
 S_1H_1+\frac{\Delta^{3/2}D^2}{Z}H_{4/3}^{3/2}
 +\frac{\Delta^2D^3}{Z^2}H_1^2\right\}+O(DZ^{-18}).
\end{equation}
\end{proposition}

The first term is the diagonal contribution of the row Gram matrix.
The other two terms control its off-diagonal part through four-cycles. For a set
of $R$ columns, the first four-cycle contribution has size proportional
to $R^{3/2}$. A weight layer of height $\mu$ therefore contributes
\[
 \mu^2R^{3/2}=(\mu^{4/3}R)^{3/2}.
\]
This is the source of the $4/3$ moment in~\eqref{geo:energy-bound}.
We prove the estimate by counting short affine relations and then showing
that many shared denominators force an especially short relation.

\subsection{Short affine relations}

\begin{lemma}[Counting pairs admitting an affine relation]\label{geo:incidence}
Let $E\subset\{0,\ldots,Z-1\}$, $R=|E|$, and $1\le V\le W$.
Let $I(E;V,W)$ count the ordered pairs $(a,a')\in E^2$ for which there are
integers $v_i$ satisfying
\[
 v_1a+v_2a'+v_3Z+v_4=0,\qquad (v_1,v_2,v_3)\ne(0,0,0),
\]
with $|v_1|,|v_2|,|v_3|\le V$ and $|v_4|\le W$.
Then, for every $\varepsilon>0$,
\begin{equation}\label{geo:incidence-bound}
 I(E;V,W)\ll_\varepsilon
 Z^\varepsilon\left(RVW+\frac{R^2V^2W}{Z}\right).
\end{equation}
\end{lemma}

\begin{proof}
If $V^2W\ge Z$, the second term already bounds $R^2$, so we may assume that
$V^2W<Z$. In particular, $W<Z$. If $v_1=0$ and $v_2\ne0$, then
$v_2a'=-v_3Z-v_4$.
There are $O(VW)$ choices for the right-hand side. When it is nonzero, the
divisor bound gives at most $O_\varepsilon(Z^\varepsilon)$ factorizations;
when it vanishes, we must have $v_3=v_4=0$ and $a'=0$.
Multiplying by the $R$ choices for $a$ gives
$O_\varepsilon(RVWZ^\varepsilon)$. The case $v_2=0$ is identical,
and these two coefficients cannot both vanish.

For the remaining relations, interchange $a,a'$ and change signs as needed,
then assign each pair a witness on a dyadic scale $S\le V$ satisfying
\[
 S\le v_2<2S,\qquad 0<|v_1|\le2S.
\]
Let $I_S$ count the pairs assigned to this scale.
Set $J=\lfloor SZ/(100W)\rfloor$.
If $SZ/W$ is bounded, then $S^2W/Z\gg1$ and the desired bound is immediate.
Otherwise, $1\le J<Z$ and $J\asymp SZ/W$. Use the Fej\'er kernel of height one,
\[
 B_J(t)=J^{-2}\left|\sum_{0\le m<J}\mathrm e(mt)\right|^2,
 \qquad f(t)=\sum_{a\in E}B_J(t-a/Z).
\]
Its Fourier support lies in $|n|<J$, and orthogonality on the full grid of
$Z$ points gives
\begin{equation}\label{geo:fejer-norms}
 \int f=R/J,\qquad \|f\|_\infty\le Z/J,
 \qquad \|f\|_2^2\le ZR/J^2.
\end{equation}
Consider the periodic operator
\begin{equation}\label{geo:periodic-operator}
 (T_Sf)(t)=
 \sum_{0<|m_1|\le2S}\ \sum_{S\le m_2<2S}
 \ \sum_{j=0}^{m_2-1}f\left(\frac{m_1t+j}{m_2}\right).
\end{equation}
For a given relation, take $m_1=-v_1,m_2=v_2,j\equiv-v_3\pmod {m_2}$.
At $t=a/Z$, the image lies within circular distance $W/(SZ)$ of $a'/Z$.
Since $J\le SZ/(100W)$, the corresponding two Fej\'er atoms are both bounded
below by an absolute positive constant on an interval of length $\gg1/J$.
Expanding $fT_Sf$ and retaining one witness for each ordered pair therefore gives
\begin{equation}\label{geo:positive-witness}
 I_S/J\ll\int_0^1 f(t)(T_Sf)(t)\,dt.
\end{equation}
The short intervals need not be disjoint: distinct pairs carry distinct
pairs of atom labels in the expansion, and all terms are nonnegative.
At the circular seam one may use a one-sided interval; the complete sum over
$j\bmod m_2$ ensures that the operator remains periodic.

For an input frequency $n$, summing over the intercept gives
\[
 \sum_{j\bmod m_2}\mathrm e\left(n\frac{m_1t+j}{m_2}\right)
 =\begin{cases}m_2\mathrm e(\ell m_1t),&n=\ell m_2,\\0,&m_2\nmid n.
 \end{cases}
\]
Thus the nonzero-frequency matrix is generated by
$n=\ell m_2$, $k=\ell m_1$, $\ell\ne0$, with weight $m_2$ for each contribution.
For fixed $n$ we have $m_2\mid n$, and for fixed $k$ we have $m_1\mid k$.
The other multiplier has $O(S)$ choices, each with weight $O(S)$.
Inputs with $|n|<J$ produce outputs with $|k|<2J$, so both row sums and column
sums are $O_\varepsilon(S^2Z^\varepsilon)$. Schur's test yields
\[
 \|T_S^{\ne0}\|_{2\to2}\ll_\varepsilon S^2Z^\varepsilon.
\]
The constant mode has eigenvalue $O(S^3)$, and, because $m_1\ne0$, no nonzero
input frequency contributes to the constant mode. Combining
\eqref{geo:fejer-norms} and~\eqref{geo:positive-witness}, we obtain
\[
 I_S\ll_\varepsilon
 J\left\{S^3(R/J)^2+S^2Z^\varepsilon ZR/J^2\right\}
 \ll_\varepsilon RSWZ^\varepsilon+R^2S^2W/Z.
\]
Summing the geometric series over $S$ and including the cases with a zero
coordinate proves the lemma.
\end{proof}

\subsection{Shared denominators and four-cycles}


For two distinct columns, their codegree is the number of denominators shared
by the two frequencies. If that number is $t$, the integer lifts of their
rational approximations lie in one plane, and the product of its successive
minima is $O(1/t)$. Their cross product consequently gives an affine relation
whose height decreases by the same factor. Summing this gain with weight
$t(t-1)$ gives the four-cycle estimate needed for the weighted energy.

\begin{lemma}[Anisotropic codegree bound]\label{geo:codegree}
Assume~\eqref{geo:det-hyp}.
For any $E\subset\{0,\ldots,Z-1\}$, let $R=|E|$ and
$t_{a,a'}=\sum_qM_{q,a}M_{q,a'}$.
Then
\begin{equation}\label{geo:fourcycles}
 \sum_{\substack{a,a'\in E\\a\ne a'}}t_{a,a'}(t_{a,a'}-1)
 \ll_\varepsilon Z^\varepsilon
 \left(R\Delta H^2+\frac{R^2\Delta H^3}{Z}\right).
\end{equation}
\end{lemma}

\begin{proof}
Fix a pair with $t=t_{a,a'}\ge2$ common denominators.
Each denominator determines a unique integer lift $(h_1,h_2,q)$ such that
\[
 e_1=h_1-qa/Z,\quad e_2=h_2-qa'/Z,
 \quad |q|\le2D,\quad |e_i|\le2CD\Delta/Z.
\]
The transformation $(h_1,h_2,q)\mapsto(q,e_1,e_2)$ has determinant of absolute
value one. By~\eqref{geo:det-hyp}, the integer determinant of any three lifted
vectors has absolute value less than one, so the vectors are coplanar.
Lifts with distinct positive denominators cannot be proportional, since both
fractions are reduced. They therefore span a lattice of rank two in this plane.

Let $\lambda_1\le\lambda_2\le1$ be the successive minima of this section lattice
with respect to the symmetric box above. The two-dimensional lattice-point
bound gives $t\ll(\lambda_1\lambda_2)^{-1}$, and hence
$\lambda_1\lambda_2\ll1/t$.
To see this counting bound, transform the box linearly to a unit cube and take
a reduced basis in the resulting plane. The two basis coordinates have
$O(1+\lambda_i^{-1})$ possible values.
Choose corresponding independent integer vectors $w,w'$, and set
$(v_1,v_2,v_3)=w\times w'$.
Writing $w=(h_1,h_2,q)$ and $w'=(h'_1,h'_2,q')$, we have
\[
 v_1=e_2q'-qe'_2,\qquad v_2=qe'_1-e_1q',
\]
and
\[
 v_1a+v_2a'+v_3Z=Z(e_1e'_2-e_2e'_1).
\]
The left-hand side is an integer. Let $v_4$ be its negative; then
\begin{equation}\label{geo:height-codegree}
 |v_1|,|v_2|,|v_3|\ll_C H/t,\qquad
 |v_4|\ll_C\Delta H/t,
 \qquad(v_1,v_2,v_3)\ne0.
\end{equation}
The product of error terms in $v_3$ is absorbed using $\Delta/Z<1$.
This also gives $t\ll_CH$: if the first height bound is smaller than one,
there is no such integer configuration.

On a dyadic range $T\le t<2T$, apply Lemma~\ref{geo:incidence} with
$V\asymp H/T$ and $W\asymp\Delta H/T$. This gives
\[
 \#\{(a,a')\in E^2:t_{a,a'}\ge T\}
 \ll_\varepsilon Z^\varepsilon
 \left(\frac{R\Delta H^2}{T^2}
       +\frac{R^2\Delta H^3}{ZT^3}\right).
\]
Multiply by $T^2$ and sum over dyadic scales. The extra logarithm in the first
term is absorbed into $Z^\varepsilon$, while the second term is a geometric
series. This proves~\eqref{geo:fourcycles}.
\end{proof}

\begin{proof}[Proof of Proposition~\ref{geo:energy}]
First take $w=\mathbf1_E$, and denote the row degrees by $\ell_q$ and the column
degrees by $d_a$. Farey separation gives $d_a\ll1+H\ll H$.
For the actual $0$--$1$ incidence matrix $M$, we have the exact identity
\begin{equation}\label{geo:trace-identity}
 \sum_{q\ne q'}(MM^*)_{q,q'}^2
 =\sum_a d_a(d_a-1)
  +\sum_{a\ne a'}t_{a,a'}(t_{a,a'}-1).
\end{equation}
Split $MM^*$ into its diagonal and off-diagonal parts. The former has operator
norm $\max_q\ell_q$, while the operator norm of the latter is at most its
Hilbert--Schmidt norm. Using
$\|M\mathbf1_E\|_2^2\le R\|MM^*\|$ and Lemma~\ref{geo:codegree}, we obtain
\begin{equation}\label{geo:unweighted-energy}
 \sum_q\ell_q^2\ll_\varepsilon Z^\varepsilon
 \left\{R\max_q\ell_q+R^{3/2}\Delta^{1/2}H
             +R^2\Delta^{1/2}H^{3/2}Z^{-1/2}\right\}.
\end{equation}
The repeated-column contribution
$\sum_a d_a(d_a-1)\ll RH^2$ is absorbed by the first four-cycle term,
since $\Delta\ge1$.

For general weights, decompose into layers
$E_\mu=\{a:\mu/2<w_a\le\mu\}$, with $\mu$ dyadic, and truncate at
$w_a\ge Z^{-10}$. There are $O(\log Z)$ layers. Write
$R_\mu=|E_\mu|$ and $\ell_{q,\mu}=\sum_{a\in E_\mu}M_{q,a}$.
Cauchy's inequality in the layer index gives
\[
 \sum_q\left(\sum_\mu\mu\ell_{q,\mu}\right)^2
 \ll(\log Z)\sum_\mu\mu^2\sum_q\ell_{q,\mu}^2.
\]
In the first term of~\eqref{geo:unweighted-energy}, use
$\mu\ell_{q,\mu}\le2S_1$ and
$\sum_\mu\mu R_\mu\le2H_1$. For the other two terms, use
\[
 \sum_\mu\mu^2R_\mu^{3/2}
 \le\left(\sum_\mu\mu^{4/3}R_\mu\right)^{3/2}
 \le4H_{4/3}^{3/2},\qquad
 \sum_\mu\mu^2R_\mu^2\le4H_1^2.
\]
This exhibits the origin of the $4/3$ moment directly.
Each row sum of the omitted tail is at most $Z^{-9}$, so its squared norm is
$O(DZ^{-18})$; the cross term is absorbed by $(u+v)^2\le2u^2+2v^2$.
Substitute $H=D^2\Delta/Z$ and absorb the fixed logarithmic loss by increasing
$\varepsilon$ slightly.
\end{proof}

\subsection{Digital row bounds on rational grids}

\begin{lemma}[Positive envelopes and separated sampling]\label{geo:sampling}
Let $Q$ be a trigonometric polynomial of degree at most $U\ge1$,
let $j=1,2$, and fix $C_0>0$. There is a nonnegative trigonometric
polynomial $E$ such that
\begin{equation}\label{geo:envelope}
 \sup_{|v|\le C_0/U}|Q(t+v)|^j\le E(t),\qquad
 \deg E=O_{C_0}(U),\qquad \int E\ll_{C_0}\int|Q|^j.
\end{equation}
If $\Theta\subset\mathbb R/\mathbb Z$ has circular spacing at least
$0<\delta\le1$, then
\begin{equation}\label{geo:separated-sampling}
 \sum_{t\in\Theta}\sup_{|v|\le C_0/U}|Q(t+v)|^j
 \ll_{C_0}(U+\delta^{-1})\int_0^1|Q(t)|^j\,dt.
\end{equation}
Moreover, every nonnegative trigonometric polynomial of degree at most
$U_0$ satisfies
\begin{equation}\label{geo:grid-positive}
 \sum_{h\bmod q}E(h/q)
 =q\sum_{\ell\in\mathbb Z}\widehat E(q\ell)
 \le(q+2U_0)\widehat E(0).
\end{equation}
\end{lemma}

\begin{proof}
Choose a smooth Fourier cutoff equal to one on $[-U,U]$ and supported
in $[-2U,2U]$. Its periodic kernel $\Psi_U$ reproduces $Q$ and satisfies
$|\Psi_U(t)|\ll_A U(1+U\|t\|)^{-2A-2}$ for every fixed $A$.
For a sufficiently large fixed integer $A$, use the positive kernel
\[
 K_U(t)=U\sum_{n\in\mathbb Z}\left\{
 \left(\frac{\sin\pi U(t+n)}{\pi U(t+n)}\right)^{2A}
 +\left(\frac{\sin\pi(U(t+n)-1/2)}{\pi(U(t+n)-1/2)}\right)^{2A}
 \right\}.
\]
It has degree $O_A(U)$ and bounded integral. The two real-line summands
have no common zero and their sum is comparable to
$(1+|Ut|)^{-2A}$. Consequently,
\[
 \sup_{|v|\le C_0/U}|\Psi_U(t+v)|\ll_{A,C_0}K_U(t).
\]
The integral form of H\"older's inequality, applied to
$Q=\Psi_U*Q$, now gives~\eqref{geo:envelope} with
$E=C_{A,C_0}K_U*|Q|^j$.
For fixed $x$, separation and the decay of $K_U$ give
\[
 \sum_{t\in\Theta}K_U(t-x)
 \ll U\sum_{m\ge0}(1+Um\delta)^{-2A}
 \ll U+\delta^{-1}.
\]
Integrating this bound against $|Q(x)|^j$ proves
\eqref{geo:separated-sampling}.
Finally, orthogonality on a complete grid gives the equality in
\eqref{geo:grid-positive}; its inequality follows from
$|\widehat E(n)|\le\widehat E(0)$ and the Fourier support of $E$.
\end{proof}

\begin{lemma}[Coprime row sums and window averages]\label{geo:rows}
Define
\[
 A_q^Z(\Delta)=\sum_{a<Z}F_Z(a/Z)M_{q,a}\qquad(q\ge1).
\]
Under~\eqref{geo:moments}, the following estimates hold uniformly.
Every sum and maximum in these three estimates is restricted to $(q,b)=1$:
\begin{align}
 \max_{q\sim D}A_q^Z(\Delta)&\ll
  \Delta^{1-\sigma}D^{1-c},&D^2\Delta&\le Z,
    \label{geo:row-small}\\
 \sum_{q\sim D}A_q^Z(\Delta)&\ll
  (D^2\Delta)^{1-\sigma},&D^2\Delta&\le Z,
    \label{geo:sum-small}\\
 \max_{q\sim D}A_q^Z(\Delta)&\ll
  Z^{\sigma-c}(D\Delta)^{1+c-2\sigma},&D^2\Delta&\ge Z,
    \label{geo:row-large}
\end{align}
where the last estimate only requires $D\Delta\ll Z$; changing the boundary
between the ranges by a fixed constant does not affect the conclusion.
The first two estimates incur no small-power loss in the ambient length $Z$.
\end{lemma}

\begin{proof}
We first prove a row-sum bound at the natural scale.
Write $Z=NM$, where $M,N$ are powers of the base, $M\le N$, and $D_0=N/M$.
By~\eqref{geo:product},
\[
 F_Z(t)=F_M(t)F_{D_0}(Mt)F_M(Nt).
\]
Apply $uv\le(u^2+v^2)/2$ to the two outer blocks and then use
\eqref{geo:envelope}, obtaining
\[
 \sup_{|v|\le C/Z}F_Z(t+v)
 \ll E_{D_0}(Mt)\{E_M(t)+E_M(Nt)\},
\]
where $\deg E_{D_0}=O(D_0)$, $\int E_{D_0}\ll D_0^{-\sigma}$,
$\deg E_M=O(M)$, and $\int E_M\ll M^{-c}$.
The first product has degree $O(N)$ and integral
$O(M^{-c}D_0^{-\sigma})$: average first conditional on $Mt\pmod1$.
Only $O(1)$ Fourier coefficients of $E_M$ at frequencies divisible by $M$
remain, and each is bounded by its mean.
For the second term, $(M,q)=1$ allows the exact permutation of $Mh\bmod q$:
\[
 \sum_{h\bmod q}E_{D_0}(Mh/q)E_M(Nh/q)
 =\sum_{h\bmod q}E_{D_0}(h/q)E_M(D_0h/q).
\]
The product on the right still has degree $O(N)$, and averaging first
conditional on $D_0t$ gives the same integral bound.
Thus~\eqref{geo:grid-positive} proves
\begin{equation}\label{geo:outer-blocks}
 \sum_{h\bmod q}\sup_{|v|\le C/Z}F_Z(h/q+v)
 \ll(q+N)M^{-c}D_0^{-\sigma}.
\end{equation}
For $\sqrt Z\ll q\ll Z$, choose a neighboring power $N\asymp q$ of the base
with $N\mid Z$ and $N\ge\sqrt Z$, to obtain
\begin{equation}\label{geo:outer-large}
 \sum_{h\bmod q}\sup_{|v|\le C/Z}F_Z(h/q+v)
 \ll Z^{\sigma-c}q^{1+c-2\sigma}.
\end{equation}
Rounding to powers of the base, and the parity of the word length, change
only constants depending on $b$.

For a wider window, choose a power $J\asymp\Delta$ of the base with $J\mid Z$,
and put $Z'=Z/J$. On the integer grid we have the exact factorization
\[
 F_Z(a/Z)=F_{Z'}(a/Z)F_J(a/J).
\]
Each rational window contains $O(J)$ consecutive integers, hence a bounded
number of complete grids modulo $J$ and at most two partial grids.
Sum the second factor over those integers first; positivity and
\eqref{geo:moments} give $O(J^{1-\sigma})$. Then take a local supremum of the first
factor:
\begin{equation}\label{geo:window-average}
 A_q^Z(\Delta)\ll\Delta^{1-\sigma}
 \sum_{\substack{h\bmod q\\(h,q)=1}}
 \sup_{|v|\le C'/Z'}F_{Z'}(h/q+v).
\end{equation}
If $D^2\Delta\ge Z$, apply~\eqref{geo:outer-large} at length $Z'$ to obtain
\eqref{geo:row-large}.

If $D^2\Delta\le Z$, choose a power $M\asymp D$ of the base such that
$M^2\mid Z'$. Discarding the middle digit block gives
\[
 F_{Z'}(t)\le F_M(t)F_M((Z'/M)t).
\]
Apply the arithmetic mean of the squares, permute the coprime residue
classes, and use~\eqref{geo:envelope}. Each envelope for $F_M^2$ has mean
$O(M^{-c})$, and the sampling cost is $O(q+M)$.
Equation~\eqref{geo:window-average} now gives~\eqref{geo:row-small}.
For the sum over all $q\sim D$, shorten the remaining word further to a length
$U\asymp D^2$. The reduced fractions are separated by $\gg D^{-2}$,
and the positive envelope has degree $O(D^2)$ and mean $O(D^{-2\sigma})$.
Lemma~\ref{geo:sampling} together with~\eqref{geo:window-average} gives
$\Delta^{1-\sigma}D^{2-2\sigma}$, which is~\eqref{geo:sum-small}.
Every loss in this argument is an absolute constant once the base is fixed;
bounded $D$ is handled by increasing that constant.
\end{proof}

\subsection{From reduced fractions to the full minor-arc sum}

The two fractions in a simultaneous approximation may have different
reduced denominators. We therefore retain both divisor sums until after
squaring. Weighted Cauchy places the divisor loss on the reduced
denominators, where the preceding estimates supply decay. The original
minor-arc restriction first removes the small rational rows; the digital
fibres can then be enlarged without restoring those rows.

\begin{proposition}[Simultaneous rational approximations with all denominators]\label{geo:all-denom}
Suppose that~\eqref{geo:moments} holds and that
\begin{equation}\label{geo:range-y}
 \sigma>1-c/2,\qquad
 0<y<\min\{2/3,\sigma,1-3g_{4/3}/2\}.
\end{equation}
For $X=b^k$ and $L=\log X$, define
\[
 \mathfrak M_B=
 \bigcup_{\substack{1\le q\le L^B\\(h,q)=1}}
 \{t\in\mathbb R/\mathbb Z:\|t-h/q\|\le L^B/X\},
 \qquad
 \mathcal R_B=\{0\le a<X:a/X\notin\mathfrak M_B\}.
\]
For every $J_0>0$ and fixed $C>0$, there is a choice of $B$ such that,
uniformly for $1\le Y\le X^y$,
\begin{equation}\label{geo:all-denom-bound}
 \sum_{\substack{a_1,a_2\in\mathcal R_B\\
 \exists,1\le q\le CY, h_1,h_2\in\mathbb Z:\
 |qa_i/X-h_i|\le CY/X\ (i=1,2)}}
 F_X(a_1/X)F_X(a_2/X)
 \ll_{J_0,C}Y L^{-J_0}.
\end{equation}
Neither fraction $h_i/q$ is required to be reduced.
\end{proposition}

\begin{proof}
We first control primitive denominators coprime to the base.
For $D\Delta\ll Z^y$, Lemma~\ref{geo:rows} gives
\begin{equation}\label{geo:small-energy}
 \frac1{D\Delta}\sum_{q\sim D,(q,b)=1}A_q^Z(\Delta)^2
 \ll\Delta^{1-2\sigma}D^{2-c-2\sigma},
 \qquad D^2\Delta\le Z.
\end{equation}
If $D^2\Delta>Z$, then
$D^3\Delta^2/Z^2\ll Z^{3y-2}=o(1)$, so Proposition~\ref{geo:energy} applies.
Insert~\eqref{geo:moments} and~\eqref{geo:row-large}, and divide by $D\Delta$.
The resulting three terms are
\[
 Z^{1-c}(D\Delta)^{c-2\sigma},\qquad
 Z^{-1+3g_{4/3}/2}D\Delta^{1/2},\qquad
 Z^{-2\sigma}D^2\Delta.
\]
Since $D\Delta\ge\sqrt Z$ and $c<2\sigma$, these are at most, respectively,
\begin{equation}\label{geo:three-powers}
 Z^{1-\sigma-c/2},\qquad Z^{y-1+3g_{4/3}/2},\qquad
 Z^{2y-2\sigma}.
\end{equation}
Choose a fixed
\[
 0<\delta_0<
 \tfrac12\min\{\sigma+c/2-1,\ 1-y-3g_{4/3}/2,\ 2\sigma-2y\}.
\]
All three entries are positive by~\eqref{geo:range-y}.
Taking the incidence exponent $\varepsilon$ sufficiently small absorbs
both $Z^\varepsilon$ and the fixed logarithmic losses, giving
\begin{equation}\label{geo:primitive-energy}
 \frac1{D\Delta}\sum_{q\sim D,(q,b)=1}A_q^Z(\Delta)^2
 \ll
 \begin{cases}
 \Delta^{-\rho}D^{-\tau},&D^2\Delta\le Z,\\
 Z^{-\delta_0},&D^2\Delta>Z,
 \end{cases}
 \quad
 \rho=2\sigma-1>0,\quad\tau=c+2\sigma-2>0.
\end{equation}

Now write a reduced denominator as $d=gq$, where $g$ is a power of $p$ and
$(q,p)=1$. Let $g_0$ be the smallest power of $b$ not less than $g$.
Then $1\le g_0/g\le p^{r-1}$, and in the present range $g_0\mid X$ for all
sufficiently large $X$.
Put $Z=X/g_0$ and write uniquely
$a=t+Z\ell$, $0\le t<Z$, $0\le\ell<g_0$.
The digital product gives
\begin{equation}\label{geo:smooth-fiber}
 F_X(a/X)=F_{g_0}\left(\frac{t/Z+\ell}{g_0}\right)F_Z(t/Z).
\end{equation}
If $a/X$ is within $C\Delta/X$ of the reduced fraction $h/(gq)$, then $t/Z$
is within $C\Delta/Z$ of the reduced fraction $(g_0/g)h/q$.
Relaxing the restriction on $\ell$ and applying the first moment on shifted
grids yields
\begin{equation}\label{geo:fiber-row}
 A_{gq}^X(\Delta)\ll g_0^{1-\sigma}A_q^Z(\Delta).
\end{equation}
If $gq\Delta\ll Y$, then
$q\Delta\ll Y/g_0\le Z^y$ and $Z\gg X^{1-y}$.
Consequently, the constants in~\eqref{geo:primitive-energy} remain uniform.

To handle unreduced common denominators, we keep the divisor loss on the
reduced denominator. For nonnegative $B_d$ and $\omega>0$, weighted Cauchy gives
\[
 \left(\sum_{d\mid q}B_d\right)^2
 \le\sigma_{-\omega/2}(q)\sum_{d\mid q}d^{\omega/2}B_d^2.
\]
Use $\sigma_{-a}(dm)\le\sigma_{-a}(d)\sigma_{-a}(m)$,
$\sum_{m\le U}\sigma_{-a}(m)\le U\zeta(1+a)$, and
$\sigma_{-a}(d)\ll_\omega d^{\omega/2}$, and interchange nonnegative sums.
This gives
\begin{equation}\label{geo:divisor-cauchy}
 \sum_{q\sim Q}\left(\sum_{d\mid q}B_d\right)^2
 \ll_\omega Q\sum_{d\le2Q}d^{\omega-1}B_d^2.
\end{equation}

Decompose the original common denominators into $q\sim Q$, and take
$\Delta\asymp\max(1,Y/Q)$, so that $Q\Delta\asymp Y$.
Let $B_d$ be the mass of row $d$ in~\eqref{geo:incidence-matrix}, still summed
only over the original set $a\in\mathcal R_B$.
The reduced denominator of every unreduced fraction divides $q$, so the
pair mass $\mathcal P_Q$ supplied by this block satisfies
\[
 \mathcal P_Q
 \le\sum_{q\sim Q}\left(\sum_{d\mid q}B_d\right)^2.
\]
The two fractions in a pair may have different reduced denominators;
the weighted divisor inequality~\eqref{geo:divisor-cauchy} accounts for
their cross terms.
First group by $d=gq'$ and $q'\sim D$.
Choose $T$ later, and take $B>T+2$.
If $gD\Delta\le L^T$, every reduced denominator in the block is at most
$2L^T$, and every window has radius at most $CL^T/X$.
Thus the original mass $B_d$ vanishes identically.
This deletion takes place before removing the original major-arc restriction
and before enlarging the sum over the digital fibers.
Divide each remaining block by $Y$. Its bound from
\eqref{geo:divisor-cauchy} and~\eqref{geo:fiber-row} is
\[
 \frac QY(gD)^{\omega-1}g_0^{2-2\sigma}
       \sum_{\substack{q'\sim D\\(q',b)=1}}A_{q'}^Z(\Delta)^2
 \ll g^{-\rho+\omega}D^\omega
       \frac1{D\Delta}
       \sum_{\substack{q'\sim D\\(q',b)=1}}A_{q'}^Z(\Delta)^2.
\]
Here $Q/Y\asymp\Delta^{-1}$ and $g_0/g$ is bounded in terms of the base.
Thus~\eqref{geo:primitive-energy} bounds the block by
\begin{equation}\label{geo:block-final}
 g^{-\rho+\omega}D^\omega
 \begin{cases}
 \Delta^{-\rho}D^{-\tau},&D^2\Delta\le Z,\\
 Z^{-\delta_0},&D^2\Delta>Z,
 \end{cases}
 \qquad Z=X/g_0.
\end{equation}
Choose
\[
 \omega<\min\{\rho/2,\tau/2,\delta_0(1-y)/(2y)\}.
\]
The large-scale part is
$O(X^{-\delta_0(1-y)/2})$, since $D\ll X^y$ and $Z\gg X^{1-y}$.
The small-scale part is at most
$(gD\Delta)^{-\lambda}$, where
$\lambda=\min\{\rho-\omega,\tau-\omega\}>0$.

Every retained small-scale block satisfies $gD\Delta>L^T$, and hence
contributes at most $L^{-\lambda T}$.
Take $\lambda T>J_0+10$ and sum over the $O(L^3)$ blocks to obtain
\eqref{geo:all-denom-bound}.
This order of operations also handles the rational point zero:
its repeated occurrence among the full denominator rows is removed by the
original major-arc restriction before the enlargement.
\end{proof}

\section{The Type II interval from moment bounds}\label{geo:typeii-section}

The two fractional moments have different roles. An estimate
$\sum_{a<X}F_X(a/X)^s\ll X^{g_s}$ controls the small-amplitude tail
and determines $\kappa=g_s/(2-s)$, while the
$4/3$ moment controls common denominators through the preceding section.
Together they give the interval $(\kappa,1-2\kappa)$ below. We use
Maynard's bilinear Cauchy reduction \cite[\S13]{Maynard2019}, inserting
the weighted common-denominator estimate in its simultaneous-approximation
alternative.

\begin{theorem}[Type II minor-arc estimate from moment bounds]\label{geo:typeii}
Assume the hypotheses of Proposition~\ref{geo:all-denom}, and suppose that,
for some $4/3\le s<2$,
\begin{equation}\label{geo:fractional}
 \sum_{a<X}F_X(a/X)^s\ll X^{(2-s)\kappa},\qquad
 \kappa<1/3,\qquad 2\kappa<y.
\end{equation}
Then every compact interval
$[\alpha_0,\beta_0]\subset(\kappa,1-2\kappa)$
is a Type II window: for each $J>0,C_0>0$, one can choose $B$ such that
\begin{equation}\label{geo:typeii-bound}
 \left|\sum_{h\in\mathcal R_B}F_X(h/X)\gamma_h
 \sum_{n\sim N}\sum_{m\sim M}u_nv_m\mathrm e(-hnm/X)\right|
 \ll_{J,C_0}X(\log X)^{-J},
\end{equation}
uniformly for $X^{\alpha_0}\le N\le X^{\beta_0}$, $NM\le C_0X$,
and all $|u_n|,|v_m|,|\gamma_h|\le1$.
The intervals of support may be shortened, and the coefficients may be
multiplied by any fixed power of the logarithm.
\end{theorem}

The geometric step separates two ways in which the bilinear kernel can be
large: its frequencies either satisfy a short affine relation or admit
simultaneous rational approximations. The common-denominator estimate
controls the second alternative; the incidence estimate and the amplitude
bounds control the first.

\begin{lemma}[The lattice dichotomy for an integer slab]\label{geo:slab-dichotomy}
Fix $C\ge1$. There is a constant $K_0=K_0(C)$ with the following property.
Let $X\ge2$ be an integer, let $N\ge1$, and suppose that
\[
 C^{-1}N/X\le\delta\le1,\qquad K\ge K_0,\qquad
 \delta N^2K\ge K_0.
\]
For $a_1,a_2\in\{0,\ldots,X-1\}$, put
$\mathbf A=(a_1,a_2,X)$ and
\[
 \mathcal S=\{v\in\mathbb Z^3:|v|\le CN,\
                         |\mathbf A\cdot v|\le C\delta X\}.
\]
If $|\mathcal S|\ge\delta N^2K$, then at least one of the following holds:
\begin{enumerate}
 \item there are integers $v_1,v_2,v_3,v_4$, with
 $(v_1,v_2,v_3)\ne0$, such that
 \[
  v_1a_1+v_2a_2+v_3X+v_4=0,\qquad
  \max_{1\le i\le4}|v_i|\ll_C X/(N^2K);
 \]
 \item there are integers $q,h_1,h_2$ such that
 \begin{equation}\label{geo:simultaneous}
  1\le q\ll_C X/(NK),\qquad
  |qa_i/X-h_i|\ll_C1/(NK)\quad(i=1,2).
 \end{equation}
\end{enumerate}
No coprimality is asserted in the second alternative.
\end{lemma}

\begin{proof}
This is the combination of the geometric reductions in
\cite[Lemmas 13.2 and 14.1]{Maynard2019}, with the lower bound on the
number of points made explicit.
Stretch the direction $\mathbf A/|\mathbf A|$ by $N/\delta$ and leave its
orthogonal complement fixed; denote this map by $\Phi$.
The lattice $\Phi(\mathbb Z^3)$ has determinant $N/\delta$, and
$\Phi(\mathcal S)$ lies in a ball of radius $O_C(N)$.
By the fixed-dimensional basis and successive-minima estimates
\cite[Section~5]{Maynard2019}, a reduced lattice basis $b_i$, with
$s_1\le s_2\le s_3$ and $s_i=|b_i|$, satisfies
\[
 \left|\sum_i z_i b_i\right|\asymp\sum_i|z_i|s_i,\qquad
 s_1s_2s_3\asymp\det\Phi(\mathbb Z^3)=N/\delta.
\]
Thus the coordinates of any point in that ball satisfy
$|z_i|\ll_C N/s_i$.
If $s_3$ is at most a sufficiently large constant times $N$, the ball
contains only
$O_C(N^3/(s_1s_2s_3))=O_C(\delta N^2)$ lattice points, contradicting
$K\ge K_0$ when $K_0$ is large enough.
Otherwise every point in the ball has $z_3=0$.
Consequently the linear span of $\mathcal S$ has dimension at most two.
It has positive dimension since $|\mathcal S|\ge K_0>1$.

If this span is a line, it passes through the origin, and its integer
points are multiples of a primitive vector $v$.
Because $\mathcal S$ is symmetric and contains the origin, its largest
multiple has coefficient $|j|\gg|\mathcal S|$.
The two defining inequalities of $\mathcal S$ therefore give
\[
 |v|\ll_C\frac1{\delta NK}\ll_C\frac X{N^2K},
 \qquad
 |\mathbf A\cdot v|\ll_C\frac X{N^2K}.
\]
Taking $v_4=-\mathbf A\cdot v\in\mathbb Z$ proves the first alternative.

If the span has dimension two, take a reduced basis of the transformed
section lattice, of lengths $s_1,s_2$.
Since $\Phi(\mathcal S)$ contains two independent vectors of length
$O_C(N)$, both basis lengths are $O_C(N)$.
The planar lattice-point bound gives
\[
 \delta N^2K\le|\mathcal S|\ll_C\frac{N^2}{s_1s_2},
 \qquad s_1s_2\ll_C\frac1{\delta K}.
\]
Pull the basis back to integer vectors $w_1,w_2$. Then
\[
 |w_i|\le s_i,\qquad
 |\mathbf A\cdot w_i|\ll_C\delta Xs_i/N.
\]
For $v=w_1\times w_2\ne0$, the projection of $\mathbf A$ onto
$\operatorname{span}(w_1,w_2)$ is
\[
 E=\frac{(\mathbf A\cdot w_1)(w_2\times v)
       +(\mathbf A\cdot w_2)(v\times w_1)}{|v|^2}.
\]
Hence
$|E|\ll_C\delta Xs_1s_2/(N|v|)\ll_C X/(NK|v|)$ and
$\mathbf A=\lambda v+E$.
Since $|v|\ge1$, $|\mathbf A|\asymp X$, and $A_3=X$, taking $K_0$
large enough gives $|v_3|\asymp|v|$.
Change the sign of $v$ so that $q=v_3>0$ and put $h_i=v_i$.
Then
\[
 q\le|v|\le s_1s_2\ll_C\frac1{\delta K}\ll_C\frac X{NK},
 \qquad
 |qa_i-Xh_i|=|v_3E_i-v_iE_3|
 \ll_C|v|\,|E|\ll_C\frac X{NK}.
\]
This is the second alternative.
\end{proof}

\begin{proof}[Proof of Theorem~\ref{geo:typeii}]
Choose $u_0$ sufficiently close to $1-2\kappa$ that
\begin{equation}\label{geo:u-choice}
 \beta_0<u_0<1-2\kappa,\qquad 1-u_0<y,\qquad
 2\alpha_0+u_0>1.
\end{equation}
These three requirements follow from $\beta_0<1-2\kappa$,
$2\kappa<y$, and $\alpha_0>\kappa$, respectively.
Choose fixed $\varepsilon_0,\xi>0$ sufficiently small that
\begin{align*}
 u_0+2\kappa+2(s-1)\varepsilon_0&<1,\\
 2\kappa-2\alpha_0+(3s-4)\varepsilon_0+\xi&<0,\\
 6\kappa-1-3\alpha_0+6(s-1)\varepsilon_0+2\xi&<0.
\end{align*}
The first inequality follows from the choice of $u_0$, the second from
$\alpha_0>\kappa$, and the third from
$6\kappa-1-3\alpha_0<3\kappa-1<0$.
These choices are made before choosing the requested logarithmic saving.
Let
$E=\{a:F_X(a/X)\ge X^{-\kappa-\varepsilon_0}\}$.
Writing $g=g_s=(2-s)\kappa$, we have
\begin{equation}\label{geo:tail}
 \sum_{a\notin E}F_X(a/X)^2\ll X^{-(2-s)\varepsilon_0},
 \qquad H_E:=\sum_{a\in E}F_X(a/X)
 \ll X^{\kappa+(s-1)\varepsilon_0}.
\end{equation}
Let $S(h)$ denote the bilinear exponential sum in~\eqref{geo:typeii-bound}.
Group the products by $nm$; their coefficients are bounded by the divisor
function and supported in $O_{C_0}(X)$.
Folding modulo $X$ introduces only $O_{C_0}(1)$ aliases.
Parseval's identity and
$\sum_{n\le CX}\tau(n)^2\ll_CX(\log X)^3$ give
\[
 \sum_{h<X}|S(h)|^2\ll_{C_0}X^2(\log X)^3.
\]
Cauchy's inequality therefore bounds the cost of discarding $a\notin E$ by
$O(X^{1-(2-s)\varepsilon_0/2}(\log X)^{3/2})$.

Split the retained frequencies into $O(\log X)$ amplitude layers
$E_T=\{a:F_X(a/X)\asymp T^{-1}\}$,
$1\le T\ll X^{\kappa+\varepsilon_0}$.
Fix one layer, retaining the restriction to $\mathcal R_B$.
Apply Cauchy's inequality in $m$, then take absolute values of the geometric
sums over $m$. The resulting square is at most $CX/N$ times the positive quantity
\begin{equation}\label{geo:cauchy-kernel}
 \mathcal K=
 \sum_{a_1,a_2\in E_T\cap\mathcal R_B}F_X(a_1/X)F_X(a_2/X)
 \sum_{1\le n_1,n_2\le C_1N}
 \min\left(\frac{C_2X}{N},
 \left\|\frac{n_1a_1-n_2a_2}{X}\right\|^{-1}\right).
\end{equation}
Our target is $\mathcal K\ll_JNX(\log X)^{-J}$.
For each pair $(a_1,a_2)$, denote the inner sum in
\eqref{geo:cauchy-kernel} by $\mathscr H_N(a_1,a_2)$.
It lies between constant multiples of $N^2$ and $XN$.
Partition the pairs according to
$\mathscr H_N(a_1,a_2)\asymp N^2K$, with $K$ dyadic and
$1\ll K\ll X/N$.
For all layers with $NK\le X^{u_0}$, equation~\eqref{geo:tail} gives directly
\begin{equation}\label{geo:small-K}
 \mathcal K_{\rm small}
 \ll N X^{u_0+2\kappa+2(s-1)\varepsilon_0}\log X
 =o\bigl(NX(\log X)^{-J}\bigr).
\end{equation}
Choose $\varepsilon_0$ so that the exponent here is strictly less than one.

The remaining layers satisfy $NK>X^{u_0}$.
Decompose distances into ranges $N/X\ll\delta\le1$:
each kernel layer is bounded by $C\delta^{-1}$ times the indicator that the
distance is at most $C\delta$.
Hence there is a choice of $\delta$ for which the integer slab
\begin{equation}\label{geo:slab}
 |v|\le C_3N,\qquad
 |v_1a_1-v_2a_2-v_3X|\le C_4\delta X
\end{equation}
contains $\gg\delta N^2K'$ lattice points, where
$K'=K/(C_5\log X)$.
These points come from the original positive integers $n_1,n_2$;
enlarging to the symmetric slab only increases their number.
Each $(v_1,v_2)$ permits only $O(1)$ choices for $v_3$, so
$\delta\ll1/K'$.
Moreover,
\begin{equation}\label{geo:slab-margins}
 K'\gg X^{u_0-\beta_0}/\log X\longrightarrow\infty,
 \qquad
 \delta N^2K'\gg X^{2\alpha_0+u_0-1}/\log X\longrightarrow\infty.
\end{equation}

The two lower bounds in~\eqref{geo:slab-margins} ensure that
Lemma~\ref{geo:slab-dichotomy} applies uniformly to every remaining
layer once $X$ is sufficiently large. Change the signs of the second
and third coordinates in~\eqref{geo:slab}; enlarge the fixed constant $C$
and choose $C_5$ so that the lemma applies with $K'$ in place of $K$.
Put $Y=X/(NK')$. The bounds on $K$ and~\eqref{geo:u-choice} give
\[
 1\ll\log X\ll Y\ll X^{1-u_0}\log X<X^y.
\]

For pairs in the simultaneous-approximation alternative,
Proposition~\ref{geo:all-denom} bounds the weighted mass by
$O_J(Y(\log X)^{-J})$.
Multiplying by the kernel bound $O(N^2K)$ costs only the ratio
$K/K'=O(\log X)$, giving a contribution
$O_J(NX(\log X)^{-J+1})$ to~\eqref{geo:cauchy-kernel}.
All constants are uniform in the amplitude, kernel-size, and distance layers.

For the line alternative, Lemma~\ref{geo:slab-dichotomy} gives a nonzero
integer relation of the form in Lemma~\ref{geo:incidence}, with all four
heights $O(V)$, where $V=X/(N^2K')$.
If $V$ is too small, integrality excludes this case; otherwise
Lemma~\ref{geo:incidence} applies with both height parameters comparable
to $V$.

Write $R=|E_T|$, $\zeta=\log_XT$, and $h=g+(s-1)\zeta$.
The $s$th moment gives
\[
 R/T\ll X^h,\qquad R\ll X^{h+\zeta}.
\]
Put $\nu=\log_XN$ and $u=\log_X(NK')$.
The weighted mass of pairs giving a line is bounded both by $(R/T)^2$ and by
\begin{equation}\label{geo:line-mass}
 X^\varepsilon T^{-2}\left(RV^2+R^2V^3/X\right).
\end{equation}
Relative to the target scale $Y=X/(NK')$, the exponent in the trivial bound is
$2h-1+u$.
Use that bound directly when $u\le1-2h-\xi$, with the $\xi$ fixed above.
In the remaining cases, the exponents of the ratios of the two terms in
\eqref{geo:line-mass} to $Y$ are at most
\[
 3h-\zeta-2\nu+\xi,\qquad 6h-1-3\nu+2\xi.
\]
Since $s\ge4/3$ and
$\zeta\le\kappa+\varepsilon_0+O(1/\log X)$,
\[
 3h-\zeta=3g+(3s-4)\zeta
 \le2\kappa+(3s-4)\varepsilon_0+o(1),\qquad
 h\le\kappa+(s-1)\varepsilon_0+o(1).
\]
The two exponents are therefore at most
\begin{equation}\label{geo:line-powers}
 2\kappa-2\alpha_0+(3s-4)\varepsilon_0+\xi+o(1),\qquad
 6\kappa-1-3\alpha_0+6(s-1)\varepsilon_0+2\xi+o(1).
\end{equation}
The initial choice of $\varepsilon_0,\xi$ makes both exponents strictly
negative. Choose the incidence loss $\varepsilon$ smaller than these
fixed margins.
The line contribution thus saves a fixed power, sufficient to absorb
$K/K'$ and all logarithmic losses from the amplitude, distance, and kernel-size
decompositions.
Combining this with~\eqref{geo:small-K} and the rank-two estimate proves
$\mathcal K\ll_JNX(\log X)^{-J}$.
Return to the initial Cauchy inequality and increase $J$ to obtain
\eqref{geo:typeii-bound}.
\end{proof}


\section{Arithmetic comparison on prime-factor regions}\label{arith:section}

The sieve uses two arithmetic consequences of the Fourier estimates.
First, a region of prime factorizations has the expected digital density
whenever a proper subproduct lies in the Type II window
(Proposition~\ref{arith:polytope}). Second, the same comparison holds for
marked products with a rough cofactor (Proposition~\ref{arith:rough}).
The latter reaches marked products of exponent $1-\alpha$ using
Type I distribution to any fixed level just beyond $1-\beta$.
The interval Type I estimate in Proposition~\ref{arith:typeI} also
applies to the octal words in Section~\ref{sec:octal-typeI}.

We prove these statements by Maynard's method of prime-factor boxes and
Harman's sieve \cite{Maynard2019,Harman2007}. The major arcs determine the
local density, the minor arcs permit a suitable grouping of the prime
factors, and a first-crossing argument supplies the rough-cofactor comparison.

Let the base $b=\ell^\nu\ge4$ be a fixed prime power, let $d$ be the excluded
digit, and write
\[
 \mathcal D=\{0,\ldots,b-1\}\setminus\{d\},\quad
 X=b^k,\quad c=\log_b(b-1),\quad A=(b-1)^k,
\]
\[
 \mathcal A_X=\left\{\sum_{j=0}^{k-1}d_jb^j:d_j\in\mathcal D\right\},
 \qquad P_X(t)=A^{-1}\sum_{n\in\mathcal A_X}e(nt),\quad F_X=|P_X|.
\]
Leading zeros are allowed here when $d\ne0$; when $d=0$, $\mathcal A_X$
contains only admissible integers with exactly $k$ digits. Put $L=\log X$
and
\begin{equation}\label{arith:localfactor}
 u_d=\#\{a\in\mathcal D:(a,b)=1\},\qquad
 \gamma_d=\frac{b\,u_d}{\varphi(b)(b-1)}.
\end{equation}
All sieve sums and prime-factor sums are restricted to $1\le n<X$.
The digital Fourier polynomial may include zero, since the arithmetic
weight paired with it vanishes there.
The number of marked prime factors, the coefficients of every polyhedral
condition, and the sieve endpoints are fixed as $X\to\infty$.
Implied constants may depend on these data and on $b,d$.
In the first-crossing argument below, an auxiliary cutoff exponent is
held fixed while $X$ tends to infinity, and only then tends to zero.

\begin{lemma}[A uniform bound for small products]\label{arith:small}
For $1\le Y\le X$, we have
\[
 \#(\mathcal A_X\cap[0,Y])\le(b-1)Y^c.
\]
Suppose that a family of marked factorizations has total absolute weight
at most $C(\log X)^K$ at each positive integer. Its combined digital and
normalized reference mass on $n\le X^{1-\delta}$ is at most
\begin{equation}\label{arith:small-bound}
 O\left(CA(\log X)^K
       \{X^{-c\delta}+X^{-\delta}\}\right).
\end{equation}
Consequently, if $\delta\log X/\log\log X\to\infty$, this quantity is
smaller than every fixed $A(\log X)^{-J}$.
\end{lemma}
\begin{proof}
Choose $b^j\le Y<b^{j+1}$. If zero is allowed, all digits above position
$j$ must vanish, leaving at most $(b-1)^{j+1}\le(b-1)Y^c$ integers.
If zero is excluded and $k>j+1$, there are no such integers; in the
remaining cases the total number of words gives the same bound.
The reference weight is $\gamma_d A/X$, so the reference mass of integers
$n\le Y$ is $O(AY/X)$. Multiplying by the bound for the absolute weight
at each integer proves the second assertion.
A product of a fixed number $r$ of marked primes has at most $r!$
markings, and its logarithmic weight is at most $L^r$. If all prime
factors exceed a fixed $X^\epsilon$, their number is at most
$\lfloor1/\epsilon\rfloor$, so this case also satisfies the stated
hypothesis.
\end{proof}

\subsection{Type I distribution and the local density}

\begin{proposition}[Type I estimate]\label{arith:typeI}
Suppose that $1/2<\sigma<c$ and that
$\int_0^1F_Z(t)\,dt\ll Z^{-\sigma}$ for every power $Z$ of the base.
Define
\[
 D_X(q)=\sup_{I,a}\left|
 \#\{n\in\mathcal A_X\cap I:n\equiv a\pmod q\}
 -\frac{\#(\mathcal A_X\cap I)}q\right|,
\]
where $I$ ranges over real intervals. For fixed $K\ge0$ and
$\epsilon>0$, there is a constant $c_1>0$ such that
\begin{equation}\label{arith:typeI-max}
 \sum_{\substack{q\le Q\\(q,b)=1}}\tau(q)^K D_X(q)
 \ll A\left(e^{-c_1\sqrt{\log X}}+
                      Q^{1+\epsilon}X^{-\sigma}\right).
\end{equation}
In particular, for every fixed $\theta<\sigma$ and $J>0$,
\begin{equation}\label{arith:typeI-unit}
 \sum_{\substack{q\le X^\theta\\(q,b)=1}}\tau(q)^K
 \left|\#\{n\in\mathcal A_X:(n,b)=1,\ q\mid n\}
                -\frac{u_d A}{(b-1)q}\right|
 \ll A L^{-J}.
\end{equation}
\end{proposition}
\begin{proof}
Decompose $I\cap[0,X)$ into aligned radix intervals, with at most
$2(b-1)$ intervals of each length $Z=b^j$. On an interval with an
admissible fixed prefix, the free lower digits form a translate of
$\mathcal A_Z$. Thus
\[
 U_X(t)=2b\sum_{Z\le X,\ Z=b^j}Z^cF_Z(t)
\]
simultaneously bounds the absolute values of the Fourier sums over all
intervals.
If a trigonometric polynomial $P$ has frequencies in $[-M,M]$, and a
set $\Theta$ has circular spacing at least $\delta$, then
\begin{equation}\label{arith:sampling}
 \sum_{t\in\Theta}|P(t)|\ll(M+\delta^{-1})\int_0^1|P(t)|\,dt.
\end{equation}
Indeed, reproduce $P$ with $2\operatorname{Fej}_{2M}-\operatorname{Fej}_M$.
The absolute value of this kernel is bounded by
$CM/(1+M\|t\|)^2$, whose sum over a $\delta$-separated set is
$O(M+\delta^{-1})$. Integration gives the assertion.

The reduced fractions $h/q$ with $R\le q<2R$ are separated by at least
$1/(4R^2)$. When $Z\ge R^2$, discard the higher-digit factors and replace
$F_Z$ by the lower-digit word whose length is comparable to $R^2$;
otherwise use length $Z$ directly. By \eqref{arith:sampling},
\[
 \sum_{q\sim R}\ \sum_{\substack{0<h<q\\(h,q)=1}}F_Z(h/q)
 \ll
 \begin{cases}R^{2-2\sigma},&Z\ge R^2,\\
 R^2Z^{-\sigma},&Z<R^2.
 \end{cases}
\]
Put $E_q=q^{-1}\sum_{0<h<q,(h,q)=1}U_X(h/q)$.
Since $c>\sigma$, summing geometrically over $Z$ gives
\begin{equation}\label{arith:primitiveI}
 \sum_{q\sim R}E_q\ll A\{R^{1-2\sigma}+RX^{-\sigma}\}.
\end{equation}

For small denominators we also use contraction along digital orbits.
The alphabet contains two consecutive digits, so there is a constant
$c_2>0$ for which the one-digit factor
$f(t)=(b-1)^{-1}|\sum_{a\in\mathcal D}e(at)|$
satisfies $f(t)\le\exp(-c_2\|t\|^2)$.
If $(h,q)=(q,b)=1$, every $b^jh/q$ is nonintegral and is at distance
at least $1/q$ from the nearest integer. When that distance is less
than $1/(2b)$, the next step multiplies it by exactly $b$.
Hence a fixed amount of contraction occurs every $O(\log(2q))$ steps,
and therefore
\[
 F_Z(h/q)\ll\exp\left(-c_3\frac{\log Z}{\log(2q)}\right).
\]
Take $R_0=\exp(t_0\sqrt{\log X})$, with $t_0$ sufficiently small.
The terms with $Z<Xe^{-\sqrt{\log X}}$ in $U_X$ already have total
mass $O(Ae^{-c_4\sqrt{\log X}})$; apply the orbit bound to the remaining
lengths. Decreasing $t_0$ if necessary gives
\[
 \sum_{2\le q\le R_0,(q,b)=1}\tau(q)^K E_q
       \ll A e^{-c_5\sqrt{\log X}}.
\]
Orthogonality and reduction of fractions yield
\[
 D_X(q)\le\sum_{r\mid q,\ r>1}\frac rq E_r.
\]
Interchange the finite nonnegative sums, using
$\tau(rj)\le\tau(r)\tau(j)$ and
$\sum_{j\le T}\tau(j)^K/j\ll(\log(2T))^{C_K}$.
For small $r$ use the preceding bound; for large $r$ use
\eqref{arith:primitiveI} and $\tau(r)^K\ll r^{\epsilon_0}$, where
$0<\epsilon_0<\min(\epsilon/4,(2\sigma-1)/2)$.
Summation over dyadic ranges, with the logarithmic factors absorbed into
the stretched-exponential decay or into $Q^{\epsilon-\epsilon_0}$,
proves \eqref{arith:typeI-max}.

Finally, put $Z=X/b$. Every admissible word coprime to $b$ has a unique
representation $n=bm+a$, where $m\in\mathcal A_Z$, $a\in\mathcal D$ and
$(a,b)=1$. For $(q,b)=1$, the condition $q\mid n$ specifies one residue
class for $m$ modulo $q$ for each final digit.
Thus the error is at most $u_dD_Z(q)$, and the main term is
$u_dZ^c/q=u_dA/((b-1)q)$.
Choose $(1+\epsilon)\theta<\sigma$ in \eqref{arith:typeI-max} to obtain
\eqref{arith:typeI-unit}. The same proof allows an interval restriction,
retaining the actual digital count in that interval as the mean.
\end{proof}

\subsection{The major-arc comparison on prime boxes}

\begin{proposition}[Prime-factor boxes with logarithmic weights]\label{arith:major}
Fix $\epsilon>0$ and an integer $r\ge1$, and put $\delta=L_2^{-1}$,
where $L_2=\log\log X$.
Take $a_i\ge\epsilon/2$ and $\sum_{i<r}a_i<1-\epsilon/2$, and define
\[
 R(n)=\sum_{\substack{p_1\cdots p_r=n\\
 X^{a_i}<p_i\le X^{a_i+\delta}\ (i<r)\\
 p_r\ge\max(X^{\epsilon/4},X^{1-\sum_{i<r}a_i-r\delta})}}
       \prod_{i=1}^r\log p_i\qquad(1\le n<X).
\]
Write $S_R(t)=\sum_{n<X}R(n)e(nt)$ and $W_R=S_R(0)$.
Let $\mathfrak M(C)\subset\{0,\ldots,X-1\}$ consist of the indices $h$
for which $h/X$ is within distance $L^C/X$ of a reduced rational point
with denominator at most $L^C$. For fixed $C,J>0$,
\begin{equation}\label{arith:majorformula}
 \frac1X\sum_{h\in\mathfrak M(C)}
       \left(\sum_{n\in\mathcal A_X}e(nh/X)\right)S_R(-h/X)
   =\gamma_d\frac AX W_R+O(A L^{-J}).
\end{equation}
The implied constant is uniform in the admissible box endpoints.
\end{proposition}
\begin{proof}
If a reduced denominator $q\le L^C$ has a prime factor other than $\ell$,
the first $k/3$ orbit points $b^ja/q$ are nonintegral. A perturbation
$|v|\le L^C/X$ reduces their distance lower bound $1/q$ by at most
a factor of two. The orbit argument used for Type I therefore gives
\[
 F_X(a/q+v)\ll_C
       \exp\left(-c_C\frac{\log X}{\log\log X}\right).
\]
There are $O(L^{3C})$ grid frequencies on the major arcs, and
$S_R\ll X L^{O_r(1)}$, so these denominators contribute an admissible
error.

The remaining denominators are $q=\ell^j$. Denote the product of the
first $r-1$ primes by $m$, and their marked logarithmic coefficient
by $C_X(m)$. We have
\[
 \sum_m\frac{C_X(m)}m\ll L^{r-1},\qquad m\le X^{1-\epsilon/3}
\]
for sufficiently large $X$.
Temporarily remove the lower bound on the last prime, obtaining
$\widetilde S_R(t)=\sum_mC_X(m)\sum_{p<X/m}(\log p)e(mpt)$.
The added products do not exceed $X^{1-\epsilon/12}+X^{1-\delta}$.
Chebyshev's bound and the preceding reciprocal sum imply
\begin{equation}\label{arith:freeprime}
 \sup_t|\widetilde S_R(t)-S_R(t)|\ll_J XL^{-J}
\end{equation}
for every fixed $J$.

The Siegel--Walfisz theorem \cite[Corollary~5.29]{IwaniecKowalski2004} gives, for
$Y=X/m\ge X^{\epsilon/3}$,
$q\le L^C$, $(a,q)=1$ and $0\le z\le Y$,
\[
 \sum_{\substack{p\le z\\p\equiv a\pmod q}}\log p
       =\frac z{\varphi(q)}+O_{C,J,\epsilon}(YL^{-J}).
\]
The trivial bound for very small $z$, together with the usual theorem
for the remaining $z$, gives this uniform form.
Partial summation then yields, for $|v|\le L^C$,
\begin{equation}\label{arith:APprime}
 \sum_{p<Y}(\log p)e(amp/q)e(vmp/X)
 =\frac{\mu(q)}{\varphi(q)}\int_0^Y e(vmt/X)\,dt
       +O(q(1+|v|)YL^{-J}).
\end{equation}
Here $(m,q)=1$, and the finitely many primes that are not units can be
absorbed into the error.

For sufficiently large $X$, we have $q\mid X$, so the offsets on the
major-arc integer grid are exactly $v/X$, with $v\in\mathbb Z$.
When $v\ne0$,
$\int_0^{X/m}e(vmt/X)dt=m^{-1}\int_0^Xe(vu/X)du=0$.
Thus \eqref{arith:freeprime}--\eqref{arith:APprime} handle all nonzero
offsets. At zero offset, the same calculation with $q=1$ gives
\[
 W_R=X\sum_m\frac{C_X(m)}m+O_J(XL^{-J}),
 \qquad
 S_R(a/q)=\frac{\mu(q)}{\varphi(q)}W_R+O_J(XL^{-J})
 \quad(q=\ell^j).
\]
Since $\mu(\ell^j)=0$ for $j\ge2$, only $0$ and $a/\ell$,
$1\le a<\ell$, remain. Put $S_{\mathcal A}(t)=AP_X(t)$.
Their digital coefficients satisfy
\[
 S_{\mathcal A}(0)-\frac1{\ell-1}
                  \sum_{a=1}^{\ell-1}S_{\mathcal A}(a/\ell)
 =\frac{\ell}{\ell-1}\#\{n\in\mathcal A_X:(n,b)=1\}
 =\gamma_d A.
\]
Multiply this identity by the arithmetic main term $W_R$ and divide
by $X$ to obtain \eqref{arith:majorformula}. All logarithmic powers
introduced by the denominators, offsets and preceding prime sums
are absorbed by choosing $J$ sufficiently large in Siegel--Walfisz.
\end{proof}

\subsection{From boxes to prime-factor regions}

We now assume that the bounded-coefficient minor-arc estimate of
Theorem \ref{geo:typeii} holds on a fixed short-factor window
$[\alpha_0,\beta_0]\subset(0,1/2)$.
Choose a sieve window strictly inside it:
\begin{equation}\label{arith:window}
 \alpha_0<\alpha<\beta<\beta_0,
 \qquad \mathcal W=[\alpha,\beta]\cup[1-\beta,1-\alpha].
\end{equation}

\begin{proposition}[Comparison on prime-factor polytopes]\label{arith:polytope}
Fix an integer $r\ge2$ and $\epsilon>0$, and consider marked prime tuples with
$p_1\cdots p_r<X$ and $p_i\ge X^\epsilon$.
One may impose any Boolean combination of finitely many fixed affine
inequalities, whose variables are either $\log p_i/\log X$ or
$\log p_i/\log(p_1\cdots p_r)$.
Suppose that each admissible tuple has a nonempty proper subproduct
whose logarithm to base $X$ lies in $\mathcal W$. Then
\begin{equation}\label{arith:polytopeformula}
 \sum_{\boldsymbol p}^{\mathcal R}
       1_{\mathcal A_X}(p_1\cdots p_r)
 =\gamma_d\frac AX\sum_{\boldsymbol p}^{\mathcal R}1
       +o(A/L).
\end{equation}
The same conclusion holds when the subproduct condition is normalized
by the total product. Finitely many such regions may be added, with
their respective multiplicities retained.
\end{proposition}
\begin{proof}
First convert conditions normalized by the total product exactly.
Write $n=p_1\cdots p_r$, $t_i=\log p_i/L$ and
$s_0=\sum_i t_i=\log n/L>0$. For any comparison sign $\bowtie$,
\begin{equation}\label{arith:homogenize}
 \lambda_0+\sum_i\lambda_i\frac{\log p_i}{\log n}\ \bowtie\ 0
 \quad\Longleftrightarrow\quad
 \sum_i(\lambda_i+\lambda_0)t_i\ \bowtie\ 0.
\end{equation}
Thus every condition becomes affine in the fixed-$X$ coordinates.
In particular, a normalized condition that is identically zero on
the simplex is decided exactly, rather than by a boundary approximation.

First consider a box from Proposition \ref{arith:major}.
If a specified proper subproduct lies in the analytic window or its
reflection throughout the box, combine its primes into a factor $u$
and the remaining primes into a factor $v$.
Each coefficient has absolute value at most $r!L^r$; repeated primes
do not affect this bound.
First insert the implicit upper bound $p_r\le X^{1-\sum_{i<r}a_i}$,
and then split $u,v$ into dyadic ranges.
Every nonempty block satisfies $UV\ll X$, with its short factor in
the analytic window.
Separate the sole remaining coupling $uv<X$ by the truncated Perron
formula with threshold $X-1/2$, real part $1/L$, and height $X^{10}$.
The logarithmic distance of an integer $uv$ from this threshold is
$\gg X^{-1}$, giving an error $O(X^{-9})$ per term, still negligible
after summation.
The total variation of the Perron measure is $O(L)$, and the inserted
factors $u^{-it},v^{-it}$ are bounded coefficients.
The minor-arc estimate therefore applies, with its arbitrary logarithmic
saving absorbing $r!L^r$ and the losses from these decompositions.
Together with \eqref{arith:majorformula}, this gives
\begin{equation}\label{arith:boxcomparison}
 \sum_{n\in\mathcal A_X}R(n)
 =\gamma_d\frac AX\sum_{n<X}R(n)+O_J(AL^{-J}).
\end{equation}

Put $\delta=1/\log\log X$ and write
$e(\boldsymbol a)=(a_1,\ldots,a_{r-1},1-\sum_{i<r}a_i)$.
Within the box, each prime-factor logarithm is
$L e_i(\boldsymbol a)(1+O_{r,\epsilon}(\delta))$.
The prime number theorem and partial summation give
\[
 \sum_{X^{a_i}<p\le X^{a_i+\delta}}\frac{\log p}{p}
       \ll\delta L,
 \qquad \sum_{n<X}R(n)\ll X(\delta L)^{r-1}.
\]
Hence the unweighted reference mass of a box is
$O(X\delta^{r-1}/L)$.
Positivity and \eqref{arith:boxcomparison} give the corresponding
bound for the digital box. After removing the logarithmic weights,
the comparison error is
\begin{equation}\label{arith:boxerror}
 O(A\delta^r/L)+O_J(AL^{-J}).
\end{equation}

By Lemma \ref{arith:small}, we may first discard total products
$n<X^{1-\delta}$. Partition the first $r-1$ logarithms of the remaining
tuples into boxes of width $\delta$.
There are $O(\delta^{-(r-1)})$ boxes, and their exponent vectors are
within $O_r(\delta)$ of the simplex points $e(\boldsymbol a)$ above.
The fixed buffer in \eqref{arith:window} ensures that, in every box
meeting the target region, the specified subproduct remains in the
larger analytic window. This also covers a subproduct condition
normalized by $n$, since $1-\delta\le s_0<1$ and the two
normalizations differ by $O_r(\delta)$.
If more than one subproduct is available, partition according to a
fixed order on the finitely many proper subsets, obtaining finitely
many regions of the same type.

Before counting boundary boxes, resolve every affine expression
that is constant on $\sum t_i=1$.
A nonzero constant has a fixed sign in a sufficiently small
neighborhood of that simplex. If the constant is zero, the
fixed-$X$ expression is $\lambda(\sum t_i-1)$: its truth value is
determined by $n<X$ when $\lambda\ne0$, and by exact equality when
$\lambda=0$. This includes the homogeneous expressions in
\eqref{arith:homogenize}.
Every remaining affine hyperplane meets at most
$O(\delta^{-(r-2)})$ boxes.
The positive box bounds show that the total boundary mass is
$O(A\delta/L)+O_J(AL^{-J})$; summing \eqref{arith:boxerror} gives the
same bound for the interior errors.
This handles both strict and nonstrict inequalities, as well as the
ordering of the prime factors.
Letting $X\to\infty$ proves the proposition.
\end{proof}

\begin{lemma}[Square factors]\label{arith:squarefree}
Assume \eqref{arith:typeI-unit} for every fixed $\theta<\sigma$,
and suppose that there is a $t>0$ with $2t<\sigma$ and $1-t<c$.
For fixed $\epsilon>0$, the integers all of whose prime factors are
at least $X^\epsilon$ and which are divisible by a prime square have
total mass $o(A/L)$ in both the digital and normalized reference
sequences.
For the applications to bases eight and nine in this section, one
may take $t=3/10$.
\end{lemma}
\begin{proof}
For $X^\epsilon\le p\le X^t$, apply \eqref{arith:typeI-unit} to the
moduli $p^2\le X^{2t}$. Their total digital mass is
$O(AX^{-\epsilon})+O_J(AL^{-J})$.
These rough integers are coprime to the base, so the final-digit
restriction applies.
For $p>X^t$, even the count over all integers is
\[
 \sum_{p>X^t}X/p^2\ll X^{1-t}=o(A/L).
\]
The reference mass is smaller after multiplication by
$\gamma_dA/X$. Marking a fixed number of prime factors introduces
only a fixed multiplicative constant.
\end{proof}

\subsection{Rough cofactors and the first crossing of the Type II interval}

The useful stopping point in the Buchstab expansion is $X^{1-\beta}$.
Each newly introduced prime is at most $X^{\beta-\alpha}$, so the first
product to cross that threshold is still at most $X^{1-\alpha}$.
It therefore lies in the reflected Type II interval. This explains why
Type I distribution beyond $1-\beta$ suffices for the rough sums below.

Put
\[
 w(n)=1_{\mathcal A_X}(n)-\gamma_d A/X\quad(1\le n<X),
 \qquad S_e(z)=\sum_{\substack{m<X/e\\P^-(m)>z}}w(em).
\]
Here $P^-(1)=+\infty$, and $S_e(p^-)$ denotes the sum with $P^-(m)\ge p$.

\begin{lemma}[Truncating a cofactor threshold]\label{arith:threshold}
Let $e$ be a product of a fixed number $r$ of marked primes, and let
$z=z(\boldsymbol p)\ge2$. Suppose that $ez\le X^{1-\delta}$ throughout
a family $\mathcal R$, for a fixed $\delta>0$. Then
\[
 \sum_{\boldsymbol p}^{\mathcal R}
 \left|S_e(z)-S_e\bigl(\min(z,\sqrt{X/e})\bigr)\right|
 \ll_r A\bigl(X^{-c\delta}+X^{-\delta}\bigr)=o(A/L).
\]
\end{lemma}
\begin{proof}
There is nothing to prove when $z\le\sqrt{X/e}$.
Otherwise a newly admitted cofactor satisfies
$P^-(m)>\sqrt{X/e}$ and $m<X/e$, so it must be a prime;
the cofactor $1$ occurs in both sums.
Every added product is at most $ez\le X^{1-\delta}$.
An integer has at most $(r+1)!$ representations by these marked
primes and the prime cofactor. The bound follows from
$|w(n)|\le1_{\mathcal A_X}(n)+\gamma_dA/X$ and
Lemma \ref{arith:small}.
\end{proof}

\begin{proposition}[Comparison for rough numbers]\label{arith:rough}
Assume \eqref{arith:window}, the comparison in Proposition
\ref{arith:polytope}, the hypotheses of Proposition \ref{arith:typeI},
and $1-\beta<\sigma$. Let $\eta=\beta-\alpha$.
For a fixed number of marked primes $p_i\ge X^\eta$, impose any
fixed Boolean combination $\mathcal R$ of affine conditions in
the variables $\log p_i/\log X$, depending only on the marked primes,
and require
$e=\prod p_i\le X^{1-\alpha}$. Then
\begin{equation}\label{arith:roughformula}
 \sum_{\boldsymbol p}^{\mathcal R}S_e(X^\eta)=o(A/L).
\end{equation}
This includes the unmarked case $e=1$.
For the direct Type II version, let $\epsilon>0$ be any fixed
number and require only $p_i\ge X^\epsilon$.
If a nonempty subproduct of the marked primes has logarithmic
exponent in $\mathcal W$, then
$\sum_{\boldsymbol p}^{\mathcal R}S_e(z)=o(A/L)$ for a threshold
$z\ge X^\epsilon$ equal to a fixed power of $X$ or to one of the
marked primes. Either $P^-(m)>z$ or $P^-(m)\ge z$ may be used.
If the selected subset is proper, no bound on the entire marked product
is needed; if it is the whole marked product $e$, the window condition
already implies $e\le X^{1-\alpha}$.
\end{proposition}
\begin{proof}
\smallskip\noindent\emph{Direct comparison inside the Type II window.}\quad
We first prove the final assertion.
With a fixed lower bound $X^\epsilon$ on all prime factors, expand
the rough cofactor into its unique nondecreasing prime factorization.
The total number of factors is bounded.
Ordering, the least-prime-factor restriction and the original
conditions all give finitely many affine conditions.
If a proper subset of the marked primes is already available, it
remains a proper subset of the full tuple even when the cofactor is
$1$, so Proposition \ref{arith:polytope} applies.
If the available factor is the entire marked product
$e\le X^{1-\alpha}$, the mass of the terms with cofactor $1$ is
$O(A(X^{-c\alpha}+X^{-\alpha}))$ by \eqref{arith:small-bound}.
After removing these terms, the same proposition applies again.

\smallskip\noindent\emph{The fundamental-lemma estimate.}\quad
Let $\tau_0=1-\beta$, and fix $h>0$ with $\tau_0+h<\sigma$.
Choose a sufficiently small fixed $0<\epsilon<\min(\eta,h)$ and
put $z_0=X^\epsilon$.
We use the fundamental lemma of the one-dimensional sieve
\cite[Section~4.2]{Harman2007} in the following form.
For sieve density $g(p)=1/p$, level $D=X^h$ and sifting limit $z_0$,
put $\Pi(z_0)=\prod_{p\le z_0,p\ne\ell}p$.
There are weights $\lambda_f^\pm$ of absolute value at most $1$,
supported on squarefree $f\mid\Pi(z_0)$ with $f\le D$, such that
\[
 \sum_{f\mid m}\lambda_f^-
 \le 1_{(m,\Pi(z_0))=1}
 \le\sum_{f\mid m}\lambda_f^+,\qquad
 \sum_f\frac{\lambda_f^\pm}{f}
 =V(z_0)\{1+O(e^{-c_0h/\epsilon})\},
\]
where $V(z_0)=\prod_{p\le z_0,p\ne\ell}(1-1/p)$.
The constants are uniform for sufficiently small $\epsilon$.
This standard exponential weakening follows immediately from the
stronger error in the fundamental lemma as a function of
$h/\epsilon$.

Apply these upper and lower weights for $e\le X^{\tau_0}$ with
$P^-(e)>z_0$, restricting the cofactor to $(m,b)=1$.
This restriction, together with coprimality to $\Pi(z_0)$, is
exactly $P^-(m)>z_0$ for sufficiently large $X$.
The digital main term on the modulus $ef$, with unit final digit, is
$u_dA/((b-1)ef)$; the reference main term is
$\gamma_d(A/X)\varphi(b)X/(bef)$, which is the same.
To see the uniformity in $e$, write
\[
 R(q)=\#\{n\in\mathcal A_X:(n,b)=1,\ q\mid n\}
       -\frac{u_dA}{(b-1)q}.
\]
The two sieve inequalities give
\[
 |S_e(z_0)|\ll
 \frac{A}{e}V(z_0)e^{-c_0h/\epsilon}
 +\sum_{\substack{f\mid\Pi(z_0)\\f\le D}}|R(ef)|
 +\frac AX\sum_{\substack{f\mid\Pi(z_0)\\f\le D}}1.
\]
The last term accounts for the $O(1)$ error in counting integers
coprime to $b$ in each progression.
By \eqref{arith:typeI-unit}, the remainders can be summed over
$ef\le X^{\tau_0+h}$.
The factorization of each $ef$ is unique: the prime factors of $e$
exceed $z_0$, whereas those of $f$ are at most $z_0$. Thus no additional
representation multiplicity arises.
The total rounding error for the reference integers is
$O((A/X)X^{\tau_0+h})$.
Mertens' estimate also gives
\[
 \prod_{p\le X^\epsilon,p\ne\ell}(1-1/p)\ll(\epsilon L)^{-1},
 \qquad \sum_{e\le X^{\tau_0},P^-(e)>X^\epsilon}\frac1e
       \ll\epsilon^{-1}.
\]
Consequently,
\begin{equation}\label{arith:fundamental}
 \sum_{e\le X^{\tau_0},P^-(e)>z_0}|S_e(z_0)|
 \ll \frac AL\epsilon^{-2}e^{-c_0h/\epsilon}
           +O_{\epsilon,h,J}(AL^{-J}).
\end{equation}

\smallskip\noindent\emph{Stopping at the dual window.}\quad
For terms whose marked product already lies in
$[X^{1-\beta},X^{1-\alpha}]$, apply the Type II comparison for rough
sums from the first paragraph of the proof.
For the remaining marked products, start with $e<X^{1-\beta}$ and
iterate the exact identities
\begin{align}
 S_e(X^\eta)&=S_e(z_0)-\sum_{z_0<p\le X^\eta}S_{ep}(p^-),\label{arith:stop1}\\
 S_e(p^-)&=S_e(z_0)-\sum_{z_0<q<p}S_{eq}(q^-).\label{arith:stop2}
\end{align}
Continue expanding while the current product is less than
$X^{1-\beta}$, and stop at the first crossing.
Every newly introduced prime is at most $X^\eta$, so the terminal
product lies in
\begin{equation}\label{arith:firstcross}
 X^{1-\beta}\le e<X^{1-\beta+\eta}=X^{1-\alpha}.
\end{equation}
The new primes strictly decrease and all exceed $z_0$, so the
expansion is finite.
Every nonterminal term is controlled by \eqref{arith:fundamental}.
Its multiplicity is bounded uniformly in $\epsilon$: the original
marked primes are at least $X^\eta$, and the new primes are at most
$X^\eta$, so this threshold determines the two parts; the decreasing
order of the new primes is then unique.
Equality introduces at most an additional factor $2$, and the
original $r$ markings at most a factor $r!$.
The terminal families satisfy \eqref{arith:firstcross}.
For fixed $\epsilon$, expand the cofactor into a finite prime tuple
and apply the first paragraph of this proof.
The requirement that the parent product has not yet crossed the
threshold, the decreasing order and the other restrictions remain
fixed affine conditions.
Their total error is $o_\epsilon(A/L)$.
It follows that
\[
 \limsup_{X\to\infty}\frac LA
       \left|\sum_{\boldsymbol p}^{\mathcal R}S_e(X^\eta)\right|
 \ll\epsilon^{-2}e^{-c_0h/\epsilon}.
\]
Letting $\epsilon\downarrow0$ proves \eqref{arith:roughformula}.
\end{proof}

The stopping rule is the dual-window version of Maynard's modified
Buchstab decomposition \cite[Section 7, proof of Proposition 6.1]{Maynard2019}.
It explains why the Type I input need only extend a fixed distance
beyond $1-\beta$, even though the marked products in
\eqref{arith:roughformula} may reach $X^{1-\alpha}$.

\subsection{The reference densities}

\begin{lemma}\label{arith:reference}
Fix an integer $r\ge1$ and $\eta>0$, and let $\mathcal R$ be a combination of
finitely many fixed affine polytopes in
$u_i\ge\eta$, $\sum_i u_i\le1$, where
$u_i=\log p_i/L$ and $e=\prod_i p_i$.
Write $\Phi^*(Y,z)=\#\{2\le m<Y:P^-(m)>z\}$.
Define Buchstab's function by
\[
 \omega(u)=0\ (u<1),\quad \omega(u)=1/u\ (1\le u\le2),\quad
 (u\omega(u))'=\omega(u-1)\ (u>2).
\]
Put $v=1-\sum_i u_i$ and $\mathcal R_+=\mathcal R\cap\{v\ge u_r\}$.
The single-prime cofactor count satisfies
\begin{equation}\label{arith:prime-density}
 \sum_{\boldsymbol p}^{\mathcal R}\#\{p\text{ prime}:p_r<p<X/e\}
 =\frac XL\int_{\mathcal R_+}\frac{d\boldsymbol u}{u_1\cdots u_r v}
   +o(X/L).
\end{equation}
The full rough sum also includes the cofactor $1$.
If its domain ensures $X/e\ge p_r\ge X^\eta$, this additional
contribution is $O_r(X^{1-\eta})$.
For the unmarked case and every fixed $0<\eta<1$, the same argument gives
\begin{equation}\label{arith:unmarked-reference}
 \Phi^*(X,X^\eta)
 =\bigl(\omega(1/\eta)+o(1)\bigr)\frac{X}{\eta\log X}.
\end{equation}
\end{lemma}
\begin{proof}
For fixed $U_0$, the bound
$\Phi^*(Y,z)\ll_{U_0}Y/\log z$ holds when $z\le Y\le z^{U_0}$.
To obtain this uniform bound, write each counted integer as a product
of at most $\lfloor U_0\rfloor$ prime factors, use reciprocal sums for
the preceding factors, and apply Chebyshev's bound to the last.
The prime number theorem and recursion on the least prime factor give
\[
 \Phi^*(Y,z)=\frac Y{\log z}
       \{\omega(\log Y/\log z)+o(1)\}
\]
uniformly for $1+\epsilon\le\log Y/\log z\le U_0$.
More explicitly, for $u\le2$ only primes occur; for larger $u$, the
recursion
\[
 \Phi^*(Y,z)=\pi(Y)-\pi(z)
       +\sum_{z<p\le\sqrt Y}\Phi^*(Y/p,p)+O(Y/z)
\]
reduces the subparameter to at most $u-1$.
Partial summation turns the main term into
$u^{-1}(1+\int_1^{u-1}\omega(v)dv)$, which is precisely $\omega(u)$.
First control the part with subparameter close to $1$ by the positive
upper bound above, and then let its fixed width tend to zero.
Induction over finitely many unit intervals proves the required
uniformity.

The prime number theorem also gives weak convergence of the finite
measures
\[
 \sum_{X^\eta\le p\le X}\frac1p\,
          \delta_{\log p/L}\ \longrightarrow\ \frac{du}{u}
 \quad\hbox{on }[\eta,1].
\]
On $1-\sum u_i\ge u_r+\epsilon$, apply the prime number theorem to
the cofactor count $\pi(X/e)-\pi(p_r)$. Its starting point contributes
$o(X/L)$ after reciprocal prime summation, since the ratio $p_re/X$
saves a fixed power on this region. The strip of width $\epsilon$
next to the boundary contributes $O(\epsilon X/L)+o_\epsilon(X/L)$,
by Chebyshev's bound on the cofactor and weak convergence of the
reciprocal prime measures. This gives \eqref{arith:prime-density}.
Fixed affine boundaries have zero limiting measure, and the reference
count involving repeated primes is at most
$\sum_{p\ge X^\eta}X/p^2\ll X^{1-\eta}$.
Finally, the cofactor $1$ corresponds to a marked product
$e\le X^{1-\eta}$, whose total marked count is at most
$r!X^{1-\eta}$.
\end{proof}

\section{A signed Buchstab decomposition and its positive remainder}\label{sieve:section}

Maynard's nine-leaf Buchstab decomposition \cite[Section 6]{Maynard2019}
expresses the prime count as its expected main term plus a sum of signed
remainders. We use the identity in~\eqref{sieve:signed}, with each
$\mathcal L_j$ a positive sum of labelled prime-factor occurrences.
An occurrence records its leaf and its marked primes, so different
markings of the same integer retain their individual contributions.
In Section~\ref{nine:sieve-section} we factor every cofactor and apply
Type II to eligible full tuples before using positivity on the
remaining occurrences.

We give the complete decomposition for the parameter range needed here.
Take
\begin{equation}\label{sieve:parameters}
 \frac14<\alpha<\frac13<\beta<\frac12,
 \qquad 2\alpha+\beta>1,
 \qquad \eta=\beta-\alpha.
\end{equation}
These imply $\eta<\alpha$ and $5\alpha/2<1$.
Write $A_*=1-\alpha$, $B_*=1-\beta$, and
\[
 z_1=X^\eta,\quad z_2=X^\alpha,\quad z_3=X^\beta,
 \quad z_4=X^{1/2},\quad z_5=X^{B_*},\quad z_6=X^{A_*}.
\]
Here $A_*,B_*$ are sieve endpoints, not cardinalities of digital sets.

\subsection{The nine labelled prime-factor regions}

\begin{definition}\label{sieve:domains}
For four variables, let $\mathcal P(u,v,w,t)$ mean that none of their
six pairwise sums belongs to $[\alpha,\beta]$.
Define the following open regions; their boundaries do not affect
the integrals:
\begin{align*}
 D_1={}&\{\eta<v<u<\alpha, v<(1-u)/2, A_*<u+v\},\\
 D_2={}&\{\eta<v<u<\alpha, \beta<u+v<B_*,\ A_*<u+2v<1\},\\
 D_3={}&\{(u,v)\in D_2, v<w<(1-u-v)/2, v+w<\alpha\},\\
 D_4={}&\{(u,v)\in D_2, v<w<(1-u-v)/2, v+w>\beta\},\\
 D_5={}&\{\eta<t<w<v<u<\alpha, \beta<u+v<B_*,\ u+2v<A_*,\\
       &\hspace{22mm}u+v+2w<1, u+v+w+2t<1, \mathcal P(u,v,w,t)\},\\
 D_6={}&\{\eta<t<w<v<u<\alpha, u+v<\alpha,\\
       &\hspace{12mm}u+v+w+t\notin[\alpha,\beta]\cup[B_*,A_*],\\
       &\hspace{12mm}\hbox{all four three-element subsums lie outside }[\alpha,\beta]\},\\
 D_7={}&\{\beta<u<1/2, \eta<v<(1-u)/2, A_*<u+v\},\\
 D_8={}&\{\beta<u<1/2, \eta<v<(1-u)/2,\ \beta<u+v<B_*,\ A_*<u+2v\},\\
 D_9={}&\{\beta<u<1/2, \eta<t<w<v,\ \beta<u+v<B_*,\ u+2v<A_*,\\
       &\hspace{22mm}u+v+2w<1, u+v+w+2t<1,\ \mathcal P(u,v,w,t)\}.
\end{align*}
\end{definition}

\paragraph{The signed leaves.}
To record the decomposition before positivity is used, let
$\widetilde D_2,\widetilde D_8$ be obtained from $D_2,D_8$ by replacing
$A_*<u+2v$ with $A_*\le u+2v$; make the same replacement in the
inherited $D_2$ condition of $D_3,D_4$. Put
$\widetilde D_j=D_j$ for the other indices.
These half-open conventions assign the splitting boundary to its
upper branch. Boundaries of the closed Type II windows belong to
the evaluated pieces.
For $j\ne2$, let $r_j$ be the dimension of $D_j$, put
$e=p_1\cdots p_{r_j}$, and define the linear functional
\[
 \mathcal L_j(f)=
 \sum_{\substack{\boldsymbol p\\
       (\log p_1/L,\ldots,\log p_{r_j}/L)\in\widetilde D_j}}
 \ \sum_{\substack{2\le m<X/e\\P^-(m)>p_{r_j}}}f(em).
\]
All outer variables are primes. For the prime-cofactor leaf, define
\[
 \mathcal L_2(f)=
 \sum_{\substack{p,q\\(\log p/L,\log q/L)\in\widetilde D_2}}
       \ \sum_{\substack{q<s<X/(pq)\\s\ {\rm prime}}}f(pqs).
\]
Every admissible marking is counted. In particular, these
functionals are sums of labelled occurrences, not indicators of
unions of integers.

\begin{proposition}[The signed Buchstab identity]\label{sieve:nine}
Assume \eqref{sieve:parameters} and the hypotheses of Propositions
\ref{arith:polytope} and \ref{arith:rough} and Lemma
\ref{arith:squarefree}. Then
\begin{equation}\label{sieve:signed}
 \#\{p\in\mathcal A_X:p\text{ is prime}\}
 =\gamma_d\frac A{\log X}+\sum_{j=1}^9\mathcal L_j(w)+o(A/\log X),
\end{equation}
The decomposition produces nine positive leaves.
Distinct leaves retain their individual multiplicities even when
they represent the same integer.
The regions $D_1,D_3$ are empty.
\end{proposition}
\begin{proof}
Lemma \ref{arith:squarefree} allows us first to remove repeated prime
factors from all relevant rough integers, at a cost of $o(A/L)$.
We may therefore order the primes with strict inequalities below.
The small-product bound and the prime number theorem give
\begin{equation}\label{sieve:prime-reference}
 \#\{p\in\mathcal A_X\}=\gamma_d A/L+S_1(z_4)+o(A/L).
\end{equation}
The contributions of $1$ and of primes at most $\sqrt X$ are
$O(1+X^{c/2}+AX^{-1/2})$, already included in the error.

\smallskip\noindent\emph{The first expansion.}\quad
Buchstab's identity gives
$S_1(z_4)=S_1(z_1)-\sum_{z_1<p<z_4}S_p(p)$.
The first sum is evaluated by Proposition \ref{arith:rough}, and the
terms with $z_2\le p\le z_3$ by the Type II comparison for rough sums.
In the remaining $S_p(p)$, replacing the threshold by
$\min(p,\sqrt{X/p})$ incurs only a small error.
To prove this, fix $\epsilon>0$.
For $p\le X^{1/2-\epsilon}$, the added cofactor can only be a prime
less than $p$, and the total product is at most $X^{1-2\epsilon}$;
Lemma \ref{arith:threshold} controls these terms.
For $X^{1/2-\epsilon}<p<\sqrt X$, use the positive rough sum with
threshold $z_1$.
Proposition \ref{arith:rough} and the reference upper bound for rough
numbers give $O(\epsilon A/L)+o_\epsilon(A/L)$.
Letting $X\to\infty$ and then $\epsilon\downarrow0$ justifies the
replacement.
Expand once more and evaluate $S_p(z_1)$ to obtain
\begin{equation}\label{sieve:two-branches}
 S_1(z_4)=
 \sum_{z_1<q<p<z_2}S_{pq}(q)
 +\sum_{z_3<p<z_4,\ z_1<q<\sqrt{X/p}}S_{pq}(q)+o(A/L).
\end{equation}
In the first branch, $q<\sqrt{X/p}$ follows automatically from
$3\alpha<1$.
In the second branch, $\beta>1/3$ gives $p>q$, and
$2\alpha+\beta>1$ also gives $q<X^\alpha$.

\smallskip\noindent\emph{Six leaves with a small first prime.}\quad
In the first branch, evaluate the part where
$pq$ lies in $[z_2,z_3]\cup[z_5,z_6]$ by the Type II comparison for rough
sums. Split the remainder into
\[
 pq>z_6,\qquad z_3<pq<z_5,\qquad pq<z_2.
\]
The first part is already a positive leaf and gives $D_1$.
Split the middle part further according to $pq^2\ge z_6$ or
$pq^2<z_6$.
When $pq^2\ge z_6$, replace the threshold $q$ by
$\min(q,\sqrt{X/pq})$.
The added prime cofactors produce total products at most
$pq^2<X^{3\alpha}$, so Lemma \ref{arith:threshold} gives a cost
$o(A/L)$.
Applying Buchstab's identity upwards gives the prime-cofactor leaf
\[
 S_{pq}(\sqrt{X/pq})
\]
and the positive sum
\[
 \sum_{q<r<\sqrt{X/pq}}S_{pqr}(r).
\]
This gives the domain $\widetilde D_2$.
Its marked product satisfies $pq<X^{1-\beta}$.
Replacing the cutoff $\sqrt{X/pq}$ for the prime cofactor by $q$
changes only products at most
\[
 \sqrt{Xpq}+pq^2<X^{1-\beta/2}+X^{3\alpha}.
\]
Lemma \ref{arith:small} therefore gives exactly the functional
$\mathcal L_2(w)$ above, up to $o(A/L)$ and the cofactor $1$.
Its reference density is \eqref{arith:prime-density}.
In the latter sum, the part with $qr\in[z_2,z_3]$ is evaluable;
the remaining parts $qr<z_2$ and $qr>z_3$ give $D_3,D_4$,
respectively.
The trivial cofactors are bounded together below.

When $z_3<pq<z_5$ and $pq^2<z_6$, expand downwards twice.
The first step is
\[
 S_{pq}(q)=S_{pq}(z_1)-\sum_{z_1<r<q}S_{pqr}(r).
\]
For every $r<q$, we have
\[
 pqr^2\le pq^3\le(pq^2)^{4/3}<X^{4(1-\alpha)/3}<X.
\]
The final strict inequality follows from $\alpha>1/4$.
Thus the support restriction $pqr^2<X$ is automatic; we retain it
in the leaf domains for a uniform description.
The same fixed power saving and Lemma \ref{arith:threshold}
justify replacing $r$ by $\min(r,\sqrt{X/pqr})$.
The second expansion gives
\[
 S_{pq}(z_1)-\sum_{z_1<r<q,\ pqr^2<X}S_{pqr}(z_1)
 +\sum_{z_1<s<r<q,\ pqr^2<X,\ pqrs^2<X}S_{pqrs}(s).
\]
Here the identities are expanded with their original thresholds before
the support restrictions are inserted. A term excluded by $pqrs^2\ge X$
has only the cofactor $1$; these terms are included in the small-product
error estimated below.
The marked products in the first two rough sums are at most
$pq^2<z_6$, so Proposition \ref{arith:rough} evaluates them.
The final sum is positive.
Remove the evaluable parts in which one of the six pairwise
subproducts lies in $[z_2,z_3]$; the remainder is $D_5$.

For the part with $pq<z_2$, make the same two downward expansions.
Since $q<p$, we have
$pq^2\le(pq)^{3/2}<X^{3\alpha/2}<z_6$,
so all rough sums requiring evaluation remain in the permitted
range.
The endpoint replacement here also satisfies
$pqr^2\le pq^3\le(pq)^2<X^{2\alpha}$.
From the final positive leaf, remove the parts where a product of
three marked primes lies in $[z_2,z_3]$, or where the entire product
of four marked primes lies in $[z_2,z_3]\cup[z_5,z_6]$.
This gives $D_6$.
A product of three marked primes cannot reach the dual interval,
since its exponent is less than
$3\alpha/2<1/2<1-\beta$.
Moreover, $u+v<\alpha$ and $t<w<v<u$ imply
$u+v+w+2t<(5/2)\alpha<1$.
Thus the cofactor support condition not written separately in this
leaf holds automatically.

\smallskip\noindent\emph{Three leaves with a central first prime.}\quad
In the second branch of \eqref{sieve:two-branches}, remove the part
with $pq\in[z_5,z_6]$. Split the remaining terms according to
\[
 pq>z_6;\qquad z_3<pq<z_5, pq^2\ge z_6;\qquad
 z_3<pq<z_5, pq^2<z_6
\]
The first two are already positive leaves and give $D_7,D_8$,
respectively.
For the third, expand downwards twice.
The bounds for the marked products and the endpoint replacements
are exactly those of the preceding argument, since $p>q$ and
$pq^2<z_6$ still hold.
Evaluate the intermediate rough sums and remove the parts with
pairwise subproducts in the short Type II window to obtain $D_9$.

\smallskip\noindent\emph{Trivial cofactors and the signed identity.}\quad
The marked product in each remaining leaf is bounded by a fixed
power of $X$ below $1$. More precisely, the bounds are
$X^{2\alpha}$ for $D_1,D_6$, $X^{1-\beta}$ for $D_2$,
$X^{1-\beta/2}$ for $D_3,D_4$, $X^{3/4}$ for $D_7,D_8$,
and $X^{4(1-\alpha)/3}$ for $D_5,D_9$.
For the middle two-prime leaves, for example,
$pq<p\sqrt{X/p}<X^{3/4}$; for the upward three-prime leaves,
$pqr<\sqrt{Xpq}<X^{1-\beta/2}$.
For the downward four-prime leaves, use
$pqrs\le pq^3\le(pq^2)^{4/3}$, or
$pqrs\le(pq)^2$ when $pq<X^\alpha$.
Lemma \ref{arith:small} therefore removes all cofactor-$1$ terms
with an $o(A/L)$ error. An equality at a cofactor-support cutoff
also leaves only such a trivial cofactor.
Restoring nonsquarefree integers in the displayed functionals costs
another $o(A/L)$ by Lemma \ref{arith:squarefree}, since the number
of markings of each rough integer is bounded.
With the half-open splitting convention in $\widetilde D_j$, the
Buchstab expansions now give \eqref{sieve:signed}.
No digital boundary contribution has been discarded in this identity.

Finally, $u+v<2\alpha<1-\alpha$ makes $D_1$ empty.
On $D_2$, we have $2v>1-\alpha-u>1-2\alpha>\alpha$, and hence
$v+w>2v>\alpha$, so $D_3$ is empty as well.
\end{proof}

\section{Base nine: full tuples and a positive sieve margin}\label{nine:sieve-section}

We first verify the Type II interval for all nine alphabets. We then
apply the signed Buchstab identity and bound its remaining loss by an
integral over ordered prime exponents, with a weight for each admissible
marking.

\subsection{Moment bounds and the Type II interval}\label{nine:fourier-section}

\begin{proposition}\label{nine:fourier}
For $b=9$ and every forbidden digit $d=0,\ldots,8$,
equations~\eqref{geo:moments} and~\eqref{geo:fractional} hold with the common choices
\begin{equation}\label{nine:moment-values}
 \sigma=\frac{127}{200},\qquad s=\frac32,\qquad
 g_s=\frac{6307}{40000},\qquad
 \kappa=\frac{6307}{20000},\qquad
 g_{4/3}=\frac{13607}{60000}.
\end{equation}
The first and $3/2$ moments both hold uniformly on grids with an arbitrary real shift.
\end{proposition}

The continuous operator inequalities are proved in
Appendix~\ref{app:nine-moments}. For the common-denominator estimate, the
useful interpolation is directly between the first and $3/2$ moments:
\[
 \sum_a F_X(a/X)^{4/3}
 \le\left(\sum_a F_X(a/X)\right)^{1/3}
      \left(\sum_a F_X(a/X)^{3/2}\right)^{2/3}.
\]
This gives $g_{4/3}=(73/200)/3+2(6307/40000)/3=13607/60000$,
which is the exponent used below.

\begin{proposition}[Choice of parameters in base nine]\label{nine:parameters}
For every forbidden digit in base nine, the common-denominator estimate holds
with $y=633/1000$. A Type II window is
\begin{equation}\label{nine:buffer}
 \alpha_0=\frac{6309}{20000},\qquad
 \beta_0=\frac{7383}{20000}.
\end{equation}
We choose the interior sieve parameters
\begin{equation}\label{nine:sieve-params}
 \alpha=\frac{631}{2000}=.3155,\quad
 \beta=\frac{3691}{10000}=.3691,\quad
 \eta=\beta-\alpha=\frac{67}{1250}=.0536.
\end{equation}
\end{proposition}

\begin{proof}
The integer comparison $8^{500}>9^{473}$ gives $c>473/500$.
The positive margins required by~\eqref{geo:three-powers} and the determinant
condition are, respectively,
\[
 \sigma+c/2-1>\frac{27}{250},\quad
 1-y-3g_{4/3}/2=\frac{1073}{40000},\quad
 2\sigma-2y=\frac1{250},\quad 2-3y=\frac{101}{1000}.
\]
We may take $\delta_0=1/2000$ in~\eqref{geo:primitive-energy}
and $\omega=10^{-6}$ for the divisor loss.
For Theorem~\ref{geo:typeii}, take
$u_0=1477/4000$, $\varepsilon_0=10^{-6}$, and $\xi=10^{-5}$.
The corresponding six positive margins are
\begin{align*}
 1-u_0-2\kappa-2(s-1)\varepsilon_0&=49/10^6,&
 u_0-\beta_0&=1/10000,\\
 y-(1-u_0)&=9/4000,&
 2\alpha_0+u_0-1&=3/20000,\\
 2\alpha_0-2\kappa-(3s-4)\varepsilon_0-\xi&=379/2000000,&
 3\alpha_0+1-6\kappa-6(s-1)\varepsilon_0-2\xi&=54227/10^6.
\end{align*}
Also $1-\beta=.6309<\sigma=.635$, so the rough-number comparison of
Proposition~\ref{arith:rough} applies; for example, one may take $.002$ as the
additional exponent in the sieve level.
The interval $[\alpha,\beta]$ lies strictly inside~\eqref{nine:buffer}.
\end{proof}

\subsection{Completing the signed remainder to full prime tuples}

To use the uniform Type II window of Proposition~\ref{nine:parameters},
we refine the nine-leaf decomposition by factoring each rough cofactor
completely. A proper subproduct involving these additional factors may
lie in the Type II window even when none of the marked subproducts does.
We evaluate all such terms before using positivity. This enlarges the
evaluable region while retaining every marked multiplicity in Maynard's
nine-leaf decomposition \cite[Section 6]{Maynard2019}.
The remaining reference loss is strictly less than one, as proved below.

Let $X=9^k$, $A=8^k$, and define
\[
 \gamma_d=\begin{cases}9/8,&3\mid d,\\15/16,&3\nmid d,
 \end{cases}
 \qquad w(n)=\mathbf1_{\mathcal W(X)}(n)-\gamma_dA/X,
 \quad 1\le n<X.
\]
The factor comes from the mass of units among the allowed terminal digits:
there are $6-\mathbf1_{3\nmid d}$ allowed terminal digits coprime to $3$.
Use the nine positive linear functionals $\mathcal L_j$ from
Proposition~\ref{sieve:nine};
their exact domains are given in Definition~\ref{sieve:domains}.
The parameters satisfy
$1/4<\alpha<1/3<\beta<1/2$ and $2\alpha+\beta>1$,
so the general sieve identity applies.

Let $\Omega_j$ be the occurrences counted by $\mathcal L_j$, including
the choice of marked primes and cofactor, and put
$\Omega=\bigsqcup_{j=1}^9\Omega_j$. If $n(\xi)$ is the integer represented
by an occurrence, then
\[
 \sum_{j=1}^9\mathcal L_j(w)=\sum_{\xi\in\Omega}w(n(\xi)).
\]
Completing a cofactor to its prime factorization changes the description
of an occurrence, while leaving its label and contribution unchanged.

For a concrete example, consider the limiting logarithmic coordinates
\[
 (t_0,t_1,t_2,t_3)=(0.16,0.19,0.27,0.38).
\]
The eighth leaf admits the marking $(u,v)=(0.38,0.16)$: its strict
inequalities include $\beta<u<1/2$, $\beta<u+v<1-\beta$ and
$u+2v>1-\alpha$. The marked exponents $0.16$, $0.38$ and their sum
$0.54$ avoid both the Type II interval and its reflection. After the
cofactor is split, however, the new subproduct has exponent
$t_0+t_1=0.35\in[\alpha,\beta]$. This occurrence can therefore be
evaluated before positivity. All inequalities are strict, so the same
gain holds on an open region of the simplex.

\begin{proposition}[Full-tuple refinement retaining the original multiplicities]\label{nine:tuple-loss}
Put
\[
 W=[\alpha,\beta]\cup[1-\beta,1-\alpha],\qquad
 \Delta_r=\{\eta\le t_0\le\cdots\le t_{r-1},\ \sum t_i=1\}.
\]
Call $\mathbf t\in\Delta_r$ bad if the coordinate sum of every nonempty proper
subset lies outside $W$.
Let $m_r(\mathbf t)$ be the total marked multiplicity of this tuple in the
nine original positive leaves.
Then
\begin{equation}\label{nine:loss-integral}
 \mathcal D=\sum_{r=3}^{18}\int_{\Delta_r}
 \mathbf1_{\rm bad}(\mathbf t)
 \frac{m_r(\mathbf t)}{t_0\cdots t_{r-1}}\,dt_0\cdots dt_{r-2}
\end{equation}
satisfies
\begin{equation}\label{nine:prime-lower-word}
 \#\{p<X:p\in\mathcal W(X)\}
 \ge\left(\gamma_d(1-\mathcal D)+o(1)\right)\frac A{\log X}.
\end{equation}
\end{proposition}

\begin{proof}
Write $L=\log X$. We start from the signed identity
\eqref{sieve:signed}, before the nonnegative digital contribution has
been discarded from any leaf.

\emph{Full factorizations and their labels.}
Every prime factor is at least $X^\eta$, so a full factorization has at
most $\lfloor1/\eta\rfloor=18$ factors. Repeated prime factors have total
signed mass $o(A/L)$ by Lemma~\ref{arith:squarefree}.
The first and third leaves are empty since $3\alpha<1$, and the second
leaf already has three prime factors.
The other signed functionals have at least two marked primes and a
nontrivial cofactor. Their cofactor-$1$ terms were already removed with a
fixed power saving in the derivation of~\eqref{sieve:signed}.
Every remaining occurrence therefore has between three and eighteen
distinct prime factors, and the marking multiplicity at each integer is
bounded by a fixed constant.

Factor each rough cofactor completely and sort the full prime tuple.
An occurrence is labelled by its original leaf and its choice of marked
positions in that sorted tuple. These labels form a disjoint union:
different occurrences of the same integer remain different terms.
All their defining conditions are finite Boolean combinations of fixed
affine inequalities in $\log p_i/L$, including the half-open endpoint
conventions inherited from the Buchstab identity.

\emph{Cancellation before positivity.}
Call a labelled occurrence good if a nonempty proper subset of its full
tuple has fixed-$X$ logarithmic sum in $W$.
Partition the good occurrences by the first available subset in a fixed
order. There are only finitely many leaf labels, numbers of prime factors,
markings, and subsets. Proposition~\ref{arith:polytope} therefore gives
$o(A/L)$ for their combined signed $w$-sum, with every original multiplicity
retained.
Let $\mathscr B_X$ be the set of remaining labelled occurrences, and let
$n(\xi)$ be the integer represented by an occurrence $\xi$.
We have the explicit inequality
\begin{equation}\label{nine:signed-deletion}
 \sum_{j=1}^9\mathcal L_j(w)
 =\sum_{\xi\in\mathscr B_X}w(n(\xi))+o(A/L)
 \ge-\gamma_d\frac AX|\mathscr B_X|+o(A/L).
\end{equation}
In particular, positivity is used only after the evaluable signed terms
have been removed. No digital boundary contribution is dropped in this step.

\emph{The reference limit.}
We claim that
\begin{equation}\label{nine:bad-reference}
 |\mathscr B_X|=(\mathcal D+o(1))\,\frac X L.
\end{equation}
For a full tuple of length $r$, mark its first $r-1$ ordered primes and
take the largest as the final prime cofactor, with the second-largest
prime as its lower threshold.
The prime-cofactor formula of Lemma~\ref{arith:reference},
\eqref{arith:prime-density}, gives density
\[
 \frac1{t_0\cdots t_{r-1}}\,dt_0\cdots dt_{r-2},
 \qquad t_{r-1}=1-\sum_{i<r-1}t_i,
\]
on the ordered simplex.

Here is the limiting argument for conditions involving the final prime.
Fix $\delta>0$ and first discard products below $X^{1-\delta}$ from the
reference count; their bounded marking multiplicity gives
$O(X^{1-\delta})=o(X/L)$.
For the remaining products, replacing the last logarithm by
$1-\sum_{i<r-1}t_i$ changes each affine test by $O(\delta)$.
Apart from the total-product cutoff already used as the upper bound for
the final prime, every boundary arising from the proper-subset tests or
the retained leaf conditions is nonconstant on the simplex, or has a
fixed sign there.
The former boundaries have strips of reference mass
$O_\eta(\delta X/L)+o_\delta(X/L)$ by the same prime-cofactor formula;
the latter are resolved before passage to the limit.
The density is bounded because every coordinate is at least $\eta$.
Thus, away from those strips, both the bad indicator and the marked
multiplicity agree with their simplex values.
Let $X\to\infty$ and then $\delta\downarrow0$.
Summing over the finitely many labels gives~\eqref{nine:bad-reference}.
This passage concerns only the positive reference count and requires no
estimate for digital mass on a bad boundary strip.

Finally insert~\eqref{nine:signed-deletion} and~\eqref{nine:bad-reference}
into~\eqref{sieve:signed}. This proves~\eqref{nine:prime-lower-word}.
\end{proof}

\subsection{The multiplicity function on the ordered simplex}

The weight $m_r$ counts all original markings compatible with a bad tuple.
The upward expansions prescribe the least remaining prime, whereas the
downward expansions retain every eligible marking. The following formula
keeps these two rules explicit.
Put $A_*=1-\alpha$ and $B_*=1-\beta$, and disregard equality boundaries of measure
zero. First, for each $x=1,\ldots,r-1$, put $u=t_x,v=t_0$.
Each of the following conditions contributes one to the multiplicity:
\begin{enumerate}
 \item Second leaf: $r=3$, $u<\alpha$,
 $\beta<u+v<B_*$, and $A_*<u+2v<1$.
 \item Fourth leaf: $r\ge4$, with the three strict inequalities in the preceding
 item. Let $w$ be the smallest coordinate remaining after $u,v$ are removed
 (when $x=1$, this is $t_2$; otherwise it is $t_1$), and require $u+v+2w<1$.
 \item Seventh leaf: $\beta<u<1/2$, $u+2v<1$, and $u+v>A_*$.
 \item Eighth leaf: $\beta<u<1/2$, $u+2v<1$,
 $\beta<u+v<B_*$, and $u+2v>A_*$.
\end{enumerate}
The fourth leaf's condition $v+w>\beta$ holds automatically on bad tuples:
$u+2v>A_*$, $u<\alpha$, and $3\alpha<1$ imply $v+w>\alpha$,
and badness excludes $[\alpha,\beta]$.
The upward Buchstab expansion of the fourth leaf requires $w$ to be exactly
the smallest remaining prime factor.

Next, for $r\ge5$ and every $1\le x<y<z\le r-1$, put
$t=t_0,w=t_x,v=t_y,u=t_z$. Write the common conditions as
\begin{equation}\label{nine:fourmark}
 \beta<u+v<B_*,\qquad u+2v<A_*,\qquad
 u+v+2w<1,\qquad u+v+w+2t<1.
\end{equation}
The fifth leaf requires~\eqref{nine:fourmark} and $u<\alpha$;
the ninth requires~\eqref{nine:fourmark} and $\beta<u<1/2$;
the sixth requires $u<\alpha$ and $u+v<\alpha$.
Every admissible choice of marked factors contributes one.
The original exclusions for subproducts of two, three, or four marked
factors hold automatically on a bad full tuple.
In these downward expansions, $w$ may be any eligible marked factor.
It is therefore necessary to enumerate every $x<y<z$, rather than fixing
$w$ as the second-smallest coordinate.

\subsection{Exact slice volumes and the weighted loss}

The integration region is an ordered simplex. Equal-bin permutations
divide the volume of a rational cell, but leave the number of admissible
markings unchanged. We combine these exact volumes with pointwise upper
bounds for the multiplicity and the prime-density factor.

\begin{proposition}[Certified sieve loss in base nine]\label{nine:certificate}
For the parameters in~\eqref{nine:sieve-params},
\begin{equation}\label{nine:certified-loss}
 \mathcal D\le\frac{923405344688153132}{10^{18}}<\frac{231}{250}.
\end{equation}
\end{proposition}

\begin{proof}
Take $N=512$ and cover the ordered simplex in~\eqref{nine:loss-integral} by
nondecreasing integer tuples $j_0\le\cdots\le j_{r-1}$, where
$j_i\ge\lfloor N\eta\rfloor=27$ and $j_i/N\le t_i\le(j_i+1)/N$.
Set $h=N-\sum_i j_i$; only boxes with $1\le h\le r-1$ have positive measure.
In fractional coordinates $y_i=Nt_i-j_i$, we have $\sum_i y_i=h$.
We use coordinate slice measure $dy_0\cdots dy_{r-2}$, with the last
coordinate eliminated by this equality. Equivalently, this is Euclidean
surface measure divided by $\sqrt r$, so it is invariant under coordinate
permutations.
Before imposing any ordering among coordinates in the same grid interval,
the unit-box slice has measure
\begin{equation}\label{nine:eulerian-volume}
 \frac{E(r-1,h-1)}{(r-1)!},\qquad
 E(n,m)=(n-m)E(n-1,m-1)+(m+1)E(n-1,m),\quad E(0,0)=1.
\end{equation}
This is the density at the integer $h$ of the sum of $r$ independent uniform
variables on $[0,1]$.
The recurrence is used for $n\ge1$ and $0\le m<n$.
All entries outside this range are zero, except for the initial entry
$E(0,0)=1$.
Inclusion--exclusion on the cube gives
\[
 \frac1{(r-1)!}\sum_{q=0}^{h}(-1)^q\binom rq(h-q)^{r-1},
\]
whose numerator is the Eulerian number above.
If a grid value $j$ occurs $\nu_j$ times, imposing the ordering divides the
volume by exactly $\prod_j\nu_j!$: permutations within an equal-grid-value
group preserve the slice, and their ordered sectors have disjoint interiors
and equal measure.
Thus the ordered box has measure, in the original coordinates,
\begin{equation}\label{nine:ordered-volume}
 \frac{N^{-(r-1)}E(r-1,h-1)}{(r-1)!\prod_j\nu_j!}.
\end{equation}

For a fixed subset of size $s$ whose grid values sum to $J$, the exact range
of its coordinate sum on the \emph{unordered} containing slice is
\begin{equation}\label{nine:subset-range}
 \left[\frac{J+\max(0,h-r+s)}N,
       \frac{J+\min(h,s)}N\right].
\end{equation}
Within-bin ordering can only shrink this range.
If the displayed interval lies in one component of $W$, the same fixed
subset certifies every tuple in the ordered cell as good, so the cell is deleted.
During recursion, the larger interval $[J/N,(J+s)/N]$ also permits safe early
deletion.
Integer bitsets indexed by subset size record all attainable sums of grid
values; when a grid value $v$ is added, update
$T_s\gets T_s\cup(T_{s-1}+v)$.
The bitset records existence of an index subset with each integer sum;
the interval bound depends only on that sum and the subset size.
The subset search thus uses only finite exact integer arithmetic.
All window endpoints are written with denominator $10000$, and every interval
test is performed by integer cross-multiplication.

For each retained box, replace each leaf condition from the preceding
subsection by a necessary condition for intersection.
Consider an affine sum involving $s$ distinct coordinates, one with
coefficient $2$, and let its integer part be $J_2$.
On the unordered containing slice, the numerators of its lower and upper
endpoints are
\begin{equation}\label{nine:double-range}
 J_2+\max(0,h-r+s),\qquad J_2+\min(h,s)+1.
\end{equation}
Since $1\le h\le r-1$, the minimum can be attained without placing fractional
mass in the coordinate with coefficient two; for the maximum, place one unit
in that coordinate first. This proves the formula.
For a constraint $L<c$, retain the box only when $\min L<c$;
for $L>c$, retain it only when $\max L>c$.
Retaining every potentially admissible marking gives an integer upper bound
$M_r(\mathbf j,h)$ for $m_r(\mathbf t)$ at every bad point in the ordered cell.
Even if the relaxed conditions cannot hold simultaneously, the resulting
value can only increase the upper bound.
Also $\prod t_i^{-1}\le N^r/\prod j_i$. We integrate this pointwise bound
over the whole ordered cell, including any good part of a cell not deleted.
Only the cell volume is divided by the equal-bin factorials; no marking
multiplicity is divided or averaged. The resulting contribution to
\eqref{nine:loss-integral} is at most
\begin{equation}\label{nine:cell-bound}
 \frac{N E(r-1,h-1)M_r(\mathbf j,h)}
 {(r-1)!\prod_i j_i\prod_j\nu_j!}.
\end{equation}
Every first-layer box crossing $\eta$ is retained, so the boundary is covered
from the outside as well.

For each $3\leqslant r\leqslant18$, the program
\path{supplement/base9/certify_full_tuple_sieve.cpp} enumerates every nondecreasing tuple
of grid values.
If the partial sum is $T$ and $r-d$ coordinates remain, the next grid
value is at most $\lfloor(N-1-T)/(r-d)\rfloor$.
When only the last coordinate remains, $h\le r-1$ also gives the lower bound
$N-r+1-T$.
The enumeration reaches each positive-volume ordered box once unless
an earlier subset test discards it. Completed tuples outside
$1\le h\le r-1$ have zero contribution.

Use $F=10^{18}$ as a common denominator.
To bound $F$ times~\eqref{nine:cell-bound}, start with $FN E(r-1,h-1)$,
divide successively by $(r-1)!$, the $j_i$, and the $\nu_j!$, rounding upward
after every division, and finally multiply by $M_r$.
Every step preserves an upper enclosure.
The program uses unsigned $128$-bit integers with overflow assertions:
$FN\,17!<2^{128}$ and $23!<2^{128}$.
The recursion counters for the specified run also remain within $64$ bits.
Running with assertions enabled gives the bounds in Table~\ref{nine:loss-table}.
Their exact sum is the numerator in~\eqref{nine:certified-loss}.
Since $19\eta>1$, there is no uncomputed tail in higher dimensions.
\end{proof}

\begin{table}[tbp]
 \centering
 \caption{Rigorous upper bounds for the base-nine sieve loss, grouped by the
 number of prime factors in the full tuple. Every entry has denominator $10^{18}$.}
 \label{nine:loss-table}
 \begin{tabular}{rr@{\qquad}rr}
 \toprule
 $r$&Upper bound&$r$&Upper bound\\\midrule
 3&$.676842122666188091$&11&$.000000000007238880$\\
 4&$.098795867909592915$&12&$.000000000003166020$\\
 5&$.129047629062754398$&13&$.000000000001492040$\\
 6&$.018176369591635255$&14&$.000000000000453596$\\
 7&$.000535666342176341$&15&$.000000000000252616$\\
 8&$.000007612230320382$&16&$.000000000000115570$\\
 9&$.000000076343360570$&17&$0$\\
 10&$.000000000529406458$&18&$0$\\\midrule
 \multicolumn{4}{c}{Total $\mathcal D\le .923405344688153132$}\\\bottomrule
 \end{tabular}
\end{table}

\subsection{Prime counts at arbitrary real endpoints}

\begin{theorem}\label{nine:theorem}
For each $d\in\{0,\ldots,8\}$, using the usual base-nine expansion without
leading zeros, we have
\[
 \pi_{9,d}(x)\asymp\frac{A_{9,d}(x)}{\log x}
 \asymp\frac{x^{\log_9 8}}{\log x}.
\]
More explicitly, for all sufficiently large real $x$, uniformly in the
forbidden digit,
\begin{equation}\label{nine:explicit}
 \frac1{320}\frac{x^{\log_9 8}}{\log x}
 \le\pi_{9,d}(x)\le400\frac{x^{\log_9 8}}{\log x}.
\end{equation}
\end{theorem}

\begin{proof}
Propositions~\ref{nine:fourier} and~\ref{nine:parameters}, together with the
general arithmetic comparisons, supply all the hypotheses of
Proposition~\ref{nine:tuple-loss}.
By Proposition~\ref{nine:certificate},
$1-\mathcal D>19/250$ and $\min_d\gamma_d=15/16$.
Hence every sufficiently large $X=9^k$ satisfies
\[
 \#\{p<X:p\in\mathcal W(X)\}\ge\frac1{40}\frac{8^k}{\log X}.
\]
For the upper bound, apply Proposition~\ref{arith:rough} with no marked factors.
Apart from primes $p\le X^\eta$, all primes are $X^\eta$-rough, and the
reference count $1+\Phi^*(X,X^\eta)$ is
$(\omega(1/\eta)+o(1))X/(\eta\log X)$ by \eqref{arith:unmarked-reference}.
Since $0\le\omega\le1$, $\gamma_d\le9/8$, and $(9/8)/\eta<21$,
the number of primes among the complete words is at most
$22\,8^k/\log X$ for sufficiently large $X$.
The small primes are absorbed by the digital counting bound $O(X^{c\eta})$.

Take $9^k\le x<9^{k+1}$.
Every prime among the complete words of length $k$ is counted by
$\pi_{9,d}(x)$.
The lower bound therefore loses at most a factor $8$, giving the lower
inequality in~\eqref{nine:explicit}.
For the upper bound, sum over all possible word lengths.
This remains an upper bound even when $d\ne0$ introduces repetitions.
Using
$\sum_{j\le K}8^j/j=(8/7+o(1))8^K/K$
with $K=k+1$ gives the stated constant $400$.
Finally, counting complete words and summing the geometric series gives
$x^c/16\le A_{9,d}(x)\le10x^c$ for sufficiently large $x$.
This proves both order-of-magnitude comparisons.
\end{proof}

\section{Octal divisor estimates and centered comparison}
\label{sec:octal-typeI}

For the octal argument put $D_d=\{0,\ldots,7\}\setminus\{d\}$,
$X=8^k$, $L=\log X$, $A=7^k$, $\rho=A/X$, and
$\delta=1-\log_8 7$. Let $\mathcal A_X$ be the length-$k$ words
over $D_d$, identified with integers in $[0,X)$. Put $e(t)=e^{2\pi i t}$
and write the one-digit and word amplitudes as
\[
 f_d(t)=\frac17\left|\sum_{a\in D_d}e(at)\right|,
 \qquad
 F_X(t)=\frac1A\left|\sum_{n\in\mathcal A_X}e(nt)\right|
       =\prod_{j=0}^{k-1}f_d(8^jt).
\]
We use the normalized transfer operator
$\mathcal T_d g(u)=\tfrac18\sum_{j=0}^7 f_d((u+j)/8)g((u+j)/8)$.
Fix any $d\in\{0,\ldots,7\}$. Put $K_d=D_d\cap\{1,3,5,7\}$ and
write $\gamma_d=2|K_d|/7$. Thus $\gamma_d=8/7$ if $d$ is even and
$6/7$ if $d$ is odd. We work with $f=1_{\mathcal A_X}1_{\rm odd}$,
of mass $\gamma_d A/2$ on a full period. If $d=0$, every word in
$\mathcal A_X$ has exactly $k$ ordinary digits. If $d\ne0$, padding
with zeros represents every allowed integer in $[1,X)$ in $\mathcal A_X$.
On $\mathbb Z/X\mathbb Z$, let $E_{X,Q}f$ denote the conditional mean
of $f$ on classes modulo a dyadic divisor $Q$ of $X$.

\begin{lemma}[Uniform Type I input]\label{octal:uniformTypeI}
For every missing octal digit, the hypotheses of Proposition~\ref{arith:typeI}
hold with $\sigma=247/400$.
\end{lemma}
\begin{proof}
The supplement supplies positive periodic piecewise linear functions
$g_d$ with $4096$ rational nodes each (Appendix~\ref{app:reproduce}).
They satisfy
\[
\mathcal T_dg_d\le (2769/10000)g_d\qquad (d=0,1,2,3).
\]
On the refined grid of $n=16384$ nodes the verifier encloses cosines
using Machin's identity and the degree-$48$ Taylor polynomial with
remainder bounded by $4^{50}/50!$. The identity
\[
 49 f_d(t)^2=7+2\sum_{j=1}^7 C_j\cos(2\pi jt),\qquad
 C_j=8-j-1_{d+j\le 7}-1_{d-j\ge0},
\]
has nonnegative coefficients. Upper enclosures for the cosines and an
upward-rounded square root thus give an upper residual at each node.
Writing $\lambda=2769/10000$, the residual
$\mathcal T_d g_d-\lambda g_d$ has Lipschitz constant at most
\[
 \tfrac18\big(\operatorname{Lip}(f_d)\max g_d+\operatorname{Lip}(g_d)\big)
       +\lambda\operatorname{Lip}(g_d).
\]
Here $\operatorname{Lip}(f_d)\le 32\pi/7<704/49$ by centering the
polynomial at $7/2$, and $\operatorname{Lip}(g_d)$ is computed from its
linear segments. The verifier adds this constant divided by $2n$ to
the largest upper nodal residual, and checks that the result is negative.
The included certificate and verifier prove these inequalities with
integer and rational arithmetic. Reflection gives the other digits.
Since $2769^{400}8^{247}<10000^{400}$ and $g_d$ is bounded above and away
from zero, iteration gives $\mathcal T_d^k1\ll X^{-247/400}$.
Thus the shifted-grid first norm is $O(X^{.3825})$. Integration in
the shift variable gives the continuous first-moment estimate.
\end{proof}


\subsection{Odd divisors and a quadratic Fourier approximation}
\label{sec:projection-typeI}

In this section $X=8^k$, $A=7^k$, $\rho=A/X$, and
$\delta=1-\log_8 7<129/2000$. Write $a=1_{\mathcal A_X}$ on
$\mathbb Z/X\mathbb Z$, and normalize the Fourier transform by
$\widehat a(b)=X^{-1}\sum_{n<X}a(n)e(-bn/X)$.


\begin{lemma}[A moment near Parseval]\label{octal:near-parseval}
Uniformly in the omitted digit and the real shift $v$,
\[
 \sum_{b<X}F_X((b+v)/X)^{127/64}
 \ll (11481/10000)^k\ll X^{\delta+141/64000}.
\]
\end{lemma}
\begin{proof}
Let $s=127/64$ and let $T_s g(u)=\sum_{j=0}^7 f_d((u+j)/8)^sg((u+j)/8)$.
The data in \path{supplement/base8/near_parseval_nodes.json} specify a positive piecewise
linear periodic function $g_d$, for each $d=0,1,2,3$, with $256$ rational
nodes. The verifier proves $T_s g_d\le\lambda g_d$ everywhere, where
$\lambda=11481/10000$. Here are the finite enclosures and the guard used.
On the refined grid of $32768$ nodes, bound $f_d^2$ above using the
cosine expansion in Lemma~\ref{octal:uniformTypeI} and directed rational
enclosures of cosines. If $u$ is an integer upper bound for $Sf_d^2$,
with $S=2^{90}$, the least nonnegative integer $w$ with
$w^{128}\ge S u^{127}$ is an upper bound for $Sf_d^s$.
The verifier finds it by directed integer square roots.
The rational nodal residual uses this upper weight and the exact linear
interpolation of $g_d$.

For the extension from nodes, $\operatorname{Lip}(f_d)<15$ by centering
at $7/2$. On $0\le x\le1$, the function $x^s$ has Lipschitz constant
at most $2$, and $x^s-x^2$ has Lipschitz constant at most $2(2-s)$.
For the latter, differentiate $t x^{t-1}$ in $t\in[s,2]$ and use
$\sup_{0<x\le1}x^{s-1}|\log x|=1/(e(s-1))$. As $T_2 1=8/7$ by
orthogonality, $T_s 1$ has Lipschitz constant at most $30(2-s)$.
For $T_s(g_d-1)$ use the branch scale $1/8$ and sum over eight branches.
The full residual $T_sg_d-\lambda g_d$ has Lipschitz constant at most
\[
 30(2-s)+30\|g_d-1\|_\infty+(1+\lambda)\operatorname{Lip}(g_d).
\]
After adding the guard of this constant divided by $2\cdot32768$,
the largest upper nodal residual is negative for each digit.
Reflection gives $d=4,5,6,7$. Iteration gives the first bound since
$T_s^k1(v)$ is the shifted grid moment and each $g_d$ is bounded above
and away from zero. The second bound follows from the integer check
$(7\cdot11481)^{64000}<8^{141}(8\cdot10000)^{64000}$.
\end{proof}

\begin{lemma}[Quadratic approximation]\label{octal:quad-approx}
Put $f(n)=a(n)1_{n\ \mathrm{odd}}$. There is a periodic function $G$ such that
\[
 \sum_{1\le n<X}|f(n)-G(n)|\ll A X^{-1/128000},\qquad
 \sum_{b<X}|\widehat G(b)|\le \rho X^{697/2000}.
\]
\end{lemma}
\begin{proof}
Keep only the frequencies $b$ with $F_X(b/X)\ge X^{-142/1000}$ in
$a$ and denote the resulting Fourier projection by $P$.
The discarded squared mass of the normalized digital modulus is
\[
 \sum_{b\ \mathrm{discarded}} F_X(b/X)^2
 \ll X^{-142/64000}X^{\delta+141/64000}=X^{\delta-1/64000}.
\]
Parseval gives total squared mass $X^\delta$, so the number of kept
frequencies is at most $X^{\delta+284/1000}<X^{697/2000}$.
As $|\widehat a(b)|=\rho F_X(b/X)$, in probability normalization
$\|P-a\|_2\ll\rho^{1/2}X^{-1/128000}$ and $\|P\|_2\le\|a\|_2=\rho^{1/2}$.
Take $G=1_{\mathrm{odd}}|P|^2$ and use $|a|^2=a$ and Cauchy--Schwarz for
the first bound. For the second, $\|\widehat P\|_1^2\le X^{697/2000}\rho$
by Cauchy--Schwarz on its support. The Fourier transform of $|P|^2$
is a convolution, and multiplication by $1_{\mathrm{odd}}$ costs a
Fourier $\ell^1$ norm of exactly one.
\end{proof}

\begin{theorem}[Type I for octal words and their periodic projections]
\label{octal:lifted-odd-divisor-marks}\label{octal:strong-typeI}
Fix $\zeta>0$, $K,J>0$, and $\theta<.6515$.
Let $X^\zeta\le Q\mid X$ be dyadic and put $F_Q=E_{X,Q}f$.
Then, uniformly in $Q$,
\[
 \sum_{\substack{e\le X^\theta\\e\,\mathrm{odd}}}\tau(e)^K
    \sup_I\left|\sum_{n\in I\cap[1,X)} F_Q(n)
                       (1_{e\mid n}-1/e)\right|\ll A L^{-J}.
\]
Here $I$ ranges over intervals. The case $Q=X$ gives the estimate for $f$.
\end{theorem}

\begin{proof}
Use the Ramanujan expansion of the centered mark. For primitive
reduced denominator $s>1$, its total divisor weight is at most
$O(\tau(s)^K L^{O_K(1)}/s)$ by summing over multiples of $s$.
For $s\le Q^{.49}$ expand
$c_s(n)=\sum_{t\mid s}\mu(s/t)t(1_{t\mid n}-1/t)$.
A second harmonic summation leaves at most a logarithmic power
times the corresponding sum of centered marks at odd $t\le Q^{.49}$
with weights $\tau(t)^K$.
Let $Y$ be the least octal power divisible by $Q$. On each
$Y$-block $F_Q$ is a density-scaled average of at most four
translates of the word at length $Y$. Each translated block
splits into at most two intervals at its cyclic endpoint, and
each mark specifies one residue class on each interval.
Slice the terminal odd digit into its allowed classes and use
the maximum-residue estimate \eqref{arith:typeI-max} at
length $Y/8$, which applies since $Q^{.49}<(Y/8)^{1/2}$ for
large $X$. The sum of these errors over the blocks has scale $A$
and arbitrary logarithmic savings.

For larger $s$ use the quadratic approximation of the word
from Lemma~\ref{octal:quad-approx}, then average and lift
from length $Y$ as above. Its counting first-norm difference
from $F_Q$ is $O(AQ^{-\epsilon_0})$ for some fixed $\epsilon_0>0$.
Its probability-normalized Fourier first norm is at most
$O((A/X)Q^{.3485})$, since cyclic translations and passage to
the period projection preserve the bound. Substitution in the large denominator
sum is justified before taking Fourier absolute values by
\[
 \sum_{s\le X^\theta}\frac{\tau(s)^K}{s}|c_s(n)|
    \ll \tau(n)^{K+1}L^{O_K(1)}\qquad (1\le n<X).
\]
This follows by expanding $c_s$ in divisors and summing over
$s$ divisible by each divisor of $n$. The pointwise divisor
bound absorbs the substitution error by a fixed power.

In a dyadic range $s\asymp S$, the distinct primitive fractions
have spacing $\gg S^{-2}$. Summing the geometric bound for the
interval test against any one Fourier frequency costs
$O(X+S^2 L)$ over these fractions. Including the divisor weight
and the Fourier first norm of the approximation, the cost
relative to $A$ is at most
\[
 L^{O(1)} Q^{.3485} (1/S+S/X) S^{\epsilon_1}
\]
for any fixed $\epsilon_1>0$. Both terms save a fixed power
since $S\ge Q^{.49}/2$, $Q\ge X^\zeta$, and
$\theta+.3485<1$. Summing the dyadic ranges proves the lemma.
\end{proof}


\subsection{An upper bound for the complete word set}\label{sec:upper-all}

Let $P_d(X)$ count primes in $\mathcal A_X$. We claim
\begin{equation}\label{octal:word-upper}
 P_d(X)\le \left(\frac{25}{4}\gamma_d+o(1)\right)\frac A L.
\end{equation}
Put $z=X^{8/25}$ and let $\mathcal S$ consist of the odd squarefree
integers at most $z$. With $G=\sum_{m\in\mathcal S}1/\varphi(m)$,
define the Selberg weights by
\[
 y_m=\frac{\mu(m)}{G\varphi(m)},\qquad
 \lambda_a=a\sum_{\substack{m\in\mathcal S\\a\mid m}}\mu(m/a)y_m
 \quad (a\in\mathcal S).
\]
Inversion gives $\sum_{a\in\mathcal S,\,m\mid a}\lambda_a/a=y_m$
and $\lambda_1=1$. Also $|\lambda_a|\le1$. Indeed its absolute value
is $a/(G\varphi(a))$ times the sum of $1/\varphi(v)$ over
$v\in\mathcal S$ with $(v,a)=1$ and $v\le z/a$. The distinct integers
$uv$ for these $v$ and $u\mid a$ lie in $\mathcal S$ and contribute
$a/\varphi(a)$ times that sum to $G$.

For an odd prime greater than $z$, the divisor sum with weights
$\lambda_a$ equals one. Expand its square against $f$ on $[1,X)$.
The modulus $[a,b]$ is odd and at most $z^2=X^{.640}$, with at most
$\tau([a,b])^2$ representations. By Theorem~\ref{octal:strong-typeI}
the total error is $O(AL^{-2})$. The main quadratic form is
\[
 \sum_{a,b\in\mathcal S}\frac{\lambda_a\lambda_b}{[a,b]}
       =\sum_{m\in\mathcal S}\varphi(m)y_m^2=G^{-1},
\]
by the divisor expansion of $(a,b)$. There are at least
\[
 \frac v2-1-\sum_{3\le p\le\sqrt v}\left(\frac{v}{2p^2}+1\right)
       \ge\frac v4-O(\sqrt v+1)
\]
odd squarefree integers at most $v$, since
$\sum_{m\ge3}m^{-2}\le1/2$. Partial summation and
$1/\varphi(m)\ge1/m$ yield $G\ge(1/4)\log z-O(1)$.
Now $P_d(X)\le z+1+\gamma_d A/(2G)+O(AL^{-2})$, proving
\eqref{octal:word-upper}.

\subsection{Periodic projections and the flat reference}
\label{sec:vaughan-signed}

For dyadic $q\mid X$ with $q\ge4$ put
$\mathcal D_q f=E_{X,q}f-E_{X,q/2}f$. Choose a smooth cutoff $\omega$
with $0\le\omega\le1$, equal to one on $[2L^{-B_1},1-L^{-B_1}]$
and supported in $[L^{-B_1},1]$. Take $B_1$ sufficiently large and
fixed. The cutoff has bounded variation and each fixed derivative
is bounded by a power of $L$. An interval of length $H$ meets at
most two aligned octal blocks of length between $H$ and $8H$, each
containing $O(H^{\log_8 7})$ words. The removed intervals therefore
have mass $O(AL^{-B_1\log_8 7})$.
Set
\[
 f_0(n)=\gamma_d\rho\,1_{n\ \mathrm{odd}},\qquad
 b(n)=\omega(n/X)(f(n)-f_0(n)).
\]
Counting the allowed last digits gives $E_{X,2}f=f_0$ and equality
of the full masses. In particular $\sum_n b(n)\ll AL^{-B_1\log_8 7}$.

\subsection{A centered divisor comparison on odd prime-factor ranks}
\label{sec:centered-odd-rank-switch}

Use the notation of Section~\ref{sec:vaughan-signed} with $\theta=.640$.
Take $r=11/20$ and a fixed smooth $J$ on $[0,1]$ with
$J(t)+J(1-t)=1$ and $J(t)=0$ for $t\ge r$. Then $J(t)=1$ for
$t\le9/20$.
For $n>1$ set
\[
 C_J(n)=\sum_{e\mid n}\mu(e)J(\log e/\log n).
\]

\begin{proposition}\label{octal:centered-odd-rank-switch}
The signed comparison takes the form
\begin{equation}\label{eq:centered-odd-rank-identity}
 \sum_{p<X}(\log p)f(p)\omega(p/X)
 =\gamma_d A-
 L\sum_{\substack{n<X,\ n\ \mathrm{squarefree}\\
                       \Omega(n)\ge3\ \mathrm{odd}}} C_J(n)b(n)+o(A).
\end{equation}
The coefficient $C_J(n)$ also vanishes on any squarefree composite in
the sum with a prime divisor at least $n^r$.
\end{proposition}

\begin{proof}
First, for $g=f+f_0$ and any fixed $K\ge0$,
\begin{equation}\label{eq:centered-odd-divisor-moments}
             \sum_{n<X}g(n)\tau(n)^K\ll_K A L^{O_K(1)}.
\end{equation}
Indeed greedily group the prime factors with multiplicity into
consecutive products at most $\sqrt X$. When no prime exceeds
$\sqrt X$, consecutive groups have product exceeding $\sqrt X$,
so there are at most three groups. The divisor count of the
integer is at most the product of those of the groups, hence
$\tau(n)^K\le\sum_{a\mid n,\,a\le\sqrt X}\tau(a)^{3K}$.
If one prime exceeds $\sqrt X$, remove it for the same bound
up to $2^K$. Apply Type~I to these odd divisor marks using
their fixed divisor powers and reciprocal sums.

Choose $\epsilon>0$ with $r+2\epsilon<\theta$.
The mass of integers divisible by some $h^2$ with $h>X^\epsilon$
is $O_M(AL^{-M})$ for arbitrary fixed $M$. Up to $X^{\theta/2}$
this follows from Type~I on $h^2$ and the harmonic square sum;
above that bound use $O(X^{1-\theta/2}+\sqrt X)$ by integer
counting. Cauchy's inequality with
\eqref{eq:centered-odd-divisor-moments} and the pointwise bounds
by $\tau(n)$ on $C_J(n)$ and the inclusion-exclusion error then
permit the replacement
\[
 \mu^2(n)\ \longmapsto\
             \sum_{\substack{h^2\mid n\\h\le X^\epsilon}}\mu(h)
\]
against $L b(n) C_J(n)$ with error $o(A)$.
Expand the latter sum and mark $[e,h^2]\le X^{r+2\epsilon}$.
The multiplicity is bounded by a divisor power. For fixed $e$,
$J(\log e/\log n)$ is a function of bounded total variation
on the product interval. Extend it by zero where $n<e$ for
$e>1$. Partial summation and interval-uniform Type~I replace
each divisor mark by its reciprocal with total error $o(A)$.

In the main term one must keep the entire harmonic coefficient
at each $n$. For fixed odd $h$ the coefficient from the sum
in $e$ is
\[
 \sum_{e\ {\rm odd}}\frac{\mu(e)}{[e,h^2]}J(\log e/\log n)
 =\frac1{h^2}\sum_{a\mid h^2}\frac{\mu(a)\varphi(a)}a
         \sum_{\substack{u\ge1\\(u,2a)=1}}
                  \frac{\mu(u)}u J(\log(au)/\log n).
\]
This follows by expanding $\gcd(e,h^2)$ by $\varphi$.
The coefficient is $O_N(L^{-N})$ uniformly on the product support.
To see this, take $\epsilon<(1-r)/4$ and put $U=X^{(1-r)/4}$.
The factor $\mu(a)$ forces $a$ to be squarefree, so $a\le h\le X^\epsilon$.
Write the restricted M\"obius sequence as $\mu*1_{\{t:\,p\mid t\Rightarrow p\mid 2a\}}$.
For $t\le U$, the taper $J(\log(atm)/\log n)$ equals one for $m\le U$
when $X$ is large, since $atU\le X^{3(1-r)/4}$ and
$\log n=L+O(\log L)$. The usual harmonic M\"obius estimate and partial
summation give $O_N(L^{-N})$ in the $m$-sum. The auxiliary harmonic
mass is $(2a)/\varphi(2a)\ll L$. For $t>U$, the harmonic tail is bounded by
\[
 U^{-1/2}\prod_{p\mid 2a}(1-p^{-1/2})^{-1}
 \ll U^{-1/2}(2a)^{1/8},
\]
where the product bound follows by comparing each sufficiently large
prime factor with $p^{1/8}$. This has a fixed power saving and the
absolute harmonic sum over $m$ costs $O(L)$. Finally
$\sum_h\tau(h^2)/h^2<\infty$, so summation over $a,h$ preserves the
arbitrary logarithmic saving.
Consequently $L\sum_n b(n)\mu^2(n)C_J(n)=o(A)$.

For a squarefree $n>1$ of even rank, complementary divisors
have the same M\"obius sign, so $2C_J(n)=\sum_{e\mid n}\mu(e)=0$.
If $p\ge n^r$ divides a squarefree composite, all divisors
containing $p$ have test value zero, and all divisors of
$n/p$ have test value one. Their signed sum vanishes.
Finally $C_J(p)=1$ for a prime, and the prime number theorem
evaluates the flat prime reference. Replacing the outside
$L$ by $\log p$ on the digital prime terms costs $o(A)$
by the ordinary upper sieve $\sum_p f(p)\ll A/L$ and the
product cutoff. This proves \eqref{eq:centered-odd-rank-identity}.
\end{proof}


\paragraph{A marked upper bound.}
For the small-prime estimates use $g=f$ or $g=f_0$. Let $a$ be a
product of at most three distinct marked prime powers of size at most
$X^{.01}$, and let $3\le z\le X^\eta$, where $\eta<.001$.
The upper main bound for
$\sum_{am<X,\,P^-(m)\ge z}g(am)\tau(m)$ is
\begin{equation}
 \ll \frac{AL}{a(\log z)^2}.
 \label{eq:vaughan-rough-hyperbola}
\end{equation}
Indeed, $\tau(m)\le2\sum_{u\mid m,\,u\le\sqrt m}1$ and every such
$u$ is $z$-rough. For fixed $u$, upper-sieve $m/u$ from odd primes below
$z$ using modulus $au$ and level $X^{.60}/(au)$. Since $au\le X^{.53}$,
the dimension-one upper sieve gives main bound $O(A/(au\log z))$
uniformly for the indicated $z$. The remaining harmonic sum is at most
$\prod_{z\le p\le X}(1-1/p)^{-1}\ll L/\log z$. Without sifting the
same hyperbola argument gives upper main bound $O(AL/a)$.
In all applications of these bounds below, the aggregate error is
$O_J(AL^{-J})$ for arbitrary fixed $J$. A remainder modulus splits into
the marked powers, the hyperbola divisor and the sieve modulus in at
most a fixed power of its divisor function ways. Theorem~\ref{octal:strong-typeI}
absorbs these divisor weights and the logarithmic weights below.
The corresponding errors for $f_0$ follow by counting odd multiples,
whose interval rounding errors are $O(A/X)$ per modulus.


\paragraph{Rough restriction at a small fixed exponent.}
Since $C_J=0$ on squarefree integers of even rank,
we may sum over squarefree integers in this restriction. Supported
primes cannot have a least prime factor at most $X^\eta$.
The contribution to \eqref{eq:centered-odd-rank-identity} from
$P^-(n)\le X^\eta$ is a signed sum whose absolute value, divided by
$A$, has arbitrarily small limsup as $\eta\downarrow0$. For $j$ specified
small prime divisors $p_i$ of the squarefree integer put
$a=\prod_i p_i$, $n=aN$, $t_i=\log_X p_i$ and $y=\log_X n$.
Group the M\"obius terms according to whether they contain each
marked prime. The resulting coefficient at $v\mid N$ is $\mu(v)$
times
\[
 (-1)^j\frac{\prod_i t_i}{y^j}
 \int_{[0,1]^j} J^{(j)}\!\left(
        \frac{\log_X v+\sum_i t_i u_i}{y}\right)\,du .
\]
For $j\le3$ this and its derivative in $y$ are
$O_J(\prod_i t_i)$ on the product support, and the coefficient
vanishes for $v>X^r$.
When there are at least three small primes, mark the three
least. The other factor is $p_3$-rough and the weight is
at most $O_J(L\tau(N)\prod_i t_i)$. Apply the marked hyperbola
bound \eqref{eq:vaughan-rough-hyperbola}, namely
$O(AL/[a(\log p_3)^2])$ without the factor $L\prod_i t_i$.
Summing the logarithmic reciprocal prime weights gives
$O_J(A\eta)+o(A)$. Indeed the upper main bound after summing is
at most a constant times
\[
 \frac A L\sum_{p_3\le X^\eta}\frac1{p_3\log p_3}
 \left(\sum_{p'<p_3}\frac{\log p'}{p'}\right)^2\ll A\eta.
\]
For exactly $j=1,2$ small primes write $n=avM$ with $v,M$ rough
past $X^\eta$, and expand the grouped divisor coefficient. Drop the
squarefree constraint on $vM$ at cost $o(A)$
since each repeated prime exceeds $X^\eta$. The square marks up to
$X^{.640}$ have total mass $O(AX^{-\eta}+AL^{-N})$ by Type~I for every
fixed $N$. The larger marks contribute $O(X^{.680}+X^{1/2})$ by
integer counting. The number of representations of any affected
integer is bounded by a divisor power, so
\eqref{eq:centered-odd-divisor-moments} and Cauchy--Schwarz absorb the
representations and the factor $L$ with error $o(A)$.

Apply the upper and lower fundamental-lemma sieves to $M$, with
the mark $av\le X^{r+2\eta}$ and level $X^h$, where $h=(\theta-r)/2$
and $2\eta<h$. A sieve divisor $\ell$ of $M$ gives the mark
$av\ell\mid n$. The main density is $V(z)/(av)$, where
$z=X^\eta$ and $V(z)=\prod_{3\le p\le z}(1-1/p)\ll(\eta L)^{-1}$.
The sieve envelopes have relative error $O(e^{-c_1/\eta})$.
Partial summation treats the dependence on $y$, and the errors in
the divisor marks sum to $O_N(AL^{-N})$ by Type~I, since the number
of factorizations of $av\ell$ is bounded by a divisor power.
The reciprocal rough $v$ sum
is $O(\eta^{-1})$ and the weighted $a$ sum is $O(\eta^j)$, giving
normalized error $O_J(e^{-c_1/\eta}\eta^{j-2})+o(1)$.
In the signed main term replace $y$ by one at a cost
$O_J((\prod t_i)(\log L)/L)$ in the coefficient bound, and use
$\sum b(n)=O(AL^{-B_2})$ with $B_2$ large. The main term is $o(A)$
for fixed $\eta$. Choose $\eta$ sufficiently small.
This leaves a fixed finite number of rough odd ranks.


\section{A uniform octal prime count}
\label{sec:octal-prime-count}

The following comparison is proved in Section~\ref{sec:octal-balanced-proof}.
It will suffice because the centered divisor sum has a simple form on
squarefree composites without a factor in the comparison interval.

\begin{lemma}\label{octal:balanced-shape-comparison}
Fix $d\in\{0,\ldots,7\}$, a positive roughness exponent $\kappa$,
and a prime factor count $r\ge2$. Write $b=\omega(f-f_0)$ as in
Section~\ref{sec:vaughan-signed}. On squarefree products of $r$
primes exceeding $X^\kappa$, consider any bounded smooth weight in
the normalized prime exponents $\log_n p_i$. Restrict the sum to
products admitting a proper divisor of exponent in $[.4,.6]$.
The resulting signed sum against $b$ is $o(A/L)$.
\end{lemma}

\begin{theorem}\label{octal:prime-count}
For every $d\in\{0,\ldots,7\}$ and all sufficiently large real $x$,
\[
 \frac9{100}\frac{x^{\log_8 7}}{\log x}
   \le \pi_{8,d}(x)
   \le 60\frac{x^{\log_8 7}}{\log x}.
\]
For $d\ne0$ the upper constant $60$ may be replaced by $56$.
\end{theorem}

\begin{proof}
Apply Proposition~\ref{octal:centered-odd-rank-switch} with its fixed
taper. The rough restriction permits us to retain only squarefree integers
with $P^-(n)>X^\kappa$ for a sufficiently small fixed $\kappa$, with
arbitrarily small error in units of $A$. At this roughness only finitely
many ranks occur. On each rank the coefficient $C_J(n)$ is a smooth
function of the labelled prime exponents. Lemma~\ref{octal:balanced-shape-comparison}
therefore gives $o(A)$ after multiplication by $L$ for all composites
having a divisor with exponent in $[.4,.6]$.

For a remaining composite, write the prime exponents as $t_i=\log_n p_i$.
If any $t_i>.6$, the large-prime cancellation of
Proposition~\ref{octal:centered-odd-rank-switch} applies. Otherwise each
$t_i<.4$. Call an exponent large if $t_i>.2$. Any pair of large exponents
has sum greater than $.4$, hence greater than $.6$ by the absence of a
divisor in the interval. Four large exponents would give two disjoint
pairs with total greater than $1$. If there are at most two large
exponents, start with one of them when present, omit the other if
present, and add the small exponents until the sum first reaches $.4$.
The available mass exceeds $.6$ and the last added exponent is at most
$.2$, so the resulting exponent lies in $[.4,.6]$. Thus there are
exactly three large exponents.

The test $J$ equals zero on subsets containing at least two of these
three primes, and one on their complements. The M\"obius signs cancel
in pairs if there is any further prime. On triples the coefficient
is $1-3=-2$. Thus the digital part of the triple contribution cannot
increase $L\sum C_J(n)b(n)$.

For the flat part, use $\omega\le1$. In the support of the cutoff,
$\log_X n=1+O(\log L/L)$. Each prime exponent relative to $n$ lies
in $[1/5,2/5]$. Thus if $x=\log_X p_1$ and $y=\log_X p_2$, the
three coordinates $x,y,1-x-y$ lie in this interval up to
$O(\log L/L)$. For fixed $p_1,p_2$, the number of choices for the last
prime is at most $\pi(X/(p_1p_2))$, which by the prime number theorem is
$(1+o(1))X/(p_1p_2(1-x-y)L)$ uniformly on this region. Reciprocal prime
summation in $p_1,p_2$ now gives, in units of $\gamma_d A$, a limiting
upper bound of
\[
 \frac13\int_{\substack{x,y,z\le 2/5\\z=1-x-y}}\frac{dx\,dy}{xyz}.
\]
The denominator is bounded away from zero. One can first enlarge
the coordinate bounds by any fixed $\epsilon>0$ and then let
$\epsilon\downarrow0$, since the affine boundaries have zero limiting
reciprocal prime measure. The factor $1/3$ is two divided by the six
orderings of the distinct primes.

Set $x_i=2/5-v_i$, so $v_i\ge0$, $\sum v_i=1/5$, and the area in
the free coordinates is $(1/5)^2/2$. The inequality
$\prod_i(1-u_i)\ge1-\sum_i u_i$ for $0\le u_i\le1$ gives
$xyz\ge(2/5)^2(1/5)$. The displayed bound is at most $5/24$.
The roughness error can be made arbitrarily small. For each fixed
choice of roughness the grouped comparison has error $o(A)$.
Since $0\le\omega\le1$ and $\log p\le L$, we obtain
\begin{equation}\label{octal:word-lower}
 \liminf_{k\to\infty}\frac{(k\log 8)P_d(8^k)}{7^k}
       \ge\frac{19}{24}\gamma_d.
\end{equation}
All counted word primes are admissible in ordinary notation, including
for $d=0$. For $8^k\le x<8^{k+1}$,
$7^k/(k\log8)\ge x^{\log_8 7}/(7\log x)$. Thus the liminf after
normalization by $x^{\log_8 7}/\log x$ is at least
$19\gamma_d/168\ge19/196>9/100$.

For the upper bound let $K=\lceil\log_8 x\rceil$. For $d\ne0$ we have
$\pi_{8,d}(x)\le P_d(8^K)$ by padding. Because $7^K\le7x^{\log_8 7}$ and
$K\log8\ge\log x$, \eqref{octal:word-upper} gives a normalized limsup
at most $(175/4)\gamma_d\le50$. For $d=0$, each positive integer has a
unique ordinary length, so $\pi_{8,0}(x)\le\sum_{j=1}^K P_0(8^j)$.
The geometric sum satisfies $\sum_{j\le K}7^j/j\sim(7/6)7^K/K$
by splitting at $K/2$ and applying dominated convergence to the upper
half after writing $j=K-i$. The same estimate sums the $o(1)$ errors
in \eqref{octal:word-upper}. The resulting normalized limsup is at
most $(175/4)(7/6)\gamma_0=175/3<60$. These strict margins give
the stated eventual bounds.
\end{proof}


\section{Fourier estimates for the octal prime comparison}
\label{sec:octal-fourier}

\begin{lemma}[Digital coefficients for dyadic characters]
\label{octal:semiprime-digital-char}
Write \(c=\log_8 7\), \(\sigma=.6175\), and \(\lambda=(c+\sigma-1)/2\). For
\(A_\chi=\sum_{n\bmod X}f(n)\overline{\chi(n)}\) one has
\begin{equation}\label{octal:semiprime-digital-moments}
 |A_\chi|\ll A q^{-\lambda}\quad(\chi\ {\rm primitive}\bmod q),
 \qquad
 \sum_{\chi\bmod q}|A_\chi|^2\ll A^2q^{1-c}.
\end{equation}
\end{lemma}

\begin{proof}
For the pointwise bound we may assume \(q\) exceeds an absolute
constant. Let \(Z\) be the least octal power divisible by \(q\),
let \(R\) be the least octal power at least \(\sqrt q\), and put
\(Q=q/R\) and \(V=Z/R\). Then \(R\mid q\), \(Q\mid R\),
\(Q\le V\le4Q\), and \(R,Q\asymp\sqrt q\). On \(1+R\mathbb Z\bmod q\)
the character is linear in the additive coordinate modulo \(Q\),
since \(R^2\equiv0\pmod q\). Primitivity gives
\(\chi(1+q/2)=-1\), so its frequency \(\ell\) is odd. With a
suitable sign for \(\ell\),
\[
 \overline{\chi(u+Rv)}=\overline{\chi(u)}e(\ell\bar u v/Q)
 \quad(u\ {\rm odd}).
\]
The lower block \(u\) of length \(R\) has its last digit in $K_d$.
The independent upper block \(v\) has length \(V\) and uses $D_d$
at every digit. Each residue of \(u\bmod Q\)
has \(O((R/Q)^c)\) allowed lifts, by fixing the digits below the
largest octal power dividing \(Q\). Inversion followed by
multiplication by \(\ell\) permutes the units modulo \(Q\).
The sum over \(v\) has absolute value \(V^cG_V(\ell\bar u/Q)\),
where \(G_V\) is the normalized absolute word transform.
Lemma~\ref{octal:uniformTypeI} and sampling give
\(\sum_{b\bmod Q}G_V(b/Q)\ll Q^{1-\sigma}\). One may convolve
with a reproducing difference of Fej\'er kernels to bound the sum
by \(O(Q+V)\int_0^1G_V(t)\,dt\).
Digits above \(Z\) contribute the factor \((X/Z)^c\), giving
\[
 |A_\chi|\ll (X/Z)^c V^c(R/Q)^c Q^{1-\sigma}
       \ll A q^{(1-\sigma-c)/2}.
\]
For the second estimate of \eqref{octal:semiprime-digital-moments},
each residue of \(n\bmod q\) has \(O((X/q)^c)\) allowed lifts.
The sum of all residue counts is \(O(A)\). By orthogonality the
character square sum is \(\varphi(q)\) times the sum of the squares
of the unit residue counts, which proves the result.
\end{proof}


For odd \(n\) and \(q\ge4\), the residue class averages also coincide
with conditional expectations on the odd unit group. Their
difference retains precisely the characters of conductor \(q\).
Consequently
\begin{equation}\label{octal:dyadic-char-expansion}
 \mathcal D_q f(n)=\frac2X
      \sum_{\substack{\chi\bmod q\\\chi\ {\rm primitive}}}
                     A_\chi\chi(n).
\end{equation}
Indeed the multiplicative coefficient at \(\chi\) in \(E_{X,q}f\)
is \(q A_\chi/(X\varphi(q))=2A_\chi/X\).


\subsection{A principal model for the long group}
\label{sec:principal-grouped-models}

We will use the interval Type~I estimate for $f$ and
$F_Q=E_{X,Q}f$ with $Q\mid X$ dyadic and $Q\ge X^\zeta$ for a fixed
$\zeta>0$. These are odd-supported, nonnegative, periodic sources of
mass at most $A$ on one full period of length $X$. We use the estimate
also on subintervals of $[1,C_0X]$ for fixed $C_0$. For $F_Q$ the proof
of Theorem~\ref{octal:strong-typeI} gives this version
directly. In the small denominator estimate there are a bounded
multiple of the original number of octal blocks. In the large
denominator estimate the geometric interval bound is
$O_{C_0}(X+S^2L)$ and the pointwise Ramanujan bound applies to all
$1\le n\le C_0X$. The $\ell^1$ approximation error increases by only
a bounded factor from the additional periods. This also gives the
stated interval estimate for $f=F_X$. By linearity it applies to
$g=f-F_Q$ and to $g=\mathcal D_Yf=F_Y-F_{Y/2}$ if both periods in
the latter expression exceed $X^\zeta$.

Let a fixed number of prime factors of scales $P_i\gg X^\kappa$ be
partitioned into two nonempty groups $I_1,I_2$, where $\kappa>0$ is
fixed. Put $U_h=\prod_{i\in I_h}P_i$ and suppose that
$X L^{-C}\ll U_1U_2\ll X$ and $U_1\le U_2\ll X^{.601}$. Let each $w_i$ be
a smooth function supported in a fixed compact subset of $(0,\infty)$,
with every fixed derivative bounded by a power of $L$. The actual
coefficient at the grouped integer $v$ is
\[
 b_h(v)=\sum_{\substack{\prod_{i\in I_h}p_i=v\\ p_i\ \mathrm{prime}}}
                 \prod_{i\in I_h}(\log p_i)w_i(p_i/P_i).
\]
To count unweighted primes, include $1/\log(P_i t)$ in $w_i(t)$.
Choose one slot $i_h\in I_h$ and define
\[
 \Phi_h(v)=\int_{(0,\infty)^{I_h\setminus\{i_h\}}}
   w_{i_h}\left(\frac{v}{P_{i_h}\prod_{i\ne i_h}t_i}\right)
   \prod_{i\ne i_h}w_i(t_i/P_i)\,\frac{dt_i}{t_i},
\]
where the products in this formula are within group $h$. The empty
integral is the single-factor test. Write $c_d(v)$ for the Ramanujan
sum $\sum_{a\bmod d}^{*}e(av/d)$, with $c_1(v)=1$. For $D=L^{K_0}$,
where $K_0>0$ is fixed, put
\begin{equation}\label{eq:principal-grouped-coefficients}
 B_2(v)=2\,1_{v\ \mathrm{odd}}\,\Phi_2(v)
       \sum_{\substack{d\le D\\ d\ \mathrm{odd}}}
                   \frac{\mu(d)}{\varphi(d)}c_d(v),
 \qquad b_2^*(v)=b_2(v)-B_2(v).
\end{equation}
All these coefficients are supported in fixed positive expansions
of the original group boxes.

\begin{proposition}\label{octal:principal-grouped-model}
Let $g$ be either of the signed sources above, and let $\Omega$ be a
smooth product test with fixed relative derivatives bounded by powers
of $L$. The constant test $\Omega=1$ is permitted. Then
\begin{align}
 \sum_{x,y} b_1(x)b_2(y)\Omega(xy/X)g(xy)
 &=\sum_{x,y} b_1(x)b_2^*(y)\Omega(xy/X)g(xy)\notag\\
 &\quad+2\sum_{n>0}g(n)\Omega(n/X)
                 \int_0^\infty \Phi_1(z)\Phi_2(n/z)\frac{dz}{z}
       +O_N(AL^{-N})                                           \label{eq:principal-grouped-representation}
\end{align}
for every fixed $N>0$. The squared counting norms of $b_1,b_2^*$ are
$O(U_1L^C),O(U_2L^C)$, respectively, and both coefficients are
$O(\tau(v)L^C)$. The exponent $C$ can be chosen independently of $K_0$.

For each fixed $K_1,N>0$, the coefficients satisfy the following
small-arc estimate uniformly in $(a,b)=1$, $1\le b\le D$, and
$U_2|\beta|\le L^{K_1}$:
\begin{equation}\label{eq:principal-grouped-low-arc}
 \left|\sum_v b_2^*(v)e\bigl((a/b+\beta)v\bigr)\right|
                    \ll_{K_0,K_1,N} U_2L^{-N}.
\end{equation}
The estimates include retained Mellin twists of polylogarithmic
height in the individual tests.
\end{proposition}

\begin{proof}
The bounds $|b_h(v)|\ll L^C$ follow from unique factorization
with boundedly many slots. The function $\Phi_h$ is supported at
scale $U_h$ and has derivative of order $j$ bounded by
$O_j(U_h^{-j}L^{C_j})$. For $B_2$ expand
$c_d(v)=\sum_{b\mid(d,v)}b\mu(d/b)$. Only squarefree $d$ occur,
so $\varphi(d)=\varphi(b)\varphi(d/b)$ for $b\mid d$. Hence
\[
 \left|\sum_{\substack{d\le D\\ d\ \mathrm{odd}}}
            \frac{\mu(d)c_d(v)}{\varphi(d)}\right|
 \le \sum_{b\mid v}\frac{b}{\varphi(b)}\sum_{e\le D}\frac1{\varphi(e)}
 \ll \tau(v)(1+\log v)(1+\log D).
\]
The last inequality follows from
$m/\varphi(m)=\sum_{s\mid m}\mu(s)^2/\varphi(s)$ and
$\sum_s\mu(s)^2/(s\varphi(s))<\infty$. Since
$1+\log D\ll_{K_0}L$, this proves the pointwise bound with a log
exponent independent of $K_0$. Also $\tau(v)^2\le\tau_4(v)$, by
checking prime powers, and $\sum_{v\le V}\tau_4(v)\ll V(1+\log V)^3$.
This gives the asserted squared norm.

In the pairing of $b_1$ and $B_2$, hold $x$ fixed and expand the
Ramanujan sum at modulus $d$. The terms indexed by $b\mid d$ are
\[
 2b_1(x)\frac{\mu(d)b\mu(d/b)}{\varphi(d)}
     \sum_{n>0} g(n)\Omega(n/X)\Phi_2(n/x)1_{xb\mid n}.
\]
For each fixed $d,b$, replace the mark $1_{xb\mid n}$ by
$1/(xb)$ using interval Type~I and partial summation. The mark
is odd and $xb\ll U_1D$ has fixed reserve below the Type~I level.
The total variation and supremum of the real-variable test in $n$
are bounded by powers of $L$, uniformly in $x$. Indeed the relevant
product scale is between $X$ and $XL^{-C}$ up to constants, and
$x\asymp U_1$. The map $x\mapsto xb$ is injective for fixed $b$,
and $|b_1(x)|\ll\tau(xb)L^C$. Summing over the polylogarithmically
many choices of $d,b$ costs only powers of $L$, so the total
replacement error is $O_N(AL^{-N})$ for every $N$.
For $d>1$ the resulting main term is zero because
$\sum_{b\mid d}\mu(d/b)=0$.

For $d=1$ the pairing in which the other coefficient is genuine
becomes
\[
 2\sum_n g(n)\Omega(n/X)\sum_x \frac{b_1(x)}{x}\Phi_2(n/x).
\]
The ordinary prime number theorem and partial summation in the
prime slots give, uniformly in $n$, up to $O_N(L^{-N})$, the
continuous reciprocal sum
$\int_0^\infty\Phi_1(z)\Phi_2(n/z)\,dz/z$. To see the formula,
replace $(\log p_i)/p_i$ by $dt_i/t_i$ in each slot and change
variables from $t_{i_1}$ to the group product $z$, fixing the
other slots. The absolute mass of $g$ is $O(A)$ on the product
support. Inserting the continuous reciprocal sum into the pairing
with $B_2$ proves \eqref{eq:principal-grouped-representation}.

To prove \eqref{eq:principal-grouped-low-arc}, fix the companion
primes in group $2$, let their product be $c$, and abbreviate
$i_2=i$ and $P_i=P$. They are coprime to $b$ for large $X$.
The selected prime variable ranges through reduced classes modulo
$b$, since its scale is a fixed positive power of $X$. The
Siegel--Walfisz theorem and partial summation give
\begin{align*}
 \sum_p(\log p)w_i(p/P)e(acp/b+\beta cp)
 &=\frac1{\varphi(b)}\sum_{u\bmod b}^{*} e(acu/b)
               \int_0^\infty w_i(t/P)e(\beta ct)\,dt
       +O_N(PL^{-N})\\
 &=\frac{\mu(b)}{\varphi(b)}\int_0^\infty
                       w_i(t/P)e(\beta ct)\,dt+O_N(PL^{-N}).
\end{align*}
Here the modulus is polylogarithmic and the variation of the
selected-prime test is a power of $L$, since $|\beta|cP\ll
L^{K_1}$. After summing the companions the total error is
$O_N(U_2L^{-N})$. In the main term put $v=ct$ and use the
ordinary prime number theorem and partial summation in the
companion slots. The result for the genuine coefficient is
\begin{equation}\label{eq:principal-grouped-genuine-arc}
 \sum_v b_2(v)e((a/b+\beta)v)
     =\frac{\mu(b)}{\varphi(b)}\int_0^\infty
                         \Phi_2(v)e(\beta v)\,dv+O_N(U_2L^{-N}).
\end{equation}

For the model, use $2\,1_{v\ \mathrm{odd}}=1-e(v/2)$ and the
defining exponential sum for $c_d$. A rational phase can pair
with $e(av/b)$ to give an integer frequency only when its
reduced denominator matches. For odd $b$ this occurs in the
$1$ part exactly at $d=b$, with coefficient $\mu(b)/\varphi(b)$.
For $b=2d$ with $d$ odd it occurs only in the $-e(v/2)$ part,
with coefficient $-\mu(d)/\varphi(d)=\mu(b)/\varphi(b)$.
For $4\mid b$ there is no such phase and $\mu(b)=0$. In each
resonant case the required odd denominator is at most $D$.
Every nonresonant rational phase has distance $\gg D^{-2}$
from an integer, which is unchanged up to constants on adding
$\beta$. Apply Poisson summation to the smooth $\Phi_2$. Its
Fourier transform at frequency $\xi$ is
$O_j(U_2L^{C_j}(1+U_2|\xi|)^{-j})$ for every fixed $j$.
Since $U_2$ is a fixed positive power and all model denominators
are polylogarithmic, the nonresonant contributions and all
noncentral Poisson images of the resonant term are $O_N(U_2L^{-N})$.
This proves \eqref{eq:principal-grouped-genuine-arc} also for
$B_2(v)$ in place of $b_2(v)$. Their difference is
\eqref{eq:principal-grouped-low-arc}.
\end{proof}

\paragraph{Dependence on the shape and the product cutoff.}
For a smooth shape depending on the normalized prime exponents,
choose a smooth partition over the eligible labelled groupings and
insert dyadic partitions in the prime slots. On each box one can
separate the resulting shape and the product cutoff by Fourier
inversion in the scaled logarithms of the prime variables. Smooth
compact extensions of these tests have fixed-order derivatives
bounded by powers of $L$. Their Fourier first norms are
polylogarithmic, and the integral can be truncated at a sufficiently
large fixed power of $L$ with arbitrary inverse-logarithmic error.
The error in the original pairing is bounded using the bounded
fixed-rank representation multiplicity and the $O(A)$ absolute
source mass. At each retained twist put the phases into the genuine
slot tests before defining \eqref{eq:principal-grouped-coefficients}.
Then \eqref{eq:principal-grouped-representation} applies with
$\Omega=1$. Its continuous term recombines by linearity and Fourier
inversion, with the same truncation estimate for the continuous
integral. This term is independent of the choice of group split
and of the prime slots selected to define $\Phi_h$.

More explicitly, for a fixed shape weight $W$ in $r$ labelled
prime slots and the product weight $\omega(n/X)$, the recombined
continuous factor at a fixed product $n$ is
\[
 \frac{\omega(n/X)}{\log n}
    \int_{u_1+\cdots+u_r=1} W(u_1,\ldots,u_r)
                        \prod_{i=1}^r\frac1{u_i}\,du_1\cdots du_{r-1}.
\]
Indeed use unweighted prime densities $dt_i/\log t_i$, fix the
first $r-1$ primes, and set $t_i=n^{u_i}$. Here the roughness
restriction is fixed, so the simplex integral is finite. A fixed
logarithmic box boundary can be tapered before separating the
weights. The strip errors are bounded by interval Type~I and the
marked upper sieve in the original sharp comparison. Also the
change from a fixed $X$-relative roughness cutoff to its
$n$-relative version is harmless, since $\log_X n=1+O(\log L/L)$
on the product support.

For the whole source $g=f-E_{X,Q_*}f$ with
$Q_*=2^{\lfloor .741\log_2 X\rfloor}$, the original and projected
word have equal full mass. On discarded endpoint intervals the
original word contributes $o(A)$ by the digit interval bound.
The projected word has mass $O(AQ_*/X)$ per period, so an interval
of length $o(X)$ contributes $o(A)$ for that source as well.
Replacing $1/\log n$ by $1/L$ on the product support incurs
$o(A/L)$ there. Thus the entire continuous term for this shape
is $o(A/L)$ after restoring the product weight. The residual
pairing may be split into the exact conductors $Y>Q_*$, and the
small-arc estimate for their actual residual coefficients is
\eqref{eq:principal-grouped-low-arc}.


\subsection{Digital weights in the lower-period affine equation}
\label{sec:ordinary-group-gram-twists}

Let $q$ be the exact dyadic period of a lower word, $\rho_L$ its
density scale, and $Q\mid q$ an octal power. Write $a=.3825$.
The off-diagonal lower Gram, after completion in the long variable,
involves sums over $i,b\bmod q$ with $pi-p'b\equiv t\pmod q$.
Its Fourier indices and rows are odd for $q>2$.
Split each normalized digital amplitude at the upper octal block of
length $Q$, taking the least octal multiple of $q$ for the word length.
On residue classes $i_0,b_0\bmod Q$ the upper factors $U(i_0),V(b_0)$
are fixed. Each lift of $i_0$ determines a lift of $b_0$ by an odd
affine permutation. Shifted Cauchy--Schwarz in the lower blocks
therefore costs $O((q/Q)^\delta)$, with bounded aliases when $q$ is
non-octal. The remaining upper weights obey
\[
 \sum U,\ \sum V\ll Q^a,\qquad \sum U^2,\ \sum V^2\ll Q^\delta
\]
by the shifted moments. A bounded mask on the original amplitudes
can be ignored in these upper bounds.

Write $s=(i_0-b_0)\overline{b_0}\bmod Q$ and
$p'=p+j$. The affine equation is $ps-j\equiv t\overline{b_0}\pmod Q$.
The total Fourier weight on a slope $s$ is at most
$B(s)=\sum_{b_0\bmod Q}^{*}U((1+s)b_0)V(b_0)$, set to zero when
$1+s$ is even. For every octal divisor $Q_1\mid Q$, split the upper
blocks again to obtain
\begin{equation}\label{eq:ordinary-group-gram-overlap}
 B(s)\ll (Q/Q_1)^\delta B_{Q_1}(s\bmod Q_1),\qquad
       \sum_{\sigma\bmod Q_1}B_{Q_1}(\sigma)\ll Q_1^{2a}.
\end{equation}
Here $B_{Q_1}$ is the overlap of the two top subwords of length $Q_1$.
For each $b_0$ class modulo $Q_1$, multiplication by $1+s$ permutes
the lifts. Shifted Cauchy--Schwarz costs $(Q/Q_1)^\delta$ on the lower
blocks and the two first moments bound the sum of the top overlaps.


\paragraph{Sampling on separated cells.}\label{sec:zero-model-wide-source-bands}
For the amplitude $F_N$ of a normalized trigonometric polynomial of degree less than $N$,
let $I_j$ have length comparable to $N^{-1}$ and bounded overlap.
For $s\ge1$, the integral and variation on each cell give
\[
 \sum_j\sup_{t\in I_j}|F_N(t)|^s
 \ll N\int_0^1 |F_N|^s+\int_0^1\big|(|F_N|^s)'\big|
 \ll N\int_0^1 |F_N|^s .
\]
The last step is Bernstein's inequality for the underlying
trigonometric polynomial and H\"older. Thus a shifted counting
moment bound $O(N^{g_s})$ also bounds these sums of suprema by
$O(N^{g_s})$, on integrating the shifted moment. The argument covers
mixed words and cells with a bounded number of dyadic aliases.

\subsection{Upper digit factors and horizontal slopes}
\label{sec:short-row-weighted-offsets}

Consider an integer bilinear sum at group scales $P,M$, with $PM\asymp X$
up to powers of $L$, with $PM\ll X$. After separating smooth product weights, write its
source in the form $F(xy)U(\lfloor xy/R\rfloor)$, where $F$ is periodic
and the upper word $U$ satisfies
\[
 \sum_{|v|\le C X/R}|U(v)|^2\ll (X/R)\rho_U L^{O(1)}
\]
for every fixed $C$. The support intervals of the two coefficient
sequences are at fixed positive multiples of $P$ and $M$. In the
application, $U$ is the exact difference of two upper digital words
or its periodic lift, and $\rho_U$ is the original upper density.

\begin{lemma}\label{octal:short-row-upper-density}
Partition the row variable into intervals $I$ of integer length $H$,
where $H\le P$ and $HM\ll R$. If $\sum_y|v_y|^2\ll M L^{O(1)}$, then
\begin{align*}
 &\left|\sum_{x,y}u_xv_y F(xy)U(\lfloor xy/R\rfloor)\right|^2\\
 &\qquad\ll \frac{PM\rho_U}{H}L^{O(1)}
    \sum_{I,y}\max_{J\subseteq I}
                       \left|\sum_{x\in J}u_x F(xy)\right|^2,
\end{align*}
where the maximum is over subintervals. The same estimate holds if
the shorter scale is assigned to either coordinate.
\end{lemma}

\begin{proof}
For each $I,y$, the index of $U$ changes at most a fixed number of
times as $x$ traverses $I$. The modulus of the corresponding row
sum is therefore bounded by the indicated maximum times a fixed
multiple of the largest $|U|$ at one of those indices. Denote this
largest value by $A_I(y)$. For fixed $y$, the indices at the left
endpoints of $I$ increase by $Hy/R$ before taking the integer part.
Each integer index, including any of its bounded neighboring indices,
is used at most $O(1+R/(Hy))=O(R/(HM))$ times. They all lie within
$O(X/R)$ of the origin. Thus
\[
 \sum_I A_I(y)^2\ll \frac{P}{H}\rho_U L^{O(1)}.
\]
Apply Cauchy to the sum over $(I,y)$, using the factors
$|v_y|A_I(y)$ and the indicated maximum. This proves the inequality.
The argument counts indices over the row intervals for a fixed column,
so it does not require the individual row centres to be below $R$.
\end{proof}

To estimate the maximum on each $I$, decompose any $J\subseteq I$ into
at most two binary intervals at each length. Cauchy bounds the square
by a logarithmic factor times the sum of squares over all binary
intervals in $I$. At each length they are disjoint, and every row
pair they contain has difference of magnitude at most $H$. Thus the
usual ordinary completion reduces to the lower row Gram of
Section~\ref{sec:ordinary-group-gram-twists}, with this shorter
range of row differences. In particular its squared normalization
loses only one factor $\rho_U^{-1}$.

\begin{lemma}[Horizontal frequencies with their exact multiplicities]
\label{octal:horizontal-weighted-offset}
In a lower row Gram with a coarse octal modulus $Q$ comparable to
the lower period, let $g$ be the dyadic additive order of a horizontal
slope and $q_g=Q/g$. Take $g\ll P$, and keep row differences
$0<|j|\ll H$ and ordinary completion indices $|t|\ll T$. Let $B$
be a dyadic shell of reduced rational denominators associated to
$j b_0+t=0\pmod{q_g}$. If $B^2\max(1,T/B)\le c_*q_g$ for a sufficiently small absolute
constant $c_*$, then this shell costs at most
\begin{equation}\label{eq:horizontal-weighted-offset}
 g^{2a-1}B^{2\delta-1}\max(1,T/H)^\delta L^{O(1)}
\end{equation}
in the absolute lower Gram, after dividing by $PH$ and the square
of the lower density. Here $a=.3825$ and $\delta=1-\log_8 7$. Bounded common masks on the lower
additive coefficients are allowed. In the complementary range the
bound is
\begin{equation}\label{eq:horizontal-weighted-overflow}
 g^{2a-1}q_g^\delta(1/B+T/q_g)L^{O(1)}.
\end{equation}
\end{lemma}

\begin{proof}
We describe the horizontal reduction to record its counting factors.
Let $b$ be the second lower Fourier index, odd modulo $Q$, and let
$s$ be the slope, so that the first index is $(1+s)b$. In the affine
row equation $ps-j=t\bar b\pmod Q$ write $q_g=Q/g$, with $s/q_g$
a unit modulo $g$. Fixing $s,b,j,t$, there are $O(P/g)$ possible
rows since the equation fixes the row modulo $g$ whenever it is
soluble. Necessarily $jb_0+t=0\pmod{q_g}$ for $b_0=b\bmod q_g$.
For each $b_0$, sum the low subword amplitudes over both sets of
lifts modulo $Q$. For fixed $b$ the frequencies $(1+s)b$ run through
a subset of the lifts of $b_0$ once. The two shifted first norms
cost $g^{2a}L^{O(1)}$. The remaining normalized amplitudes, say
$V_1(b_0)V_2(b_0)$, are those of the upper subwords at modulus
$q_g$. Bounded aliases from the lower period divided by $Q$
are harmless by the shifted second moment. Dividing by $PH$
gives the bound
\[
 \frac{g^{2a-1}}{H}L^{O(1)}
 \sum_{b_0\bmod q_g}^{*} V_1(b_0)V_2(b_0)
       \sum_{\substack{0<|j|\ll H,\ |t|\ll T}}
                    1_{j b_0+t=0\ (q_g)}.
\]

Fix one reduced rational $a_0/B_0$, $B_0\asymp B$, obtained by
reducing $(j b_0+t)/(j q_g)$ as a real rational. Then $B_0\mid j$.
Writing $D_0=|b_0-q_g a_0/B_0|$, the exact identity
$|j|D_0=|t|$ shows that the number of possible $j$, divided by $H$,
is at most
\begin{equation}\label{eq:horizontal-count-per-offset}
 O\left(\frac1B \min\left(1,\frac{T}{H D_0}\right)\right).
\end{equation}
The statement includes $D_0=0$, by taking the minimum to be one.
For each such $j$ the value of $t$ is fixed by the real rational.
A nonempty band has $B\ll H$ and $D_0\ll T/B$.

If $D_0\le V$, the total $V_1V_2$ mass in the bands over the
denominator shell is at most
\[
                  O\bigl((B^2\max(1,V))^\delta L^{O(1)}\bigr).
\]
To prove this, split the word at a bottom block of length
comparable to $B^2$ and a terminal block of length comparable to
$\max(1,V)$. The rational centres are separated by $\gg B^{-2}$
on the circle, with only bounded repetitions from the real lifts
in this range. The sum of local suprema of the squared bottom
amplitude is bounded by its shifted square norm at this density,
using Bernstein's inequality as in
Section~\ref{sec:zero-model-wide-source-bands}. On the offset grid
the terminal block has square sum $O(\max(1,V)^\delta)$ by shifted
Parseval. Cauchy gives the claim for the two possibly different
amplitudes. Dyadic $q_g$ differing from an octal power by a factor
at most four gives only bounded aliases.

Use \eqref{eq:horizontal-count-per-offset} on $D_0\le1$ and each
dyadic shell $D_0\asymp V>1$. Since $0<\delta<1$, the weighted
sum of $V^\delta$ against $\min(1,T/(HV))$ over $V>1$ is
$O(\max(1,T/H)^\delta)$. This proves
\eqref{eq:horizontal-weighted-offset}. In the complementary range
one can count for each $b_0$ at most $O(1+BT/q_g)$ reduced real
rational centres within distance $O(T/(Bq_g))$ of $b_0/q_g$, by
their real separation. Each uses $O(H/B)$ row differences. The
complete square norms on $q_g$ then give
\eqref{eq:horizontal-weighted-overflow}. This counting also includes
the real lifts when an offset crosses the boundary of a period.
\end{proof}


\paragraph{A bilinear bound at a single rational approximation.}
\label{sec:dynamic-rational-source-bands}
For coefficients $u_x,v_y$ on intervals at scales $P\le M$, with squared
counting norms at most $PL^C$ and $ML^C$, Cauchy--Schwarz in the long
variable gives
\[
 \frac{\left|\sum_{x,y}u_xv_y e(\alpha xy)\right|^2}{(PM)^2}
 \ll \frac{L^{O(1)}}{PM}\sum_{|j|\ll P}\min(M,\|\alpha j\|^{-1}).
\]
The inner geometric sum is bounded by the minimum displayed, and the
coefficient autocorrelation at every $j$ is at most $PL^C$ by Cauchy.
For $\alpha=a/b+\beta$ and $S=X|\beta|$, suppose $b(1+S)\ll P$ with a
small fixed constant, allowing a logarithmic reserve when $PM\asymp X$
up to powers of $L$. For $j$ not divisible by $b$ the displacement stays
below a fixed fraction of $1/b$, and the geometric sums on these classes
total $O(P L^{O(1)})$. For $j=bk$ the arguments are $\beta bk$ without
wraparound. Summing their reciprocals with the cutoff at $M$ gives
$O((M+PM/[b(1+S)])L^{O(1)})$, by the trivial class bound when $S\le1$.
Thus the bilinear sum divided by $PM$ is
$O((b(1+S))^{-1/2}L^{O(1)})$ in this range.


\subsection{The full high-conductor comparison for grouped products}
\label{sec:near-balanced-high-conductors}

\begin{proposition}\label{octal:near-balanced-high-conductors}
Use the genuine short coefficient $u_x=b_1(x)$ and the long residual $v_y=b_2^*(y)$
of Proposition~\ref{octal:principal-grouped-model}, with any retained
logarithmic Mellin phases placed in the prime-slot weights. Let the shorter
group have reference scale $P=X^\nu$, with $.399\le\nu\le.500$, and let
$XL^{-C}\ll PM\ll X$. Take an exact dyadic conductor $Y=X^\eta\mid X$ with
$.740\le\eta\le1$. For every $J>0$, one may take the fixed power $K_0$ in
the principal model large enough that
\[
             \left|\sum_{x,y}u_xv_y\mathcal D_Y f(xy)\right|
                          \ll_J AL^{-J}.
\]
The choice is uniform in $Y$ and across all bounded sets of prime-slot
derivative bounds of the stated logarithmic size.
\end{proposition}

\begin{proof}
Use $a=.3825$, $h_0=.067$, and $\delta=1-\log_8 7<\delta_0=.0643$,
the last inequality following from $7^{10000}>8^{9357}$.
When the actual smaller scale exceeds $\sqrt X$ by a bounded factor,
take the reference scale to be $\sqrt X$ and absorb that factor in its
support constants. The coefficients are odd supported, have squared
counting norms $\ll PL^C,ML^C$, and are pointwise $\ll X^\epsilon$ for
every fixed $\epsilon>0$. The power of $L$ in the squared norms is
independent of $K_0$. The only further arithmetic information used in
this proof is that, on any specified polylogarithmic range for $b$ and
$S$, and for any fixed $N$, we may arrange
\begin{equation}\label{eq:near-balanced-small-model-arc}
    \left|\sum_y v_y e(\alpha xy)\right|\ll_N M L^{-N}
   \quad\bigl(\alpha=a_0/b+O(S/X),\ x\asymp P\bigr).
\end{equation}
Indeed multiplication by $x$ gives an angle with reduced denominator
at most $b$ and offset $O(S/M)$, up to powers of $L$, so this follows
from \eqref{eq:principal-grouped-low-arc}. The choice of the
polylogarithmic range is made after the logarithmic losses in the
elementary estimates below have been fixed.

Put $\mu=1-\nu$ and $\lambda=\max(.589,\mu+h_0)$, and choose
an octal cut $R\asymp X^\lambda$. The exact source
$(E_{X,Y}-E_{X,Y/2})f$ has physical factorization
\begin{equation}\label{eq:near-balanced-source-factorization}
   L_R(n\bmod R) U(\lfloor n/R\rfloor\bmod Y/R),
\end{equation}
times $(X/Z)^{-\delta}$, where $Z$ is the least octal multiple of $Y$.
Indeed $R\mid Y/2$, so both expectations act on the upper word, and
averaging the digits above $Z$ contributes their density. After
factoring out this constant, the upper normalized transform has
the shifted local moment bounds of the word at length $Z/R$ on
the $Y/R$ frequency grid, with bounded octal aliases. The
upper signed word has $O((Y/R)\rho_U)$ jumps in its step
extension, where $\rho_U=(Z/R)^{-\delta}$ is its density before
taking the exact projection. This follows by taking the union of
the supports of the boundedly many translates needed for both
conditional expectations below the least octal multiple of $Y$.
The squared values on one period sum to $O((Y/R)\rho_U)$.
Decompose $L_R$ into exact dyadic periods $q\mid R$, combining
the periods at most two. Denote a lower component by $L_q$ and
its density scale by $\rho_L=(|K_d|/7)R^{-\delta}$, the density of
$L_R$ before projection. In units of this density the lower amplitude has the shifted word
moment bounds at length comparable to $q$. The digits above that
length contribute their density. In the least significant position,
the normalized $K_d$ factor has modulus at most one. Dropping it
in the first moment costs at most eight copies of a shortened
shifted grid. For the square moment, shifted Parseval costs the
factor $7/|K_d|$. These bounds are uniform since $|K_d|\ge3$.
In particular its squared mean is $O(\rho_L^2 q^\delta)$.
The comparison scale $X\rho_L\rho_U$ times the omitted density
$(X/Z)^{-\delta}$ is comparable to $A$.


\paragraph{Lower row completion.}
For a lower period $q$ put $\gamma=\log_X q$ and take
$H\asymp\max(X^{.065},q/M)$ up to powers of $L$, decreasing
it by a logarithmic factor if needed so that $HM\ll R$.
This is possible since $\lambda\ge\mu+.065$ and $q\le R$.
Lemma~\ref{octal:short-row-upper-density}
and its binary interval decomposition reduce the squared
comparison to the completed lower Gram for pairs of rows at
distance $O(H)$, losing just $\rho_U^{-1}$ in the relative
squared estimate. The diagonal is at most
$O(X^\epsilon \rho_U^{-1} q^\delta/H)$, by summing the squared
lower values by the integer product with the divisor bound.
The smooth column majorant leaves only
$|t|\ll T=(1+q/M)X^{\epsilon_0}$, at arbitrary power accuracy.
The exponents $\epsilon,\epsilon_0$ will be taken sufficiently
small with respect to all fixed margins below.

We now bound the contribution of nonhorizontal slopes to this Gram.
Choose an octal divisor $Q$ of $q$ with $Q\asymp\min(q,PH)$.
Use the least one exceeding a sufficiently large fixed multiple
of $PH$ if available, and otherwise the largest octal divisor
of $q$.
Slopes $s$ of additive order $g=Q/(s,Q)\ll P$ are called
horizontal, using a fixed comparison constant large enough
to include every slope with a horizontal difference of length
at most the diameter of the row box. For the other slopes the
absolute contribution at each fixed twist $t$, in units of
$\rho_L^2$, is at most
\begin{equation}\label{eq:near-balanced-capped-affine}
              (q/Q)^\delta Q^{2a}(1+PH/Q)L^C.
\end{equation}
To see this, rescale the difference lattice of the affine
equation by $P,H$. Its determinant is $Q/(PH)$. If $\ell$ is
the length of a shortest nonzero vector, the affine count in
the rescaled box is $O(1+PH/Q+1/\ell)$. Indeed parallel
lattice lines have spacing $Q/(PH\ell)$ and the points on
each have spacing $\ell$. If two points occur, $\ell\ll1$.
Use the full two digit first norms at multiplicities
$O(1+PH/Q)$. For multiplicity $J\gg1+PH/Q$, a shortest
vector gives $0<|v|\ll P/J$ and $0<|w|\ll H/J$. A purely
vertical vector has length at least $Q\gg H$, and horizontal
vectors were excluded. Put $c=(v,Q)$, which divides $w$, and
divide out the full greatest common divisor of $(v,w)$. Orient
the resulting primitive pair with first coordinate $v_0>0$.
Its first coordinate is odd and $s\equiv w_0/v_0\pmod{Q/c}$.
Choose an octal $Q_1\mid Q/c$ comparable to $PH/(Jc)^2$ with a
sufficiently large comparison constant. It exists on nonempty
vector ranges since $J\gg 1+PH/Q$. Equality of two ratios modulo
$Q_1$ would give a vanishing integer cross determinant, whose
absolute value is less than $Q_1$, hence the same primitive pair.
The ratio-overlap bound \eqref{eq:ordinary-group-gram-overlap} gives
$\ll c (Q/Q_1)^\delta Q_1^{2a}L^C$. Summing in the dyadic $c$,
with multiplicity $O(J)$, gives
$\ll Q^\delta(PH)^{2a-\delta}J^{1-2(2a-\delta)}L^C$.
Since $2(2a-\delta)>1$, the tail sums to
$Q^{2a}(1+PH/Q)L^C$ on the coarse grid. In the original
congruence each lift of one index determines the other lift
by an odd affine permutation. Shifted Cauchy therefore gives
the factor $(q/Q)^\delta$ in \eqref{eq:near-balanced-capped-affine}.
Bounded masks on the
lower coefficients only decrease these positive majorants.

For $.119\le\gamma\le\mu-.0001$ take $H\asymp X^{.065}$ up to powers of $L$
and retain only $t=0$, by smooth completion. In this untwisted
case fix the odd greatest common divisor $c$ of a distinct row
pair and write $B_c=PH/c^2$. Here $c\ll H$ since it
divides the nonzero row difference. The reduced positive pair is
primitive. At a coarse modulus $Q_c\mid q$ chosen octal with
$Q_c\asymp\min(B_c,q)$, each ratio is attained by
$O(1+B_c/Q_c)$ primitive pairs. Rescale the primitive box
by $P/c,H/c$. If the shortest lattice vector is shorter than
a suitably small multiple of the determinant $Q_c/B_c$,
all points in the bounded box are collinear and at most one
is positive primitive. Otherwise counting points on parallel
lattice lines proves the asserted multiplicity. The digit
first norms and the lifted second norms give at fixed $c$
the bound $(q/Q_c)^\delta Q_c^{2a}(1+B_c/Q_c)L^C$.
Put $F=PH$ and $u=2a-\delta>1/2$. The bound at fixed $c$ is
$O(q^\delta B_c\min(q,B_c)^{u-1}L^C)$.
Since $\min(q,B_c)\ge c^{-2}\min(q,F)$ and $u<1$, summing $c^{-2u}$
gives $O(q^\delta F\min(q,F)^{u-1}L^C)$.
After dividing by $PH$ and including the upper density cost, the
relative exponent is bounded by $(2a-1)(.119)+.0643(1-\lambda)<0$
when $q\ll PH$, and by $.0643-.2993(.399+.065)<0$ when $q\gg PH$.
These bounds also apply at the constant-factor transition.
The diagonal exponent is at most $.0643-.065<0$.

For $\gamma\ge\mu-.0002$ put $h=\max(.065,\gamma-\mu)$
and retain the $T$ ordinary completions. If $q\gg PH$, the
relative nonhorizontal squared exponent from
\eqref{eq:near-balanced-capped-affine} is at most
\begin{equation}\label{eq:near-balanced-nonhorizontal-exponent}
 \delta(\gamma+\eta-\lambda)
        +(2a-\delta-1)(\nu+h)+(\gamma-\mu)_++O(\epsilon_0+\epsilon).
\end{equation}
We may use $\delta_0=.0643$ in this expression because the
net coefficient of $\delta$ is nonnegative when $q\ge PH$.
The expression is increasing in $\gamma$, so its largest
value is bounded by its value at $\lambda$, namely
\[
 \delta_0\eta-(1+\delta_0-2a)\nu+(2a-\delta_0)(\lambda-\mu).
\]
When $\lambda=\mu+h_0$ this is at most $-.000392$ uniformly
by evaluation at the left endpoint in $\nu$ and $\eta=1$. When $\lambda=.589$ its
maximum is at $\nu=.5$ and is smaller. If $q\ll PH$, the
corresponding exponent is
$(2a-1)\gamma+(\gamma-\mu)_++\delta(\eta-\lambda)$.
This is negative below $\mu$ and is bounded above $\mu$
by $2a\lambda-\mu+\delta(\eta-\lambda)<0$.
The transition with $q\asymp PH$ is covered by these
same bounds. The diagonal is paid since $H\gg X^{.065}L^{-C}$.

For a coarse horizontal slope, take the additive order $g$
of its true slope modulo $q$. It satisfies $g\ll(q/Q)P$.
Put $q_g=q/g$ and keep the exact affine equation rather than
projecting this contribution to $Q$. For $g\ll P$,
Lemma~\ref{octal:horizontal-weighted-offset} applies on the
full lower grid with at most bounded dyadic aliases.

For $g$ larger than a fixed multiple of $P$, write the two
indices as $b_0+k' q_g$ and $b_0+k q_g$. The exact order $g$
makes $k'-k$ odd. The equation for the rows $p,p+j$ becomes
\begin{equation}\label{eq:near-balanced-affine-true-order}
      p(k'-k)-j k=(t+j b_0)/q_g\pmod g,
               \qquad jb_0+t=0\pmod{q_g}.
\end{equation}
At $b_0,j,t$ fixed, split the digit transform on the lift
variable at a dyadic modulus $g_0\mid g$ comparable to $P$.
For each pair $k,k'\bmod g_0$, the congruence determines $p$
modulo $g_0$ because their difference is odd. Thus only
$O(1)$ integers $p$ in the actual row box occur. Fixing one,
each remaining lift of $k$ determines the lift of $k'$ by
an affine permutation, since both $p$ and $p+j$ are odd.
The shifted square norms on these lifts cost $(g/g_0)^\delta$.
The two upper first norms at the coarse lift scale cost
$g_0^{2a}L^C$. If $g$ is nonoctal, retain the lower block
through its greatest octal divisor $g_-$ and the upper block
starting from its least octal multiple. The retained lower
digits split at $g/g_0$, choosing $g_0/(g/g_-)$ octal.
The coarse first norm is over a bounded number of aliases
on the $g_0$ grid, and the remaining square sum is unchanged.
The top blocks at size $q_g$ are then treated as in the
horizontal lemma. Consequently both estimates in that lemma
hold in this range upon replacing $g^{2a-1}$ by
$P^{2a-1}(g/P)^\delta$.

For these large orders $q\gg PH$ up to constants, so
$g\ll q/H\ll M$ and $q_g\gg H$. Both the separated and
dense contributions save with room. In the separated case
use $(g/P)^\delta\ll(M/P)^\delta$. In the dense case
combine $(g/P)^\delta q_g^\delta=(q/P)^\delta$ and use
$T/q_g\ll X^{O(\epsilon_0)}$. After the upper density cost,
both exponents are at most
$\delta_0\eta-(1+\delta_0-2a)\nu+O(\epsilon_0+\epsilon)<0$.

For the remaining orders $g\ll P$, the separated part
of the horizontal bound with upper density included costs
\[
 X^{\delta(\eta-\lambda)+O(\epsilon_0+\epsilon)}
                   g^{2a-1} B^{2\delta-1}L^C,
\]
where $B$ is the denominator shell of the reduced centres
and we used $T/H\ll X^{O(\epsilon_0)}$. Define
\begin{equation}\label{eq:near-balanced-order-cutoff}
     \Gamma(\nu,\eta)=.002+\frac{\delta_0}{1-2a}(\eta-\lambda).
\end{equation}
The parts $g>X^{\Gamma(\nu,\eta)}$ or $B>X^{.037}$ save
powers, because
$(1-2a)\Gamma(\nu,\eta)-\delta_0(\eta-\lambda)\ge.00047$ and
$(1-2(.0643))(.037)=.0322418>.0264273$.
In the dense case of that lemma, $B\ll H$ and $T/H\ll
X^{O(\epsilon_0)}$ force $g\gg q/H^2$ up to small powers.
Thus $g>X^{.34}$ and $T/q_g\ll X^{.0002+O(\epsilon_0)}L^C$,
because $g\ll P$. The complete square bound in the lemma
pays the dense case with a fixed margin.

It remains to remove, before the row bound, the lower source
frequency mask
\begin{equation}\label{eq:near-balanced-lower-band}
 \left|\frac{i}{q}-\frac{u}{gd}\right|_{mathbb R/\mathbb Z}
        \ll \frac{T}{dq},\qquad
               g\le C X^{\Gamma(\nu,\eta)}\text{ dyadic},\quad
               d\le C X^{.037}.
\end{equation}
Indeed both horizontal frequencies agree modulo $q_g$, and
the lower index condition in the horizontal lemma puts each
of them in this mask. Choose the fixed constants large
enough and include the half-period translates to preserve
odd physical support. Deleting this mask kills the indicated
terms in the remaining row Gram. We prove its comparison
directly with the unchanged upper word.

\paragraph{Fourier lifting of the lower mask.}
Group the reduced original centres in \eqref{eq:near-balanced-lower-band}
by the exact dyadic part $G$ of their denominator and the
dyadic size $D_0$ of its odd part. Also group $d\asymp B$.
Write $b=|\operatorname{den}(u/(gd))|\asymp GD_0$;
after reducing we may decrease $g,d$ and enlarge the bands.
Assign duplicates to one allowed representation of largest
width. Distinct centres are disjoint at scale $T/q$ by their
rational separation and the fixed upper bound $.1145+.037$
on their denominator exponent. For the squared mass keep an initial
block of length comparable to $G$, one of length comparable to $D_0^2$,
and a terminal block of length comparable to $1+T/B$. They fit below
$q$ by the fixed power bounds. Shifted Parseval applies to the initial
dyadic lifts and the terminal offset grid. After removing the dyadic
part, the odd fractions are separated by $\gg D_0^{-2}$. Sum the local
suprema of the squared middle amplitude using Bernstein's inequality,
as in Section~\ref{sec:zero-model-wide-source-bands}. The squared lower
mass in units of $\rho_L^2$ is consequently
\[
                  \ll G^\delta D_0^{2\delta}(1+T/B)^\delta L^C.
\]
Since $G/B\ll g$, $G\ll X^{.1145+.037}$ and
$T\le X^{.089+O(\epsilon_0)}$, its exponent is less than $.0180$
up to small completion losses. For $T\ge B$ the bound uses
$.0643(.1145+.089+2\cdot.037)<.0180$, and for $T<B$ it uses
$.0643(.1145+3\cdot.037)<.0180$.
For the first norm keep two middle blocks of length comparable
to $D_0^2$ and use Cauchy--Schwarz and the shifted square moment on
each. Octal dilation preserves the odd-denominator separation for
both blocks. The initial and terminal factors use shifted first norms. Thus the corresponding
first norm on a lower index radius $V\ll T/B$, enlarged to
$V\ge1$ when needed, is
\begin{equation}\label{eq:near-balanced-band-first-norm}
           G^a D_0^{2\delta} \max(1,V)^a L^C .
\end{equation}

Smooth the upper step word at scale $X^{-.020}R$ using
a compactly supported smooth approximate identity. In each
affected boundary interval, Cauchy and lower periodicity
bound the lower absolute mass by
\[
 (X^{-.020}R)^{1/2}(q+X^{-.020}R)^{1/2}
                       \rho_L X^{.0180/2+O(\epsilon_0)}L^C .
\]
There are $O((X/R)\rho_U)$ relevant boundaries in the
physical product range, including period repetitions.
The integer product divisor estimate makes the smoothing
error power-small relative to $X\rho_L\rho_U$.

Set $Q_R=Y/R$ and $Q_q=Y/q$. The upper step coefficient at
index $\ell$ is $\widehat U(\ell\bmod Q_R)$ times
\begin{equation}\label{eq:near-balanced-sinc-block}
 \frac1R\int_0^R e(-\ell u/Y)\,du
   = \left(\frac{q}{R}\sum_{j<R/q}e(-qj\ell/Y)\right)
                    \frac1q\int_0^q e(-\ell u/Y)\,du.
\end{equation}
The factor in parentheses is the normalized Fourier
transform of the unrestricted block from $q$ to $R$.
The smoothing multiplier cuts the tail past $Q_R X^{.025}$
to arbitrary power accuracy. The discrete coefficients on the integer
grid are the sums of the continuous coefficients in each class modulo
$Y$. The retained interval has length $o(Y)$ and the discarded aliases
have arbitrarily small total mass by the smoothing multiplier. Accordingly a full ordinary
frequency $\alpha=h/Y$ from one lower band remains within
effective offset $S=X\|\alpha-u/(gd)\|\ll X^{\nu+.0002+O(\epsilon_0)}/B$.
The padding exponent $1-\lambda+.025\le\nu-.042$ is
smaller than $\nu-.037$ throughout the range.

For the convolution, write $h=i_0 Q_q+r$ with $0\le r<Q_q$,
and in its lower coefficient put $i=i_0-j$. At
$\ell=r+jQ_q$ the last factor in
\eqref{eq:near-balanced-sinc-block} is $O((1+|j|)^{-1})$.
The other two upper factors are unchanged if $h$ is replaced
by $h-jQ_q$. The lower transform is evaluated at $(i_0-j)/q$,
within $1/q$ of that shifted argument divided by $Y$.
We can take local suprema on the retained lower digit blocks.
If $q$ is not octal, use its greatest octal divisor for the
last retained lower position and its least octal multiple
for the first retained unrestricted upper position, dropping
only bounded radix factors. The unrestricted geometric
factor and the upper word in
\eqref{eq:near-balanced-sinc-block} have complete first norm
$O(Q_q^a)L^C$ in their joint period and shifted local first
norm $O(V^a)L^C$ on any index interval of length $V\le Q_q$.
This follows by keeping their terminal block at length
comparable to $V$ for the local norm. For every shell of
effective offsets of size $S\ge1$ it follows that the
normalized first norm on the full grid is at most
\begin{equation}\label{eq:near-balanced-lift-first-norm}
                  G^a D_0^{2\delta}\max(1,YS/X)^a L^C .
\end{equation}
For clarity, put $V=\max(1,YS/X)$. A lower index within
$O(1+V/Q_q)$ of its centre costs the bound
\eqref{eq:near-balanced-band-first-norm} at that width,
and the convolution in the upper factor costs
$O(\min(V,Q_q)^a L^C)$ by the harmonic weight in $j$.
For lower index offsets $A\gg1+V/Q_q$, each term has
sinc weight $O(A^{-1})$. There are $O(1+V/Q_q)$ upper
period cells, with first norm $O(\min(V,Q_q)^a L^C)$ each.
Summing the lower first norm at each dyadic $A$ against
$A^{-1}$ gives the same bound, since $a<1$.

\paragraph{Resolving secondary rational centres.}
For the inner range $b(1+S)\ll P$, the short-row bilinear
Cauchy estimate of Section~\ref{sec:dynamic-rational-source-bands}
gives a pointwise normalized bound $(b(1+S))^{-1/2}L^C$,
using only the counting second norms of $u_x,v_y$.
In conjunction with \eqref{eq:near-balanced-lift-first-norm}
this pays with arbitrary logarithmic savings except at
polylogarithmic $G,D_0,S$, since $a<1/2$ and $2\delta<1/2$.
There the coordinate model estimate
\eqref{eq:near-balanced-small-model-arc} applies to the
actual coefficient, uniformly in the opposite physical
coordinate.

Choose a reduced Dirichlet approximant $a'/b'$ to the
actual full angle, with $b'\le P$ and error
$\ll1/(Pb')$. If this is the original centre and
$b(1+S)\gg P$, it has $bS\ll M$. In this case the
short Cauchy estimate gives $O((P^{-1}+(b(1+S))^{-1})^{1/2}L^C)$.
Indeed the rational short-difference classes have displacement
$O(L^C/b)$. The classes at larger distance from the integers
cost $O(P L^C)$ in the sum of column geometric sums.
In the polylogarithmically many exceptional classes the phases have
gap $bS/X$ and cross at most polylogarithmically many integers. Their
cost is $O(M+X/(b(1+S)))L^C$, using the trivial estimate
when $S<1$. This proves the bound. For $bS\gg P$ its
$P^{-1/2}$ factor pays \eqref{eq:near-balanced-lift-first-norm}
by $a(\Gamma(\nu,\eta)+.0002)+(a-.5)\nu<0$. The maximum
for each parameter choice is at $\eta=1$ and the lower
endpoint of $\nu$.

Otherwise the cross determinant is nonzero and gives
\begin{equation}\label{eq:near-balanced-cross-determinant}
               1\ll b b' S/X + b/P.
\end{equation}
The $b/P$ term is power-small. In a secondary shell
$b'\asymp B'=X^\sigma$, write $S=X^\tau$ up to rounding,
$G=X^{\Gamma_G}$ and $D_0=X^\Delta$. Then
$\sigma+\tau+\Gamma_G+\Delta\ge1-O(L^{-1})$.
The same short Cauchy estimate at each such approximant
is $O((B')^{-1/2})L^C$, since $b'\le P$ and the
approximation error has size at most $1/(P b')$. Put
$\sigma_0=\min(\nu-.001,.490,\eta-\nu-.001)$.
For $\sigma\ge\sigma_0$ the direct bound
\eqref{eq:near-balanced-lift-first-norm} gives exponent
\[
  a(\Gamma(\nu,\eta)+\nu+.0002+\eta-1)-\sigma_0/2
     +O(\epsilon_0)\le-.00989+O(\epsilon_0).
\]
For the three choices in the minimum defining $\sigma_0$,
the maximizing endpoint pairs $(\nu,\eta)$ are respectively
$(\nu_0,1),(.5,1),(.5,.740)$, where $\nu_0$ is the applicable
lower endpoint. This follows by using the two affine pieces
of $\lambda$ and their respective monotonicities.

For $\sigma<\sigma_0$, fix one original centre, the
secondary denominator shell and the common integer $j$
in the sinc expansion. By
\eqref{eq:near-balanced-cross-determinant} the distinct
secondary centres all lie in an interval of length $O(S/X)$
about the original centre, and their mutual spacing is
$\gg (B')^{-2}$. Each full frequency window around them
has length $O(V')$ in index units, where $V'=Y/(PB')$.
Put an octal start position $B_0\asymp X/S$ and keep a
block of octal length $N_0\asymp (B')^2 S/X$ there, as
well as a terminal block of length comparable to $V'$.
Both lengths are polynomially large in this range.
They fit disjointly since $B_0N_0V'\asymp YB'/P\ll Y$
with a fixed power margin. The original-centre bottom
blocks fit below $B_0$, since the exponent of $GD_0^4S$
is at most $.1145+.5002+3(.037)+O(\epsilon_0)<1$.

After dilation by $B_0$ the secondary centres lie in a
bounded real interval and have spacing $\gg N_0^{-1}$.
The common translate $-j/q$ preserves this property and
variation within their windows stays within the sampling
cells, since $B_0N_0V'/Y\ll B'/P$. A block entirely below
or above $q$ has sampled first norm $O(N_0^a)L^C$.
The lower rounding by $1/q$ changes its arguments by
at most the size of a sampling cell. If the block crosses
$q$, split it into retained lengths $N_L,N_U$ on either
side, dropping a bounded radix factor if necessary.
Partition the dilated real angles into cells of width
$1/N_L$. In each cell the upper arguments lie in a
bounded interval and are $\gg N_U^{-1}$ separated. The
upper sampled first norm is $O(N_U^a)L^C$. Sum the
suprema of the lower factor on the enlarged cells at
cost $O(N_L^a)L^C$. Thus the joint norm is again
$O(N_0^a)L^C$. The upper part uses the unrestricted
middle digits in \eqref{eq:near-balanced-sinc-block};
they have the same shifted first norm bound. The determinant bound also gives $\nu+\sigma\ge1-\Gamma(\nu,\eta)-.0002-o(1)>\lambda$,
so the terminal block begins beyond the lower period. It contributes
$(V')^a L^C$ on each
secondary integer frequency window. Sum the bottom
suprema over original centres as in
\eqref{eq:near-balanced-band-first-norm}, and sum the
harmonic weights in $j$. The normalized joint first norm
is bounded by
\[
         G^a D_0^{2\delta}(N_0V')^a L^C
             \le G^a D_0^{2\delta}(B'S/P)^a L^C .
\]
Multiplying by the pointwise $(B')^{-1/2}$ and using
\eqref{eq:near-balanced-cross-determinant}, together with
$\tau\le\nu+.0002-\log_X d+O(\epsilon_0)$ and
$\Gamma_G\le\log_X g+\log_X d-\Delta+O(L^{-1})$,
gives the exponent
\[
  (a-1/2)(1-\nu)+.0001+\tfrac12\log_Xg
                       +(2\delta-a)\Delta+O(\epsilon_0)
          \le-.00142+O(\epsilon_0).
\]
This completes the comparison of the removed horizontal
mask with the original upper word.

\paragraph{The small lower periods.}
For $q\le X^{.120}$ use all the lower frequencies, without
a mask. Their squared mass in lower density units is at most
$q^\delta L^C$. Smooth the upper step at $X^{-.009}R$;
the boundary error again saves a power since $.120\delta<.009$.
Truncating at $Q_R X^{.012}$ now leaves disjoint source
frequency windows of effective width at most $X^{\tau_0+o(1)}$,
where $\tau_0=1-\lambda+.012\le\min(.423,\nu-.053)$,
around $i/q$. In each such window there is exactly one
lower coefficient, with total normalized first norm
$O(q^a)L^C$ over all centres. For $q>2$ the exact conductor forces $i$ to be odd,
so $i/q$ has reduced denominator $q$. Apply the same inner and
secondary estimates with $G=q$, $D_0=1$ and effective
offset exponent at most $\tau_0$; at the bounded periods
there are only boundedly many denominator families.
The coincident-centre outer case uses
$a(.120+\tau_0)-\nu/2<0$. The direct large-secondary branch uses
$a(.120+\tau_0+\eta-1)-\sigma_0/2<0$, and the joint
small-secondary branch uses
$a(1-\nu)-1/2+(.120+\tau_0)/2<0$, uniformly with fixed
margins. For the direct estimate at $\sigma_0=\nu-.001$,
use $\tau_0\le\nu-.053$ and $\nu\ge.399$. For the other
two choices in $\sigma_0$, use $\tau_0\le .423$.
The latter sampling block starts past the lower
word and uses the unrestricted run to $R$ if needed.
This also includes the combined parity projection.

The small, untwisted, and twisted lower ranges overlap.
Summing their bounds over dyadic $q$ proves Proposition~\ref{octal:near-balanced-high-conductors}.
The exponent inequalities are checked in exact rational arithmetic by
\path{supplement/base8/verify_exponents.py}.
\end{proof}



\section{The grouped comparison for octal prime-factor regions}
\label{sec:octal-balanced-proof}

\begin{proof}[Proof of Lemma~\ref{octal:balanced-shape-comparison}]
Work first off fixed thin strips about the exponent boundaries. A
smooth partition over the possible groupings can be supported on
labelled prime subsets and their complements with scale exponents
at least $.399$ and at most $.601$, with logarithmic rounding
absorbed by the reserve. Express the weight as a finite sum in
these choices and insert the dyadic partitions in the prime scales.
Fourier inversion in the scaled logarithms of the variables separates
the smooth shape and product tests into smooth prime-slot weights,
as in Section~\ref{sec:principal-grouped-models}.
Reflection interchanges the two groups
if their order by size changes.

The distinctness condition can be discarded in such boxes. For
any of the nonnegative sources used below, use Type~I on the
divisors $p^2\le X^{.640}$ with $p\ge X^\kappa$, and use the
integer count $O(X^{1-.640/2}+X^{1/2})$ for larger squares.
The first main term is $O(A\sum_{p\ge X^\kappa}p^{-2})$, and
both bounds are $o(A/L)$. The fixed-rank representation count
and the interval cutoffs preserve that error.

Put $Q_*=2^{\lfloor .741\log_2 X\rfloor}$. First replace $f_0$
by $E_{X,Q_*}f$. We only need the classical fixed-modulus Ingham
density bound in this conductor range. For $q\le L^C$, fix all but one
prime and apply Siegel--Walfisz to the odd classes in that prime
column. The exact conductor contribution has zero mean on these
classes. The total absolute source mass over a period is
$O(Aq/X)$, and the reciprocal sum over fixed factor choices costs
only powers of $L$. For $L^C<q\le X^{c_0}$ with $c_0>0$ sufficiently
small in terms of $\kappa$, use the ordinary prime exponential
sum bound \cite[Theorem~23.8]{Koukoulopoulos2019} in the same column. In normalization
by the prime scale $N$, its three bounds are $q^{-1/2}$, $N^{-1/5}$, and
$(q/N)^{1/2}$, up to logarithms. The denominator remains $q$ since the
fixed factors are odd. The shifted digital first norm costs $q^{.3825}$.
Since $N\gg X^\kappa$, choosing $c_0$ sufficiently small gives arbitrary
inverse-logarithmic saving, on taking $C$ sufficiently large.

For larger $q$, apply the primitive character expansion
\eqref{octal:dyadic-char-expansion}. For each character in the
common zero layer, the smooth prime polynomial for a factor at
scale $N_i$ is bounded by $N_i^{\beta_\chi}L^{O(1)}$ by the smooth explicit formula. Use height $H=X^{3\epsilon}$ and let
$\beta_\chi$ be the maximum of $1/2$ and the real parts of zeros of
$L(s,\chi)$ at ordinates of absolute value at most $H$. The retained
Mellin twists have height at most $L^C$. Integration by parts in the
smooth test controls the tail above $H$ to any fixed power of $X$. The
unit-height zero count is $O(\log(q(3+|t|)))$, and the test decays at
least quadratically in the shifted zero ordinate, so the contribution
of retained zeros is $O(N_i^{\beta_\chi}L^{C_1})$. Prime powers cost
$O(N_i^{1/2}L^{C_1})$ and the primitive character is nonprincipal. The same common
$\beta_\chi$ is used for all factors and all the small Mellin
twists. Thus the product of the polynomials is bounded by
$X^{\beta_\chi}L^{O(1)}$, since the scales multiply to $X$ up to
powers of $L$.

Put $y=1-\beta$ for a layer with lower real-part endpoint $\beta$,
$\delta=1-\log_8 7$, $\sigma=.6175$ and $\lambda_0=(\sigma-\delta)/2$.
Ingham's fixed-modulus zero-density bound, in the form recorded in
\cite[(1.4)]{ChenGuptaLi2026}, gives at most
$X^{\gamma D_I(y)+O(\epsilon)}$ characters in the layer, where
$\gamma=\log_X q$ and $D_I(y)=3y/(1+y)$ for $0\le y\le1/2$.
Choose a finite real-part mesh of width at most $\epsilon$ and use its
upper endpoints in the prime bounds. The powers lost to the mesh and
to the height are $O(\epsilon)$. At the layer
with $\beta=1/2$ use the trivial character count, which agrees
with $D_I(1/2)=1$. The pointwise and square norms in
Lemma~\ref{octal:semiprime-digital-char} give the relative exponent
\[
 -y+\gamma\min\{D_I(y)-\lambda_0,(D_I(y)+\delta)/2\}.
\]
At fixed $y$ this is affine in $\gamma$, so it suffices to check
$\gamma=c_0$ and $\gamma=.742$. When $c_0<1/3$, the first term
of the minimum bounds the exponent above by $-\lambda_0c_0$,
since $D_I(y)\le3y$. At $.742$ the two branches meet at
$y_*=\sigma/(3-\sigma)$. The derivative in $y$ of the pointwise
branch is positive up to $y_*$, and that of the square branch
is negative from $y_*$ onward. These inequalities follow from
$3(.742)/2<(1+y_*)^2<3(.742)$. The maximum is at most
\[
    -y_*+.742(\sigma+.0643)/2<0.
\]
The strict savings absorb the mesh and height losses. Summing
over the smooth boxes and conductors pays $E_{X,Q_*}f-E_{X,2}f$
with error $o(A/L)$, and $E_{X,2}f=f_0$.

For the remaining source $g=f-E_{X,Q_*}f$, apply the principal models
of Proposition~\ref{octal:principal-grouped-model} after separation
into smooth prime-slot weights. The reconstructed continuous term in
\eqref{eq:principal-grouped-representation} is $o(A/L)$ on the fixed
smooth shape, by equality of the total masses and the endpoint bounds.
In the residual pairing split $g$ into its exact conductors $Y>Q_*$.
The grouped coefficients are the same for all these conductors.
Proposition~\ref{octal:near-balanced-high-conductors} bounds each such
pairing with arbitrary logarithmic saving. The scales have a fixed
positive margin in the group bounds used there. Summation over all
retained twists, dyadic boxes, and conductors gives $o(A/L)$.

Finally remove the strips and the intermediate smoothing. There
are only finitely many relevant subsets of the prime labels at fixed
$r$. The boundary at $.4$ or $.6$ can be marked by the indicated
subproduct or its complement. It has a fixed Type~I reserve. At
fixed roughness the marked upper sieve costs $O(A/(mL))$ for each
mark $m$ and has negligible total remainder, including on lifted
nonnegative sources by Lemma~\ref{octal:lifted-odd-divisor-marks}.
The reciprocal mass of a strip of width $\epsilon$ is
$O_\kappa(\epsilon)+o(1)$ by fixing all but one marked factor and
summing reciprocals over the last exponent interval. The same
argument applies at the individual roughness boundaries. Thus the
cost of the strips is $O_\kappa(\epsilon A/L)+o(A/L)$. Letting
$\epsilon$ tend to zero proves the estimate on the sharp region.
\end{proof}

\part{Asymptotics in bases ten and eleven}

\section{Digital hypotheses and affine correlations}
\label{dec:section}

The exact asymptotic is obtained by extending bilinear cancellation
strictly below the short-factor exponent $1/4$. The gain comes from a
correlation between two digital blocks linked by an affine permutation.
We first state the digital criterion that isolates the quantitative inputs,
and then prove the contraction principle behind this correlation.
Section~\ref{dec:bilinear-section} turns it into a minor-arc estimate;
Section~\ref{dec:prime-section} compares the digital set with a single
nonnegative model and applies the Ford--Maynard asymptotic sieve.

Throughout these three sections, fix $b\in\{10,11\}$ and a forbidden
digit $d$. Write $\mathcal D_k=\mathcal W_k$ for the suffix words of
Section~\ref{sec:preliminaries}, and put
\[
 q=b^k,\qquad A=(b-1)^k=q^\alpha,\qquad
 \alpha=c_b=\frac{\log(b-1)}{\log b},\qquad f=f_{b,d}.
\]
We use $\widehat u(t)=\sum_nu(n)e(nt)$; for a function on
$\mathbb Z/q\mathbb Z$, the sum runs over $0\le n<q$.
Thus $F_q(t)=A^{-1}|\widehat{1_{\mathcal D_k}}(t)|$.
Write $e_q(t)=e(t/q)$. Fixed powers of logarithms, divisor bounds,
and factors $\exp(O(\sqrt{\log q}))$ are absorbed in $q^{o(1)}$.
All parameters are fixed before $q\to\infty$.

\subsection{A criterion for the exact asymptotic}
\label{dec:inputs}

The quantitative argument uses the parameters
\begin{equation}\label{dec:constants}
 \gamma_0=\frac{27}{77},\quad
 p_0=\frac{235}{154},\quad g_0=\frac{59}{433},\quad
 \varepsilon_0=\frac{23}{500},\quad g_4=\frac3{5000},\quad
 \kappa=\frac{23}{80}.
\end{equation}
The exponents $\gamma_0=27/77$ and
$(p_0,g_0)=(235/154,59/433)$ are Maynard's decimal inputs
\cite[Lemmas 10.3--10.4]{Maynard2019}.
The additional $3/2$ moment is paired with an affine-correlation
exponent $\delta$; both will enter the reverse digit decomposition.
For every $Y=b^j$ with $j\ge0$, the moment estimates are
\begin{align}
 \sup_x\sum_{a<Y}F_Y((a+x)/Y)&\ll Y^{\gamma_0},\label{dec:moment1}\\
 \sup_x\sum_{a<Y}F_Y((a+x)/Y)^{p_0}&\ll Y^{g_0},\label{dec:momentold}\\
 \sup_x\sum_{a<Y}F_Y((a+x)/Y)^{3/2}&\ll Y^{g},\label{dec:momentnew}\\
 \sum_{a<Y}F_Y((a+x)/Y)^2&=Y^{1-\alpha},\qquad
 \sup_x\sum_{a<Y}F_Y((a+x)/Y)^4\ll Y^{g_4}.
 \label{dec:evenmoments}
\end{align}
where
\begin{equation}\label{dec:groups}
 (g,\delta)=
 \begin{cases}
 (14387/100000, 53/10000),&b=10, d\in\{4,5\},\\
 (287/2000, 1/250),&\text{otherwise}.
 \end{cases}
\end{equation}
The second and fourth moments in \eqref{dec:evenmoments} follow from
\eqref{prelim:even}. The scalar inequalities
$1/100<1-\alpha<\varepsilon_0$ and $\alpha>20/21$ hold in both bases;
the exact comparisons are recorded in Appendix~\ref{app:exact-inputs}.

For the digit laws, put
\[
 P_i(x)=\frac1{b(b-1)}
          \left|\sum_{\substack{0\le j<b\\j\ne d}}e(j(i+x)/b)\right|^2,
 \qquad \sum_{i=0}^{b-1}P_i(x)=1,
\]
and define
\begin{align*}
 H(x,E)&=\sum_{i=0}^{b-1}\sqrt{P_i(x)P_i(x+E)},
 &\rho(E)&=\sup_{0\le x\le1}H(x,E),\\
 K_t&=\sup_{\substack{x\in[0,1],\,c\in[0,b]\\a\in\{1,\ldots,b-1\}}}
 \sum_{i=0}^{b-1}P_i(x)\rho(ai+c)^t.
\end{align*}
The function $H$ is one-periodic in $x$, while $\rho$ is even and
$b$-periodic. 
The exponent $\delta$ measures the contraction of affine correlations.
The following criterion separates that input from the bilinear and sieve
arguments.

\begin{theorem}[Digital criterion for bases ten and eleven]
\label{dec:digital-criterion}\label{dec:asymptotic}
For $b\in\{10,11\}$ and a forbidden digit $d$, take the pair
$(g,\delta)$ in \eqref{dec:groups}. Assume that
\eqref{dec:moment1}--\eqref{dec:momentnew} hold uniformly in the shift
and in every power $Y$ of $b$, and that
\begin{equation}\label{dec:one-digit-hypothesis}
 K_2\le b^{-3\delta}.
\end{equation}
Then, for ordinary base-$b$ expansions and arbitrary real $X\to\infty$,
\[
 \pi_{b,d}(X)\sim\kappa_{b,d}\frac{A_{b,d}(X)}{\log X},
 \qquad
 \kappa_{b,d}=\frac{b\bigl(\varphi(b)-\mathbf1_{(d,b)=1}\bigr)}
                    {(b-1)\varphi(b)}.
\]
\end{theorem}

For the decimal central digits, the stronger affine contraction
compensates for the slightly weaker $3/2$ moment. The estimates use
each pair in \eqref{dec:groups} jointly.
In addition to the stated hypotheses, the proof uses the
digital product structure, the elementary even moments, and the
one-sided shift comparison of Lemma~\ref{prelim:shift}.
The latter is proved for these deleted-digit products in
Appendix~\ref{app:shifts}.

\begin{proposition}[Verification of the digital hypotheses]
\label{dec:verified-inputs}\label{eleven:inputs}
For every deleted digit in each of bases ten and eleven,
\eqref{dec:moment1}--\eqref{dec:momentnew} and
\eqref{dec:one-digit-hypothesis} hold with the pair
$(g,\delta)$ in \eqref{dec:groups}.
\end{proposition}

Appendix~\ref{app:exact-inputs} proves this proposition by continuous
positive-operator inequalities and one-digit correlation bounds, with
exact rational certificates. Theorem~\ref{dec:digital-criterion} is
proved in Section~\ref{dec:prime-section}; together these statements
will prove Theorem~\ref{main:exact}.

\subsection{Two-ended interpolation on probability laws}

A row-by-row Hellinger estimate produces a product along a digital path.
The next lemma converts bounds for that product under both endpoint laws
into a bound for the complete correlation.

\begin{lemma}[Two-sided probability interpolation]\label{dec:twoended}
Let $P,Q,R$ be probability laws on a finite set, and let $0\le U\le1$.
Suppose that $\sqrt{P(\omega)Q(\omega)}\le R(\omega)U(\omega)$ and,
for some $t>0$,
\[
 \mathbb E_P U^t\le\eta,\qquad \mathbb E_Q U^t\le\eta.
\]
Then
\[
 \sum_\omega\sqrt{P(\omega)Q(\omega)}
 \le\eta^{1/(1+t)}.
\]
\end{lemma}
\begin{proof}
Put $a=1/(1+t)$ and $w=\sqrt{PQ}$. The pointwise inequalities imply
\[
 w=w^a w^{1-a}
 \le(PU^t)^{a/2}(QU^t)^{a/2}R^{1-a},
\]
because $at=1-a$. H\"older's inequality with exponents
$2/a,2/a,1/(1-a)$ gives the assertion.
\end{proof}

\subsection{Affine contraction with exact carries}

An affine map transports the successive digit laws, so its correlation
can be evaluated on the two paths simultaneously. The number of levels
that contribute a gain is determined by the first digit at which the
slope differs from $1$.

For $B=b^L$, put
\[
 \mu_{B,x}(z)=((b-1)/b)^L F_B((z+x)/B)^2,
 \qquad z\in\mathbb Z/B\mathbb Z.
\]
This is a probability law whose conditional digit probabilities, read from
the least significant digit upwards, are $P_i(x_j)$, with
$x_0=x$ and $x_{j+1}=(i_j+x_j)/b$.
For a unit $\lambda$ modulo $B$ and an integer $c$, write
\[
 \mathcal H_L(\lambda,c;x,y)
 =\sum_{z\bmod B}\sqrt{\mu_{B,x}(z)\mu_{B,y}(\lambda z+c)}.
\]
We use $v_b(n)=\max\{j:b^j\mid n\}$, with $v_b(0)=\infty$;
this definition also applies when $b$ is composite.

\begin{lemma}[Affine Hellinger contraction]\label{dec:hellinger}
Let $\lambda$ be a unit modulo $B$ with $\lambda\equiv1\pmod b$, and let
$m=\min(L,v_b(\lambda-1))$. For every integer $c$, every $x,y\in[0,1]$,
and every $t>0$,
\begin{equation}\label{dec:hellingerbound}
 \sum_{z\bmod B}\sqrt{\mu_{B,x}(z)\mu_{B,y}(\lambda z+c)}
 \le K_t^{(L-m)/(1+t)}.
\end{equation}
In particular, the numerical input $K_2\le b^{-3\delta}$ gives the
upper bound $(B/b^m)^{-\delta}$.
\end{lemma}
\begin{proof}
For $m=L$, apply Cauchy's inequality. Otherwise write
$\lambda-1=b^ma$, where $b\nmid a$. Along the source path, write
\[
 r_j=z\bmod b^j,\quad r'_j=(\lambda r_j+c)\bmod b^j,
 \quad C_j=(\lambda r_j+c-r'_j)/b^j.
\]
At digit $j$, the permutation is $i\mapsto i+C_j\pmod b$.
Its Hellinger row sum is at most $\rho(E_j)$, where
\[
 E_j=C_j+\frac{r'_j+y-r_j-x}{b^j}
 =\frac{(\lambda-1)r_j+c+y-x}{b^j}.
\]
When $j=m+h$, we have
\[
 E_{m+h}\equiv ai_h+\frac{ar_h}{b^h}
                  +\frac{c+y-x}{b^{m+h}}\pmod b.
\]
Since $m\ge1$, this congruence determines $E_{m+h}$ modulo $b$
from $i_0,\ldots,i_h$. Given the preceding digits, averaging its
factor $\rho(E_{m+h})^t$ over $i_h$ costs at most $K_t$.
Conditioning in decreasing order of $h$ gives
\begin{equation}\label{dec:sourceexpect}
 \mathbb E_{\mathrm{src}}\prod_{h=0}^{L-m-1}\rho(E_{m+h})^t
 \le K_t^{L-m}.
\end{equation}
For the reverse calculation, choose integers $\lambda',c'$ with
$\lambda'\lambda\equiv1\pmod B$ and $c'\equiv-\lambda'c\pmod B$.
The first nontrivial digit of $\lambda'$ is again $m$.
At the corresponding source and destination prefixes, the reverse error
$E'_j$ satisfies
\[
 E'_j+E_j
 =(1-\lambda')C_j+
 \frac{(\lambda'\lambda-1)r_j+c'+\lambda'c}{b^j}
 \in b\mathbb Z\qquad(j<L).
\]
This identity includes the real initial states $x,y$.
Since $\rho$ is even and $b$-periodic, the reverse conditioning
argument gives \eqref{dec:sourceexpect} under the destination law
with exactly the same product.

To finish, set $P(z)=\mu_{B,x}(z)$,
$Q(z)=\mu_{B,y}(\lambda z+c)$ and
$U(z)=\prod_{h=0}^{L-m-1}\rho(E_{m+h})$.
Normalize each row of the geometric-mean transition kernel by its
row sum $H_j(z)$. This produces a probability law $\mathbb R$ on
paths with
\[
 \sqrt{P(z)Q(z)}=\mathbb R(z)\prod_{j=0}^{L-1}H_j(z)
 \le\mathbb R(z)U(z).
\]
A zero row may be completed to any probability row; the identity
remains valid because its unnormalized path weight is zero.
Apply Lemma~\ref{dec:twoended} with $\eta=K_t^{L-m}$.
\end{proof}


\section{Bilinear cancellation below one quarter}
\label{dec:bilinear-section}

We assume the digital hypotheses of Theorem~\ref{dec:digital-criterion}.
The objective of this section is the following minor-arc estimate for
arbitrary bounded coefficients. The proof combines the affine contraction
with rational-grid estimates and the classification of large bilinear
sums developed by Maynard \cite[Sections 10, 13--15]{Maynard2019}.

For $H\ge1$, let $\mathfrak M(q;H)$ be the set of indices
$a\in\mathbb Z/q\mathbb Z$ for which there are integers $c,r$ with
$1\le r\le H$ and $(c,r)=1$ such that
\[
 \left\|\frac aq-\frac cr\right\|\le\frac Hq.
\]
For a set $\mathfrak M$ of frequency indices define
\[
 (P_{\mathfrak M}u)(n)
 =\frac1q\sum_{a\in\mathfrak M}\widehat u(a/q)e_q(-an).
\]

\begin{proposition}[Minor-arc bilinear estimate]\label{dec:minorarc}
Fix an admissible positive prefix $v$ and $T>0$.
There is $D_0>0$ such that, with
$\mathfrak M=\mathfrak M(q;(\log q)^{D_0})$, one has
\begin{equation}\label{dec:minorarcII}
 \begin{split}
 &\sum_{m\in I_M,r\in I_R}\alpha_m\beta_r
 1_{\{vq\le mr<(v+1)q\}}1_{(mr,b)=1}
 ((I-P_{\mathfrak M})1_{\mathcal D_k})(mr-vq)\\
 &\hspace{35mm}\ll_{v,T}A\log^{-T}q
 \end{split}
\end{equation}
whenever $MR\asymp_vq$, $q^{0.249999}\le R\le q^{0.4244}$,
$I_M\subset[M,2M]$ and $I_R\subset[R,2R]$ are intervals, and
$|\alpha_m|,|\beta_r|\le1$.
\end{proposition}

The three ranges in \eqref{dec:rawwindow} correspond to different
bilinear estimates. For intermediate amplitudes, the positive kernel
below separates into a low-digit correlation and a high-digit
arithmetic estimate. Reversing the digit splitting provides the
correlation gain required for short factors below exponent $1/4$.

\subsection{The positive kernel and short integer products}
\label{dec:shortproduct}

For a nonnegative frequency weight $w$, define
$W=\sum w$, $V_w=\sum w^2$ and $|u|_q=q\|u/q\|$, and put
\[
 \mathcal N_q(w;R)=\sum_{r,s\in I_R}\sum_{a,c\bmod q}w(a)w(c)
 1_{\{|ra-sc|_q\le C_0R\}},\quad I_R\subset[c_1R,c_2R].
\]
The multiplier set $I_R$ may be restricted, in particular to units
modulo $b$.

For arbitrary $|\alpha_r|,|\beta_m|,|\gamma_a|\le1$ and
$M\ll_vq/R$, Cauchy's inequality in the long variable gives
\begin{align}
 \left|\sum_aw(a)\gamma_a\sum_{r,m}\alpha_r\beta_m e_q(-arm)\right|^2
 &\ll_v\frac qR\mathcal K_q(w;R),\label{dec:cauchy}\\
 \mathcal K_q(w;R)&=\sum_{a,c,r,s}w(a)w(c)
 \min(q/R,\|(ra-sc)/q\|^{-1}).\notag
\end{align}
Partition the output into cells of length $R$, and apply cyclic Cauchy
and harmonic summation to the distances between cells. This gives
\[
 \mathcal K_q(w;R)\ll(q/R)(\log q)\mathcal N_q(w;R).
\]
Thus the intermediate-amplitude bound
$\mathcal N_q\ll R^2q^{-c}$ bounds the original Fourier pairing
by $O_v(q^{1-c/2}\sqrt{\log q})$.
Returning to the digital indicator introduces the factor $A/q$.
The following observation permits decompositions of the source weight
in the complete output norm, including all cross terms.

\begin{lemma}[Positive output norms]\label{dec:outputnorm}
Let $1\le R\le q$. For nonnegative weights
$w,w_1,\ldots,w_J$ on $\mathbb Z/q\mathbb Z$,
if $w\le\sum_{j=1}^Jw_j$, then
\[
 \mathcal N_q(w;R)\ll J\sum_{j=1}^J\mathcal N_q(w_j;R).
\]
The same assertion holds for $\mathcal K_q$ in \eqref{dec:cauchy}.
The constants depend only on the fixed window constants.
Both statements concern the complete pair sum over the chosen
multiplier set $I_R$.
\end{lemma}
\begin{proof}
Use the actual output mass
$\mu_w(z)=\sum_{r,a:ra=z\bmod q}w(a)$.
Partition the circle into consecutive cells of length comparable to $R$
and at most $C_0R$. Then
\[
 E_R(w)=\sum_I\left(\sum_{z\in I}\mu_w(z)\right)^2
 \ll\mathcal N_q(w;R)\ll E_R(w).
\]
The first inequality counts pairs in the same cell. For the second,
a pair at distance at most $C_0R$ lies in one of a bounded number of
neighboring cell pairs, and Cauchy's inequality applies.
The map from $w$ to its cell masses is linear and positive. Hence
$E_R(w)\le J\sum_jE_R(w_j)$.

For $\mathcal K_q$, use the circle measure
$\widetilde\mu_w=\sum_z\mu_w(z)\delta_{z/q}$, put $Q_*=q/R$, and take
\[
 H(t)=\sum_{\substack{1\le L\le Q_*\\L\text{ a power of }2}}
 \frac1L\left|\sum_{n<L}e(nt)\right|^2.
\]
The geometric-series formula gives
$H(t)\asymp\min(Q_*,\|t\|^{-1})$: for the lower bound choose
$L\asymp\min(Q_*,\|t\|^{-1})$ small enough that $L\|t\|\le1/4$,
and use $L=1$ if necessary. Each summand has nonnegative Fourier
coefficients. Thus
\[
 E_H(\mu)=\iint H(t-u)\,d\mu(t)\,d\mu(u)
\]
is a squared Hilbert norm. Since $H\ge0$, the pointwise majorization
of the weights first gives
$E_H(\widetilde\mu_w)\le E_H(\sum_j\widetilde\mu_{w_j})$;
the Hilbert norm inequality then gives
$E_H(\widetilde\mu_w)\le J\sum_jE_H(\widetilde\mu_{w_j})$.
Pointwise comparison with $H$ proves the second assertion.
\end{proof}

In applications, pair restrictions enter only after this decomposition,
so the norm comparison always concerns the complete multiplier sum.

\begin{lemma}[Short products and congruent multipliers]\label{dec:fejer}
For $1\le V\le q$, we have
\begin{equation}\label{dec:shortbound}
 \sum_{1\le|r|,|s|\le C_1V}\sum_{a,c}w(a)w(c)
 1_{\{|ra-sc|_q\le C_2V\}}
 \ll q^{o(1)}V^2V_w+\frac{V^3W^2}{q}.
\end{equation}
If only pairs with $r\equiv s\pmod D$ are retained, where $D\le R$, then
\begin{equation}\label{dec:deep}
 \mathcal N_{q,D}(w;R)
 \ll q^{o(1)}R(1+R/D)(V_w+RW^2/q).
\end{equation}
\end{lemma}
\begin{proof}
Choose a sufficiently large fixed $C_3$.
If $V\ge q/(2C_3)$, the trivial bound $O(V^2W^2)$ is already
$O(V^3W^2/q)$. Otherwise put $L=\lfloor q/(C_3V)\rfloor\ge1$.
The normalized Fej\'er kernel
\[
 H_L(x)=L^{-2}\left|\sum_{j<L}e_q(jx)\right|^2
 =\sum_{|t|<L}\frac{L-|t|}{L^2}e_q(tx)
\]
has a fixed positive lower bound on $|x|_q\le C_2V$.
The zero frequency contributes $O(V^3W^2/q)$.
For the nonzero frequencies, apply Cauchy's inequality in $r$.
Since $0<|rt|<q/4$, each product is an actual nonzero integer with at most
$2\tau(|rt|)$ factorizations. Hence
\[
 \sum_{r,\,0<|t|<L}|\widehat w(rt/q)|^2
 \le q^{o(1)}\sum_{\xi\bmod q}|\widehat w(\xi/q)|^2=q^{1+o(1)}V_w.
\]
Multiplication by $V/L\ll V^2/q$ gives \eqref{dec:shortbound};
for $V\asymp q$, the trivial estimate suffices.
Under the restriction $r\equiv s\pmod D$, the multiplier graph has degree
$O(1+R/D)$. Replacing the preceding application of Cauchy's inequality
by the Schur bound for this graph gives \eqref{dec:deep}.
\end{proof}

When $w\le\tau$, the same lemma bounds the total mass of all
relations of small height:
\begin{equation}\label{dec:relation}
 \sum_{a,c}w(a)w(c)
 1_{\{v_1a+v_2c+v_3q+v_4=0\text{ for some }0\ne v\in[-V,V]^4\}}
 \ll q^{o(1)}(\tau WV^2+W^2V^3/q).
\end{equation}
For $v_1v_2\ne0$, apply \eqref{dec:shortbound} directly.
If one of the first two coefficients is zero, fix $(v_3,v_4)$ and use
the divisor bound for nonzero integer products. When the corresponding
source point is zero, its contribution is at most $\tau W$.
If both of the first two coefficients vanish and $V<q$, no nonzero
relation is possible. For $V\ge\sqrt q/2$, the inequality
$W^2\le q\tau W\ll\tau WV^2$ suffices.

\subsection{Forward carries and the low-digit second moment}
\label{dec:forward}

The forward splitting turns a low-digit incidence relation into a
bounded-degree graph. We work with amplitude subweights: at exponent $u$,
these are functions satisfying
\[
 0\le w(a)\le F_q(a/q)
       1_{\{q^{-u}\le F_q(a/q)<2q^{-u}\}}.
\]
Replacing the factor $2$ by another fixed constant has no effect.
This convention keeps the amplitude constraints attached to actual
nonempty digit layers, even when a smaller subweight is estimated.

Write $R=q^\nu$. Take $q=BQ$, with $Q,B$ powers of $b$,
$Q=q^{\vartheta}$, and $\vartheta=10\nu/3+O(1/\log q)$.
Write $a=z+BA$.
The exact factorization is
\begin{equation}\label{dec:forwardfactor}
 F_q(a/q)=F_B(z/B)F_Q((A+z/B)/Q).
\end{equation}
On an actual amplitude layer $F_q\asymp q^{-u}$ with $u\le\kappa$,
we have
\begin{equation}\label{dec:shift}
 F_Q((A+z/B)/Q)\le D_qF_Q(A/Q),\qquad
 D_q=\exp(O(\sqrt{\log q})).
\end{equation}
Indeed, the high factor in \eqref{dec:forwardfactor} is at least
$F_q(a/q)\gg q^{-\kappa}$, and $z/B\in[0,1)$.
Lemma~\ref{prelim:shift} therefore compares the shifted value with
the unshifted grid value at a cost $\exp(O(\sqrt{\log q}))$,
uniformly in $A,z$. Its proof in Appendix~\ref{app:shifts} uses the
unit-circle zeros of the one-digit polynomial to control the digit
positions at which such a comparison could otherwise be large. The comparison is used only on
this actual amplitude layer.

Decompose according to the low amplitude $F_B\asymp q^{-a_0}$, and put
\[
 l(z)=F_B(z/B)1_{\{q^{-a_0}\le F_B(z/B)<2q^{-a_0}\}},
 \quad \beta=\frac{q^{-u}}{C D_qq^{-a_0}},
 \quad f_Q^*(A)=\beta1_{\{F_Q(A/Q)\ge\beta\}}.
\]
The original weight is bounded by $q^{o(1)}$ times these low--high
products, and $x=u-a_0+o(1)$. The constant-height superlevel weight
ensures both pointwise bounds in \eqref{dec:highweight}.
For fixed $r,s,A,A'$, the original near-resonance condition
\[
 |B(rA-sA')+rz-sz'|_q\ll R
\]
first forces $|rA-sA'|_Q\ll R$. For each fixed $z$, the condition
$s\asymp R$ leaves only $O(1)$ possible values of $z'$, and conversely.
The low-digit bipartite graph therefore has bounded degree on both sides.
More explicitly, fixing $z$ confines $z'$ to intervals of bounded
length, separated by $q/s\gg B$, and only a bounded number meet
$[0,B)$. The same argument with $r,s$ exchanged bounds the opposite
degree. Schur's inequality gives
\begin{equation}\label{dec:forwardcarry}
 \mathcal N_q(l\otimes f_Q^*;R)
 \ll\|l\|_2^2\mathcal N_Q(f_Q^*;R).
\end{equation}
Put $\ell=\log_q\|l\|_2^2$ and $s=\log_q\sum_Af_Q^*(A)$.
Up to $o(1)$ relaxations, put $x=u-a_0$ and $\vartheta=10\nu/3$.
The full continuous parameter region needed is
\begin{align}
 &0.257\le\nu\le0.29,\quad \nu-0.07\le u\le\kappa,
 \quad 0\le a_0\le u,
 \label{dec:forwardregion}\\
 &-2a_0\le\ell\le\min\{\varepsilon_0(1-\vartheta),
        g_0(1-\vartheta)-(2-p_0)a_0,
        g_4(1-\vartheta)+2a_0\},\label{dec:lowconstraints}\\
 &-x\le s\le\min\{\gamma_0\vartheta,
        g_0\vartheta+(p_0-1)x,
        \varepsilon_0\vartheta+x,
        g_4\vartheta+3x,
        \vartheta-x\}.\label{dec:highconstraints}
\end{align}
The upper bounds come from the first, second, $p_0$th, and fourth
moments and the amplitude cap; the lower bounds for nonempty layers
come from the presence of at least one actual atom.
In the rank-two inequalities of Section~\ref{dec:ranktwo}, we will
take $p=p_0,g=g_0$.

\subsection{Reverse carries: crossing the one-quarter threshold}
\label{dec:reverse}

Now let $0.249999\le\nu\le0.258$, still with
$\vartheta=10\nu/3$ and $q=BQ$, but write $a=A+Qz$.
The product formula becomes
\begin{equation}\label{dec:reversefactor}
 F_q(a/q)=F_Q(A/Q)F_B((z+A/Q)/B).
\end{equation}
Decompose the high factor according to its actual amplitude $q^{-x}$.
The low amplitude is $q^{-a_0}$, where $a_0=u-x$.
The low profile $l_A(z)$ may depend on $A$, and the initial state
$A/Q$ is retained throughout. For each contributing $r,s,A,A'$,
write the unique centered high error as
$rA-sA'=Qk+e$, with $|e|\ll R<Q/2$.
Then the original near-resonance condition is equivalent to
\begin{equation}\label{dec:actualreversecarry}
 |e|\ll R,\qquad rz-sz'+k\equiv0\pmod B.
\end{equation}
Since the multipliers are units, the low-digit relation is an actual
affine permutation. Hence
\begin{equation}\label{dec:lowcross}
 \sum_zl_A(z)l_{A'}((r/s)z+k/s)
 \le\min\{V_{
 \rm low},B^{1-\alpha}
 \mathcal H_L(r/s,k/s;A/Q,A'/Q)\},
\end{equation}
Here the profiles obey
\[
 0\le l_A(z)\le F_B((z+A/Q)/B)
       1_{\{F_B((z+A/Q)/B)\asymp q^{-a_0}\}},
\]
and $V_{\rm low}$ is the supremum, over the actual initial states
$A/Q$, of the squared mass of the full low-amplitude band.
The first bound in \eqref{dec:lowcross} is Cauchy's inequality.
For the second, discard the amplitude restrictions and use
$F_B((z+x)/B)=B^{(1-\alpha)/2}\mu_{B,x}(z)^{1/2}$.
Thus the factor $B^{1-\alpha}$ and both initial states are exact.

First split the original output measure according to the
$\varphi(b)$ unit classes of the multiplier modulo $b$.
The Hilbert norm argument in Lemma~\ref{dec:outputnorm} costs only
this fixed number of classes. Within a class, $r/s\equiv1\pmod b$,
so Lemma~\ref{dec:hellinger} applies. Put $\ell_0=\log_qV_{\rm low}$.
We first majorize each subweight by the full nonempty amplitude band,
which justifies the constraint $\ell_0\ge-2a_0$.
Retain every constraint in \eqref{dec:lowconstraints}, with $\ell$
replaced by $\ell_0$, together with \eqref{dec:highconstraints}, and add
\begin{equation}\label{dec:newconstraints}
 \ell_0\le g(1-\vartheta)-a_0/2,\qquad s\le g\vartheta+x/2.
\end{equation}
Take $\eta=10^{-6}$, and choose a power of $b$ of the form
\begin{equation}\label{dec:deepchoice}
 D_*=q^{d_*+o(1)},\qquad d_*=\max(0,\ell_0+s-x)+\eta.
\end{equation}
When $D_*\mid r-s$, Lemma~\ref{dec:fejer} bounds the contribution
of these deep slopes by
\begin{equation}\label{dec:deepcost}
 \frac{\mathcal N_{\rm deep}}{R^2}
 \ll q^{\ell_0-d_*+o(1)}
       (q^{s-x}+q^{\nu+2s-\vartheta}).
\end{equation}
The first term is at most $q^{-\eta+o(1)}$.
The remaining slopes satisfy $v_b(r-s)<\log_bD_*$, so Hellinger
contraction acts on at least $\log_b(B/D_*)$ levels.
Equation \eqref{dec:lowcross} gives an effective low-digit load
$q^\ell$, where
\begin{equation}\label{dec:effectiveell}
 \ell\le\ell_0,\qquad
 \ell\le\varepsilon_0(1-\vartheta)-\delta(1-\vartheta-d_*).
\end{equation}
We take the minimum of these two upper bounds; the relative Hellinger
saving is not multiplied by $V_{\rm low}$.
All mass parameters and $D_*$ are chosen from the original full nonempty
amplitude bands and then held fixed. Restricting a source set decreases
the sums to be estimated; it does not change these parameters or require
a new lower bound on the mass of the restricted set.
To justify source deletion in this setting, perform the source
decomposition of Lemma~\ref{dec:outputnorm} in the original output
norm, before imposing any slope restriction.
Explicitly, a set $E$ of high frequencies lifts to the subweight
$w_E(A+Qz)=w(A+Qz)1_E(A)$.
For each resulting subweight, the deep contribution is bounded by
\eqref{dec:deepcost}, using the full masses as upper bounds.
For every other slope, \eqref{dec:lowcross} bounds the low cross term
by $q^\ell$, and enlarging the high multiplier sum gives the complete
kernel $\mathcal N_Q$ for that source subweight.
The estimates of Section~\ref{dec:ranktwo} will therefore apply to
this complete kernel.
Thus the deletion estimate follows from the positive complete kernels
$\mathcal N_q$ and $\mathcal N_Q$.

\subsection{From digital moments to arbitrary rational grids}
\label{dec:grids}

Suppose that $1\le p\le2$ and
$\sup_x\sum_{a<Y}F_Y((a+x)/Y)^p\ll Y^g$.
These are the base-$b$ sampling estimates underlying
\cite[Lemmas 10.5--10.7]{Maynard2019}; we keep the common shift explicit.
Write $F_Y=|P_Y|$ for the normalized digital trigonometric polynomial.
Integration over the shift, followed by
$\int|(|P_Y|^p)'|\le p\|P_Y\|_p^{p-1}\|P_Y'\|_p$
and Bernstein's inequality, gives
\begin{equation}\label{dec:variation}
 \int_0^1F_Y^p\ll Y^{g-1},\qquad
 \int_0^1|(F_Y^p)'|\ll Y\int_0^1F_Y^p.
\end{equation}
Consequently, a local supremum may be taken independently in each sampling
cell. Suppose that the centers have mutual distance $\gg1/Z$ and that
the integer neighborhood of each center has width $O(E)$, where $Z,E\ge1$.
Writing $\mathcal U$ for the union of these neighborhoods on the integer
grid modulo $Y$, we obtain
\begin{equation}\label{dec:separated}
 \sum_{a\in\mathcal U}F_Y(a/Y)^p\ll\min(Y^g,(ZE)^g).
\end{equation}
When $ZE<Y$, choose powers $U,V$ of $b$ with
$U\asymp Z$, $V\asymp E$, and $UV\le Y$. The product formula gives
\[
 F_Y(a/Y)^p\le F_U(a/Y)^p F_V(a/V)^p.
\]
Each neighborhood has diameter $O(1/U)$ in the first factor.
By \eqref{dec:variation}, the sum of its local suprema over the
separated centers is $O(U^g)$. In the second factor, a neighborhood
contains only $O(1)$ periods of the integer grid modulo $V$,
and contributes $O(V^g)$. This proves \eqref{dec:separated}.
For $ZE\ge Y$, use the full-grid moment.
If representations of the centers are counted with multiplicity, each
point has multiplicity at most $O(1+ZE/Y)$, giving the corresponding bound
$(ZE)^g+ZEY^{g-1}$. In particular, for $w\le\min(F_Y,\tau)$,
\begin{equation}\label{dec:squarearc}
 \sum_{a\in\mathcal U}w(a)^2
 \ll\tau^{2-p}\min(Y^g,(ZE)^g).
\end{equation}

We also need two mixed estimates. The first uses two equal-sized digit
blocks. Write an arbitrary denominator as $s=S v$, where every prime
factor of $S$ divides $b$ and $(v,b)=1$. For a common shift $\beta$,
\begin{equation}\label{dec:adaptiverow}
 \sum_{j<s}\sum_{\|a/Y-\beta-j/s\|\le E/Y}w(a)
 \ll (1+E)\tau^{1-p/2}s
 \min\{v,(Y/(1+E))^{1/2}\}^{g-1}.
\end{equation}
To prove this, choose a power of $b$ with
$U\asymp\min\{v,(Y/(1+E))^{1/2}\}$ and use
$F_Y^{p/2}\le F_U(t)^{p/2}F_U(Ut)^{p/2}$.
The two transformed grids have denominators $s$ and
$s/(s,U)$, both at least $v\ge U$; repeated points are counted with
their original multiplicity. This is a statement about the entire
grids, and does not require each numerator to be coprime to $s$.
Moreover, the displaced neighborhoods have width $O(1/U)$ in both
blocks, since $U^2(1+E)\ll Y$.
By \eqref{dec:variation}, the sum of the $p$th powers of the
centerwise suprema in either block is $O(sU^{g-1})$.
Apply Cauchy's inequality, allow for the $O(1+E)$ grid points in each
neighborhood, and use $w\le\tau^{1-p/2}F_Y^{p/2}$.
For $E\gg Y$, the pointwise bound suffices.

The second is the corresponding base-$b$ proof of Maynard's mixed
rational-grid lemma \cite[Lemma 10.7]{Maynard2019}.
Fix $q_1\asymp Q_1$ with $(q_1,b)=1$, and fix $d_0\asymp D$
whose prime factors all divide $b$. If $DE\ll Y$, then
\begin{align}
 &\sum_{\substack{q_2\asymp Q_2\\(q_2,b)=1}}
 \sum_{\substack{a<d_0q_1q_2\\(a,d_0q_1q_2)=1}}
 \sum_{\substack{|\eta|\le E/Y\\Y(\beta+a/(d_0q_1q_2)+\eta)\in\mathbb Z}}
 F_Y(\beta+a/(d_0q_1q_2)+\eta)\notag\\
 &\hspace{25mm}\ll (DE)^{\gamma_0}(Q_1Q_2^2)^{1/21}
       +E^{5/6}D^{3/2}Q_1Q_2^2Y^{-10/21}.
 \label{dec:maynardmixed}
\end{align}
Here $\beta$ is fixed throughout the sum. To explain the roles of the base
and the common shift, choose powers of $b$ with
$D'\asymp D,E'\asymp E,V\asymp(Y/(D'E'))^{1/2}$, and factor
$d_0=d_1d_2d_3$ so that $d_3=(d_0,D')$ and
$d_2d_3=(d_0,VD')$. Enumeration by the Chinese remainder theorem
and the digital product formula first contribute $E^{\gamma_0}$ from
the terminal block. The center-separation parameters for the two remaining
equal-sized blocks are $H_1=d_1d_2Q_1Q_2^2$ and $H_2=d_1Q_1Q_2^2$,
respectively. Parseval's identity and the squared version of
\eqref{dec:variation} give the full bound
\[
 E^{\gamma_0}\{d_3^{\gamma_0}(H_1^{1-\alpha}+H_1V^{-\alpha})
 +(d_2d_3)^{\gamma_0}(H_2^{1-\alpha}+H_2V^{-\alpha})\}.
\]
Now use $1-\alpha<1/21$, $\alpha/2>10/21$,
$\gamma_0+10/21<5/6$, and $1+10/21<3/2$ to obtain
\eqref{dec:maynardmixed}. In the two blocks the common shift becomes
$D'\beta$ and $D'V\beta$; it changes neither circular spacing nor integrals
over a full period. This proves the same estimate in base eleven and
makes explicit why the shift must not vary with $q_2$.

\subsection{Uniform exponent inequalities for all rank-two denominators}
\label{dec:ranktwo}

All exponents in this subsection are measured relative to the original
$q$. Let $R=q^\nu$ and $Q=q^{\vartheta+o(1)}$, where
$\vartheta=10\nu/3$, so that $R=Q^{3/10+o(1)}$.
The high-frequency weight $f_Q^*$ satisfies
\begin{equation}\label{dec:highweight}
 0\le f_Q^*\ll\min(F_Q,q^{-x}),\qquad
 f_Q^*\ll q^{(p-1)x}F_Q^p,\qquad
 W_Q=\sum f_Q^*\ll q^s.
\end{equation}
There is also a low-digit load $q^\ell$.
After division by $R^2$, the two terms in Lemma~\ref{dec:fejer}
have exponents
\begin{equation}\label{dec:generaltargets}
 \ell+s-x,\qquad \ell+\nu+2s-\vartheta.
\end{equation}
The zero Fourier term will be controlled throughout the parameter domain.
If the general term has exponent below $-10^{-6}$, the layer is already
settled; hence we retain the constraint $\ell+s-x\ge-10^{-6}$ below.

Fix a frequency pair $(A_1,A_2)$, and let $J=q^j\ge1$ be the actual
number of near-resonant multiplier pairs. We decompose this count
dyadically; factors bounded independently of $q$ are absorbed in the
fixed width margin below. In particular, $0\le j\le2\nu+o(1)$. Any three independent integer
points in the convex body
\[
 |r|,|s|\ll R,\qquad |rA_1-sA_2+nQ|\ll R
\]
have determinant $O(R^3/Q)=o(1)$, so the rank is at most two.
If the actual points are collinear, the primitive relation has height
$O(R/J)$. Taking the smaller of \eqref{dec:relation} and the trivial
mass bound $W_Q^2$, and optimizing in $J$, bounds the normalized
contribution of this part by
\begin{equation}\label{dec:linetargets}
 q^{\ell+(3s-x)/2-\nu+o(1)}+
 q^{\ell+2s-\nu-\vartheta/3+o(1)}.
\end{equation}
Here we use $\min(t,c/t)\le c^{1/2}$ and
$\min(t,c/t^2)\le c^{1/3}$, respectively.

If there are two independent actual points, the lattice-area argument
of \cite[Lemma 14.1]{Maynard2019} gives
\begin{equation}\label{dec:rationalwitness}
 1\le q_0\ll\mathcal R:=R^2/J,\quad (b_1,b_2,q_0)=1,
 \qquad A_i/Q=b_i/q_0+O(\mathcal R/(Qq_0)).
\end{equation}
Here is the calculation in the present normalization. Write
$e(r,s,n)=rA_1-sA_2+nQ$ and transform the section lattice by
$(r,s,n)\mapsto(r/R,s/R,e(r,s,n)/R)$.
Its two successive minima satisfy $\lambda_1\lambda_2\ll J^{-1}$,
by the planar lattice-point bound, since the section contains $J$
points and two independent ones. Choose independent integer vectors
$v,v'$ attaining these minima up to fixed constants, and put
$w=v\times v'$. Then
\[
 |w_3|\ll R^2/J,\qquad
 \left|w_1-w_3A_1/Q\right|+
 \left|w_2+w_3A_2/Q\right|\ll R^2/(JQ).
\]
The last two inequalities follow by substituting
$n=(e-rA_1+sA_2)/Q$ in the cross product.
If $w_3=0$, integrality and $R^2/(JQ)=o(1)$ would force $w=0$,
a contradiction. Divide $w$ by the gcd of its coordinates and choose
its sign so that $w_3>0$; then take
$(b_1,b_2,q_0)=(w_1,-w_2,w_3)$.
This proves \eqref{dec:rationalwitness}, including the case $J=O(1)$.
Let $\lambda=2\nu-j$, so that $0\le\lambda\le2\nu+o(1)$.
Put $(b_i,q_0)=g_i$, interchange the coordinates so that $g_1\le g_2$,
and write
\[
 g_1=d_1g'_1,\quad q_0=d_0d_1q'g'_1g_2,\qquad
 g_2=d_2g'_2,\qquad(q'g'_1g'_2,b)=1,
\]
where all prime factors of $d_0,d_1,d_2$ divide $b$.
The number of their possible values is bounded by a fixed power of $\log q$.
Let $(a_1,a_2,a_3,a_4,a_5,a_6)$ be the scale exponents of
$(q',g'_1,g_2,d_0,d_1,d_2)$, respectively. To keep the matrix notation
concise, we denote them below by $(a,c,h,D,E,G)$. Thus
\begin{equation}\label{dec:geometryconstraints}
 \begin{gathered}
 a,c,h,D,E,G\ge0,\qquad 0\le\lambda\le2\nu,\\
 \rho_0=a+c+h+D+E\le\lambda,\qquad c+E\le h,\qquad G\le h.
 \end{gathered}
\end{equation}
Put $w_0=10^{-6}$ and $e_0=\lambda-\rho_0+w_0$, and define
\[
 z_1=\lambda+c-h+w_0,\quad
 z_2=\lambda+a+h-c-E+w_0,\quad
 z_{\rm arc}=\lambda+a+c-h+w_0.
\]
This fixed width margin absorbs all geometric constants.

We now define the rational-grid sums whose exponents enter the linear systems.
For an integer denominator $D_1$ and a width $E_1\ge1$, define
\[
 \mathcal S(D_1,E_1)=
 \sum_{\substack{a<D_1\\(a,D_1)=1}}
 \ \sum_{\substack{A\bmod Q\\\|A/Q-a/D_1\|\le C E_1/Q}}
 f_Q^*(A),
\]
where $C$ is fixed. Every rational representation is counted.
Within the chosen dyadic scales and fixed smooth parts, put
\[
 I(q')=\sum_{g'_1}\mathcal S(d_0d_1q'g'_1,q^{e_0}),\qquad
 O(q')=\sum_{g_2}\mathcal S(d_0q'g_2,q^{e_0}),
\]
and write $\sup_{q'}I(q')\ll q^{s_1+o(1)}$,
$\sum_{q'}O(q')\ll q^{s_2+o(1)}$.
These denominators are exactly $q_0/g_2$ and $q_0/g_1$ from the
simultaneous approximation, respectively.
Estimates \eqref{dec:separated}, \eqref{dec:adaptiverow}, and
\eqref{dec:maynardmixed} give
\begin{align}
 s_1\le\min\{&\gamma_0z_1,z_1-x,gz_1+(p-1)x,
                \max(A_1,B_1),\max(U_1,V_1)\},\label{dec:s1}\\
 s_2\le\min\{&\max(\gamma_0z_2,z_2-(1-\gamma_0)\vartheta),z_2-x,
 \notag\\[-2mm]
 & (p-1)x+\max(gz_2,z_2-(1-g)\vartheta),
   \max(A_2,B_2),\max(U_2,V_2),s+\max(0,z_2-\vartheta)\},
 \label{dec:s2}
\end{align}
where all the linear forms are as follows:
\begin{align*}
 A_1&=\gamma_0(\lambda-a-c-h+w_0)+(a+2c)/21,\\
 B_1&=5(\lambda+w_0)/6+a/6+7c/6-5h/6+2(D+E)/3-10\vartheta/21,\\
 U_1&=D+E+a+2c+e_0-(1-p/2)x-(1-g)(a+c),\\
 V_1&=D+E+a+2c+e_0-(1-p/2)x-(1-g)(\vartheta-e_0)/2,\\
 A_2&=\gamma_0(\lambda-a-c-h-E+G+w_0)+2(a+h-G)/21,\\
 B_2&=5(\lambda+w_0)/6+7a/6-5c/6+7h/6+2D/3-5E/6-G/2-10\vartheta/21,\\
 U_2&=D+2a+2h-G+e_0-(1-p/2)x-(1-g)(a+h-G),\\
 V_2&=D+2a+2h-G+e_0-(1-p/2)x-(1-g)(\vartheta-e_0)/2.
\end{align*}
For example, the separation parameter for the inner grid is
$d_0d_1q'(g'_1)^2$; multiplication by the width gives exactly the
exponent $z_1$. The orientation constraint gives
$z_1\le\lambda+w_0\le2\nu+w_0<\vartheta$ throughout the windows
used below, so the inner grid has no saturation term. In the outer grid, combine $q'g'_2$, losing only divisor
multiplicity. Thus \eqref{dec:maynardmixed} is applied with smooth part
$d_0d_2$ and coprime part $q'g'_2$, giving $A_2,B_2$.
The last term in \eqref{dec:s2} includes all outer representations:
for a fixed source point, the reduced centers are separated by at least
$\gg[d_0(Q_1G_2)^2]^{-1}$, so the number of representations is at most
$q^{o(1)}(1+q^{z_2-\vartheta})$.

For each actual frequency pair in a $J$-layer, first select one witness
from the finite set in \eqref{dec:rationalwitness}, for example the
lexicographically least $(q_0,b_1,b_2)$, with the coordinate orientation
also fixed by a deterministic rule. Group the pairs by the scales of
this chosen witness. Thus each packet $\mathcal P$ consists of ordered
source pairs, each carrying its one selected witness.
Empty packets are discarded. Its mass is
$P=\sum_{(A_1,A_2)\in\mathcal P}f_Q^*(A_1)f_Q^*(A_2)$.
For each pair, the selected witness supplies one term in
$I(q')O(q')$. Positivity therefore gives
$P\le\sum_{q'}I(q')O(q')$, while the set interpretation gives
$P\le W_Q^2$. Consequently,
\begin{equation}\label{dec:packetmass}
 P\le W_Q^2,\qquad
 P\ll \bigl(\sup_{q'} I(q')\bigr)\sum_{q'}O(q'),\qquad
 t:=\log_qP\le\min(2s,s_1+s_2)+o(1).
\end{equation}
The contribution of this packet is at most $q^\ell JP$.
Since $J/R^2=q^{-\lambda}$, its normalized target exponent is
$\ell+t-\lambda$.
The rational-grid estimate enlarges to all representations only after
the packet has been selected. This order makes it compatible with
the independent bound $P\le W_Q^2$.

We may first delete entire source arcs whose contribution is already
controlled. Fix $d_0,d_1$ and the scales above, and take the union of
the neighborhoods of width $O(q^{e_0}/Q)$ over all inner denominators
$d_0d_1q'g'_1$. Their separation parameter is $d_0d_1(Q_1G_1)^2$,
so \eqref{dec:squarearc} bounds their squared mass by
$\ll q^{gz_{\rm arc}-(2-p)x+o(1)}$. If
\[
 \ell+gz_{\rm arc}-(2-p)x<-10^{-6},
\]
let $\mathcal A$ denote this union of source arcs and apply
Lemma~\ref{dec:fejer} to $f_Q^*1_{\mathcal A}$.
Its general term is bounded by the displayed squared mass; its
zero Fourier term is bounded using the full mass $W_Q$ and the
global second target in \eqref{dec:generaltargets}.
These bounds pay for the full output norm of this source subweight.

The number of choices of scales and smooth parts is bounded by a
fixed power of $\log q$. Delete the union of their favorable source
arcs by Lemma~\ref{dec:outputnorm}, retaining its complement as one
further subweight. Restrict the previously selected packets to pairs
in that complement, without reselecting witnesses. If the inner
coordinate of a surviving pair lay in a favorable arc, the pair
would already have been deleted. Thus every remaining packet
satisfies
\begin{equation}\label{dec:arcfilter}
 \ell+gz_{\rm arc}-(2-p)x\ge-10^{-6}.
\end{equation}

\subsection{The forward and reverse bilinear windows}

We now combine the carry reductions with the rank-two estimates.
Each contribution is bounded by a power of $q$ whose exponent is a
piecewise linear function of the scale parameters. The constraints above
retain both the actual amplitude information and every rational
representation. It remains to show that all the relevant exponents are
strictly negative throughout their continuous parameter regions.

\begin{lemma}[Forward window]\label{dec:forwardlemma}
Throughout \eqref{dec:forwardregion}, every nonnegative subweight
of an actual amplitude layer satisfies
\[
 \mathcal N_q(w;R)\ll R^2q^{-c}
\]
for a fixed $c>0$.
\end{lemma}
\begin{proof}
Combine \eqref{dec:generaltargets}--\eqref{dec:arcfilter} and
\eqref{dec:forwardregion}--\eqref{dec:highconstraints} into rational
linear systems. Split each of the five binary maxima into branches,
adding the constraint that the selected term dominates the other.
In \eqref{dec:s2}, the transitions for the basic moment, the $p$th
moment, and representation multiplicity all occur at $z_2=\vartheta$.
There are therefore $32$ consistent branches.
The rational certificates in the supplementary material give
\[
\begin{array}{c|c}
 \text{target}&\text{strict upper bound}\\ \hline
 \ell+\nu+2s-\vartheta\quad(\text{unfiltered})&-6061839073/100023000000\\
 \ell+(3s-x)/2-\nu&-8736163/63218000000\\
 \ell+2s-\nu-\vartheta/3&-4100926993/142240500000\\
 \ell+t-\lambda\quad(\text{all remaining rank-two branches})&-43400641/189654000000
\end{array}
\]
For each feasible branch, a nonnegative rational linear combination
proves the stated strict upper bound; the remaining branches have
Farkas certificates. The general term has already been controlled
where $\ell+s-x<-10^{-6}$. All deleted source arcs and collinear
contributions are covered by the preceding lemmas.
Apply Lemma~\ref{dec:outputnorm} to the low-amplitude subweights and
the deleted source arcs. Within each retained subweight, add the
nonnegative contributions of the selected source-pair packets.
Both decompositions have only a logarithmic-power number of parts,
so the fixed power saving is retained.
\end{proof}

\begin{lemma}[Reverse window]\label{dec:reverselemma}
Let $w$ be a nonnegative subweight of the amplitude layer of size
$q^{-u}$ and restrict the multipliers in $\mathcal N_q$ to units
modulo $b$. For either input pair in \eqref{dec:groups}, throughout
$0.249999\le\nu\le0.258$ and $\nu-0.07\le u\le\kappa$, we have
\[
 \mathcal N_q(w;R)\ll R^2q^{-c}
\]
for a fixed $c>0$.
\end{lemma}
\begin{proof}
Take $p=3/2$ in Section~\ref{dec:ranktwo}, add
\eqref{dec:newconstraints}--\eqref{dec:effectiveell}, and split into
the two branches of the maximum in \eqref{dec:deepchoice}.
Each input pair gives $64$ rank-two branches and $12$ basic targets.
The exact rational certificates for the two parameter pairs give
\[
\begin{array}{c|c|c}
 \text{target}&(g,\delta)=(.14387,.0053)&(g,\delta)=(.1435,.004)\\ \hline
 \ell+(3s-x)/2-\nu&-568609129/16755000000000&-2061939/6275000000\\
 \ell+t-\lambda&-110754136637/2010600000000000&-1645105237/3012000000000\\
 \ell_0+\nu+2s-\vartheta&-130413037/2500000000&-52781879/1000000000
\end{array}
\]
The second collinear target is also strictly negative.
The unfiltered depth constraints give
$d_*-(1-\vartheta)\le-.093999$ and $d_*-\nu\le-.203998$.
Thus $D_*<B,R$, and the second term of \eqref{dec:deepcost}
is controlled as well. The remaining contributions are combined
as in the forward window.
\end{proof}

All these computations use only integers and rationals.
The verifier reconstructs the variables, constraints, targets, and all
branches from the fixed parameters in the text, compares them entry by
entry with the certificate matrices, and then checks the nonnegative
dual or Farkas combinations. There are $36$ forward systems and $76$
systems for each reverse input pair, for a total of $188$.
The correspondence between the mathematical variables and the entries
of each certificate matrix is fixed in the supplementary specification.
Since there are finitely many dual certificates, the sums of their
coefficients are bounded.
Consequently, rounding the original scales and allowing $q^{o(1)}$
losses changes each target by only $o(1)$, which the strictly negative
margins absorb.

\subsection{The classical rank-two estimate for large amplitudes and the right-hand window}
\label{dec:oldranktwo}

To give the full extension to base eleven, we record the precise form
of Maynard's rank-two proposition that is needed.
Let $\mathcal E=\{a:F_q(a/q)\ge q^{-\kappa}\}$, and let
$\mathcal F(Q_{
 \rm arc},E)$ be the rational-arc grid with denominators
$\asymp Q_{\rm arc}\le\sqrt q$ and offsets $\asymp E/q$,
where $E\le100\sqrt q/Q_{\rm arc}$; the case $E=0$ is allowed.
Let $\mathcal B_2(N,K,\Delta)$ be the set of frequency pairs for which
there is a rank-two integer lattice containing at least $\Delta KN^2$
points satisfying
\[
 \|\mathbf n\|_2\ll N,\qquad
 |n_1a_1+n_2a_2+n_3q|\ll\Delta q,
\]
and these points are not all collinear.

\begin{lemma}[A common form of the classical rank-two input]\label{dec:classicalrank}
If $N,K\ge1$, $\Delta\ge N/q$, and $NK\ge q^{17/40}$,
then for some fixed $c>0$,
\begin{equation}\label{dec:classicalrankbound}
 \sum_{\substack{(a_1,a_2)\in\mathcal B_2(N,K,\Delta)\\
 a_1,a_2\in\mathcal E\cap\mathcal F(Q_{\rm arc},E)}}
 F_q(a_1/q)F_q(a_2/q)
 \ll\frac{q(\log q)^{O(1)}}{NK(Q_{\rm arc}+E)^c}.
\end{equation}
This holds for both $b=10$ and $b=11$.
\end{lemma}
\begin{proof}
For base ten this is \cite[Proposition 13.3]{Maynard2019}.
We spell out the base-dependent part of its proof.
The lattice and gcd decomposition is
\cite[Lemmas 14.1--14.3]{Maynard2019}: it uses integer determinants,
reduced fractions, and the Chinese remainder theorem.
For a different fixed base, its smooth denominators are simply those
supported on the primes dividing $b$.
The estimates below verify all the digital and exponent inputs of
\cite[Lemma 14.4]{Maynard2019}. First, \eqref{dec:momentold} gives
\[
 \sum_{a\in\mathcal E}F_q(a/q)
 \ll q^{g_0+(p_0-1)\kappa}=q^\sigma,\qquad
 \sigma=1533559/5334560<23/80.
\]
The lattice-area argument gives a common rational denominator and width.
Use the gcd orientation of Section~\ref{dec:ranktwo}, so that
$G_1\ll G_2$. The denominator and error scales obey
\[
 Q_0=Q_1G_1G_2D_0D_1,\qquad E_0\ge1,\qquad
 Q_0E_0\ll q/(NK)\le q^{23/40}.
\]
The decomposition bounds the source-pair mass, up to logarithmic
factors, by $\min(S_1S_2,S_1S_3)$.
Here $S_1$ is the supremum of the inner grid sum.
The outer projection of the selected source pairs has two alternative
upper bounds: $S_2$ counts all rational representations, while $S_3$
counts a denominator parameter only once for each source frequency
and uses the total mass on $\mathcal E$.
The latter need not bound the enlarged sum $S_2$; both bound the
same selected source-pair contribution before enlargement.
The common-shift estimates of Section~\ref{dec:grids} give
\begin{align*}
 S_1&\ll(Q_0E_0)^{\gamma_0},\\
 S_2&\ll\left(\frac{Q_0^2E_0}{D_0D_1^2G_1^2}\right)^{\gamma_0}
       +\frac{Q_0^2E_0}{q^{50/77}D_0D_1G_1},\\
 S_3&\ll q^{23/80}\left(1+\frac{Q_0}{D_0D_1NK}\right),\\
 S_1&\ll Q_0^{1/21}(D_0D_1E_0)^{\gamma_0}
       +\frac{Q_0G_1(D_0D_1)^{1/2}E_0^{5/6}}{G_2q^{10/21}}.
\end{align*}
The bound for $S_3$ follows from the difference of the actual rational
approximations and the integer divisor bound.
For a fixed reduced center, the cross-difference associated with the
other center is an integer of size $O(Q_0/(D_0D_1NK))$.
Once this integer is fixed, the Chinese remainder theorem and divisor
bounds control all representations.

If the first term in the last estimate dominates, multiplication
by the bound for $S_2$ gives
\[
 S_1S_2\ll
 Q_0^{54/77+1/21}E_0^{54/77}
 +Q_0^{2+1/21}E_0^2q^{-50/77}
 \ll(Q_0E_0)^{1-c}.
\]
For the last step, use $Q_0E_0\ll q^{23/40}$ and
$\frac{22}{21}\frac{23}{40}<\frac{50}{77}$.
If the second term dominates, use $G_1\ll G_2$ and $E_0\ge1$ to obtain
\[
 S_1\ll\min\left((Q_0E_0)^{\gamma_0},
          \frac{Q_0E_0(D_0D_1)^{1/2}}{q^{10/21}}\right).
\]
For the first term of $S_2$, interpolate the two displayed
bounds for $S_1$ with weights $1/3,2/3$; for its second term,
use the second bound directly. This gives
\[
 S_1S_2\ll Q_0^{3/2}E_0^{6/5}q^{-3/10}
                +Q_0^3E_0^2q^{-9/8}.
\]
The rounded exponents are upper bounds: in particular,
$50/77+10/21>9/8$.
Since $Q_0E_0\ll q^{23/40}$, the first term is
$\ll Q_0E_0q^{-1/80}$ and hence $\ll(Q_0E_0)^{1-c}$.

Likewise, the two terms of $S_3$ give
\[
 S_1S_3\ll q^{23/80}(Q_0E_0)^{27/77}
             +\frac{Q_0^2E_0}{q^{3/16}NK}.
\]
Here $10/21-23/80>3/16$.
The second term is at most
$Q_0q^{13/16}/(NK)^2\ll Q_0q^{-3/80}$, so it too is
$\ll(Q_0E_0)^{1-c}$. We have therefore proved
\begin{align*}
 S_1S_2&\ll(Q_0E_0)^{1-c}+Q_0^3E_0^2q^{-9/8},\\
 S_1S_3&\ll(Q_0E_0)^{1-c}+q^{23/80}(Q_0E_0)^{27/77}.
\end{align*}
Taking the smaller of the last two terms and interpolating with weights
$23/100,77/100$ gives
\[
 \min\{Q_0^3E_0^2q^{-9/8}, q^{23/80}(Q_0E_0)^{27/77}\}
 \le Q_0^{96/100}E_0^{73/100}q^{-299/8000}.
\]
Thus every packet is bounded by $(Q_0E_0)^{1-c}$.
Separation of the common rational approximation from the original arc,
as in \cite[Lemma 14.2]{Maynard2019}, gives
$Q_{\rm arc}+E\ll(q/(NK))^2$.
Writing $T=q/(NK)$, the packet bound is $O(T^{1-c})$,
which is $O(T(Q_{\rm arc}+E)^{-c/2})$.
After decreasing $c$, this proves \eqref{dec:classicalrankbound}.
For fixed base, the number of smooth denominators is bounded by a
fixed power of a logarithm. The argument uses only
\eqref{dec:moment1}, \eqref{dec:momentold}, \eqref{dec:maynardmixed},
and comparisons of rational exponents, so it applies in base eleven
as well.
\end{proof}

Take $\zeta=.001$, and put
\[
 u_0(\nu)=\nu-\tfrac32(1-\alpha)-\zeta.
\]
On an amplitude layer with $.249999\le\nu\le.29$ and
$u\le u_0(\nu)$, the second moment gives $W\ll q^{1-\alpha+u}$.
Decompose the full Cauchy kernel according to thickness and multiplicity,
write the number of points as at least $\Delta K'R^2$, where
$\Delta\ge R/q$, and put $\mathcal R=q/(RK')$.
Fix $0<\eta<\zeta/2$. If $\mathcal R\ge q^\eta W^2$, the total
mass directly gives the bound $q^{-\eta}\mathcal R$.
Otherwise,
\begin{equation}\label{dec:largeamplitudeNK}
 RK'\gg q^{1-2\nu+(1-\alpha)+2\zeta-\eta-o(1)}.
\end{equation}
Since $1-\alpha>.01$, this exponent is still greater than $.4315$
at $\nu=.29$. Thus $RK'>q^{17/40}$, and $K'$ is bounded below
by a fixed power of $q$. This lower bound absorbs the
$O(\Delta R^2)$ points outside the plane in the lattice-point
classification of \cite[Lemma 13.2]{Maynard2019}.
Apply Lemma~\ref{dec:classicalrank} to the rank-two part.
For the collinear part, the actual nonzero multiples give height
$V\ll\mathcal R/R$. By \eqref{dec:relation}, its mass divided by
$\mathcal R$ is at most
\[
 q^{o(1)}\left(q^\eta\frac{q^{-u}W^3}{R^2}
          +q^{2\eta}\frac{W^6}{qR^3}\right).
\]
The two exponents are at most $\eta-2\zeta$ and
$2\eta+3\nu-1-3(1-\alpha)-6\zeta$, respectively, both strictly negative.
Since $1-\alpha<.046$, we have $u_0(\nu)>\nu-.07$.
This range overlaps the intermediate-amplitude bands in the two
preceding windows.

The right-hand window requires only the earlier moment estimate. Put
\[
 \nu_0=(3\sigma-\kappa)/2=3066991/10669120.
\]
On each actual amplitude layer retain
$W\ll q^{g_0+(p_0-1)u}$. Since $3p_0-4>0$, this gives
$q^{-u}W^3\ll q^{3\sigma-\kappa}$.
For $\nu_0+\varepsilon_1\le\nu\le.425-\varepsilon_1$, choose
$0<\eta<\min(\varepsilon_1,\kappa-\sigma)$.
The nontrivial multiplicity layers satisfy
$RK'>q^{1-2\sigma-\eta-o(1)}>q^{.425}$, while the two collinear
exponents are at most
\[
 \eta+2\nu_0-2\nu<0,\qquad
 2\eta+6\sigma-1-3\nu<0.
\]
At $\nu=\nu_0$ and $\eta=0$, the second expression is
$-293477/2133824$. Taking $\varepsilon_1=.0004$ and using the fixed
joining points $.2575,.288$, the windows together cover
\begin{equation}\label{dec:rawwindow}
 [.249999,.4244].
\end{equation}
The following subdivision records which estimates supply this range:
\begin{center}
\begin{tabular}{ll}
\toprule
Short-factor exponent $\nu$ & Estimates used\\
\midrule
$[.249999,.2575]$ & Reverse carries and the large-amplitude bound\\
$[.2575,.288]$ & Forward carries and the large-amplitude bound\\
$[.288,.4244]$ & The classical right-hand window\\
\bottomrule
\end{tabular}
\end{center}
In the first two rows, $u_0(\nu)>\nu-.07$ ensures that the
large-amplitude and carry estimates overlap. The ordinary spectrum
$F_q<q^{-\kappa}$ is treated separately in the completion below.
All lattice-point classifications retain the actual nonzero multiplier
points. Absorbing the fixed geometric constants into
$N\mapsto CR$, $\Delta\mapsto C\Delta$, and
$K'\mapsto cK'/C^3$ preserves the lower bound for the thickness
and the fixed power margin in $NK'$.
Moreover, a nonempty layer satisfies $\Delta K'\ll1$, so the
power lower bound for $K'$ ensures $\Delta<1$ after rescaling.

\subsection{Completion of the minor-arc estimate}
\label{dec:typeii}

\begin{proof}[Proof of Proposition~\ref{dec:minorarc}]
For the ordinary spectrum $F_q<q^{-\kappa}$, we retain the earlier
moment estimate, since
\begin{equation}\label{dec:ordinary}
 \sum_{F_q<q^{-\kappa}}F_q^2
 \ll q^{-\delta_0},\qquad
 \delta_0=\kappa(2-p_0)-g_0=127/5334560>0.
\end{equation}
The product coefficient of any bilinear convolution with bounded
coefficients is at most $\tau_2(n)$. In the physical support of a fixed
prefix, each residue class modulo $q$ contains only $O_v(1)$ integers.
Parseval's identity therefore bounds the squared frequency sum by
$\ll_vq^2\log^3q$. The ordinary spectrum consequently contributes
$O_v(Aq^{-\delta_0/2}\log^{3/2}q)$.

Choose a Dirichlet approximation with denominator at most $\sqrt q$
for every frequency, and group the frequencies by denominator
$Q_{\rm arc}$ and offset $E$. The number of packets is bounded by
a fixed power of $\log q$. If the logarithmic exponent $D_0$ in the definition of
$\mathfrak M$ is sufficiently large, every frequency outside
$\mathfrak M$ has $Q_{\rm arc}+E\gg(\log q)^{D_0}$.
The constants implicit in the packet widths are absorbed by increasing
$D_0$.
The large-amplitude and right-hand windows retain the negative power
in \eqref{dec:classicalrankbound}, while the intermediate-amplitude
range has the stronger fixed power saving.
The major arcs may subsequently be enlarged, since an arbitrary
frequency mask is already allowed in $\gamma_a$.

Lemma~\ref{dec:outputnorm}, applied to $\mathcal K_q$, combines
all frequency and amplitude packets with only a logarithmic loss.
This also handles every cross term between different packets;
no disjointness of their output supports is needed.

For a fixed admissible prefix $v$, separate the product restriction
$vq\le mr<(v+1)q$ by a truncated Perron formula with half-integer
endpoints. Take real part $1/\log q$ and height $q^4$.
The truncation error for each product is $O_v(q^{-3})$, and the
integral kernel has total absolute mass $O(\log q)$.
The integral merely multiplies the arbitrary coefficients on the two
sides by $m^{it},r^{it}$ and bounded real-part factors.
The factor $1_{(mr,b)=1}$ is absorbed separately into the two coefficient
sequences. Hence \eqref{dec:minorarcII} holds throughout the range
\eqref{dec:rawwindow}. The argument permits a further restriction of
the frequency mask, so the major arcs may be enlarged as needed in
Section~\ref{dec:prime-section}.
\end{proof}


\section{A common comparison model and the prime asymptotic}
\label{dec:prime-section}

Assume the digital hypotheses of Theorem~\ref{dec:digital-criterion}.
We now pass from Proposition~\ref{dec:minorarc} to a comparison of prime
counts. The comparison model retains the initial digits, which control
physical intervals, and the terminal digits, which control local
congruences. Averaging only the middle block makes the model nonnegative
and permits the same sequence to satisfy Type I and Type II estimates.

\subsection{The model and the common analytic estimates}
\label{dec:model}

Fix an admissible positive prefix $v$, and let $x=(v+1)q$.
Choose a sufficiently large fixed $C$ according to the final error
precision required, and put
\[
 L=\lceil C\log_b\log q\rceil,\quad B_0=b^L,\quad
 Y=q/B_0^2,\quad \rho=(b-1)^{k-2L}/Y.
\]
For sufficiently large $k$, we have $2L<k$.
For $n=vq+a(q/B_0)+B_0m+c$, with $0\le a,c<B_0$ and
$0\le m<Y$, define
\begin{equation}\label{dec:modeldef}
 f_*(n)=1_{n-vq\in\mathcal D_k}1_{(n,b)=1},\qquad
 G_C(n)=\rho\,1_{a\in\mathcal D_L}1_{c\in\mathcal D_L^*},\qquad
 W_C=f_*-G_C,
\end{equation}
and set these functions equal to zero outside this support.
Here $\mathcal D_L^*=\mathcal D_L\cap\{n:(n,b)=1\}$.
If $h=\#\{j:0\le j<b, j\ne d, (j,b)=1\}$, then
\[
 \sum_n f_*(n)=h(b-1)^{k-1}=\sum_nG_C(n).
\]
The masses also agree exactly on each cylinder obtained by fixing
the initial $L$ suffix digits. Thus the model preserves the actual
initial-digit structure at that scale.

Put $\Gamma_0=50/77$ and
\[
 \Gamma_*=\frac{\Gamma_0\alpha}{\alpha+\Gamma_0(1-\alpha)}>.629.
\]

\begin{proposition}[A common comparison sequence]\label{dec:commonmodel}
Fix an admissible positive prefix $v$, an exponent $\Gamma<\Gamma_*$,
and $D,K>0$. There is $C_0$ such that, for every fixed $C\ge C_0$,
the same sequence $W_C=f_*-G_C$ satisfies
\begin{equation}\label{dec:weightedI}
 \sum_{r\le x^\Gamma}\tau(r)^D\max_{J\ {\rm interval}}
 \left|\sum_{m\in J}W_C(rm)\right|
 \ll_{v,D,K}A\log^{-K}q
\end{equation}
and
\begin{equation}\label{dec:fullII}
 \left|\sum_{m\in I_M,r\in I_R}\alpha_m\beta_rW_C(mr)\right|
 \ll_{v,D,K}A\log^{-K}q
\end{equation}
whenever $MR\asymp_vq$, $q^{.249999}\le R\le q^{.4244}$,
$I_M\subset[M,2M]$, $I_R\subset[R,2R]$ are intervals, and
$|\alpha_m|\le\tau(m)^D$, $|\beta_r|\le\tau(r)^D$.
The latter estimate also allows a common interval restriction on $mr$.
The constants are uniform in the truncations and coefficient
sequences for fixed $v,\Gamma,D,K,C$.
\end{proposition}

The Type I estimate is proved directly from the first Fourier moment.
For Type II, the minor arcs have already been treated; the remaining
step is to compare the full major-arc projection with $G_C$.
We then extend the coefficient class using a divisor-moment bound.

\subsection{Maximal Type I with divisor weights}

\begin{lemma}[Maximal weighted Type I]\label{dec:typeI}
Given $\Gamma<\Gamma_*$ and $D,K>0$, if $C$ is sufficiently
large, then the maximal estimate \eqref{dec:weightedI} holds.
\end{lemma}
\begin{proof}
The lower bound for $\Gamma_*$ follows from $\alpha>477/500$ in
both bases and
$\Gamma_*>23850/37879>629/1000$.
For a free suffix of scale $b^s$ and $(r,b)=1$, put
\[
 U_s(r)=\frac1r\sum_{a=1}^{r-1}
   |\widehat{1_{\mathcal D_s}}(a/r)|.
\]
Additive orthogonality shows that the same nonnegative quantity
$U_s(r)$ controls the error relative to $(b-1)^s/r$ for every prefix
shift and every target residue class.
Reduce the nonzero frequencies and apply the first-moment estimate
on rational grids to obtain
\[
 \frac1R\sum_{t\asymp R}\sum_{(a,t)=1}F_{b^s}(a/t)
 \ll R^{-23/77}+R/(b^s)^{50/77}.
\]
For a reduced fraction $a/t$ with $(t,b)=1$ and $1<t<Z^{1/3}$,
the digital product gives
\[
 F_Z(a/t+\eta)\ll\exp(-c\log Z/\log t),\quad
 t<Z^{1/3},\quad(t,b)=1,\quad t>1,\quad |\eta|<Z^{-2/3}/2.
\]
This is the base-$b$ proof of \cite[Lemmas 8.1--8.2]{Maynard2019}:
in every block of $O_b(\log t)$ digits, the multiplicative orbit
$b^ja/t$ has at least one point at a fixed positive distance from the
integers. Otherwise, the recurrence between successive nearby integers
would force $t$ to divide a power of $b$.
The corresponding one-digit factor is strictly below $1$.
A small perturbation preserves this conclusion for an initial fixed
proportion of the digits. Thus, whenever the upper modulus bound lies
below $(b^s)^{\Gamma_0}$ by a fixed power, the sum $\sum_rU_s(r)$
has any prescribed logarithmic saving.
With a unit terminal digit, the corresponding progression error is
bounded by $U_s^*(r):=h_{b,d}U_{s-1}(r)$ for $s\ge1$, by fixing
each allowed unit digit.

Choose
\[
 \Gamma(1-\alpha)/\alpha<\Delta<1-\Gamma/\Gamma_0,\qquad
 Z\asymp q^{1-\Delta}
\]
with $Z$ a power of $b$. Any initial truncation decomposes into at most
$b-1$ complete cylinders of each free suffix length and one incomplete
boundary cylinder of length $Z$.
Even if the truncation endpoint is chosen arbitrarily for each $r$,
the complete-cylinder errors are bounded by
$(b-1)\sum_{b^s\ge Z}U_s^*(r)$.
If $r\le Z$, choose $b^j\le r<b^{j+1}$ and fix the digits above
position $j$. A
residue class modulo $r$ then has at most one low-digit candidate.
Thus each class contains at most $(b-1)Z^\alpha r^{-\alpha}$
points in the boundary cylinder.
Summing over $r\le x^\Gamma$, the boundary error is at most
$Z^\alpha x^{\Gamma(1-\alpha)}\ll Aq^{-c}$.

Divisor weights do not change this power range.
When reducing $r=tj$, use
$\tau(tj)^D\le\tau(t)^D\tau(j)^D$.
For large denominators, the bound $\tau(t)^D\ll_\epsilon t^\epsilon$
is absorbed by the fixed power margin; for small denominators, use
the exponential decay above.
At the boundary, use
$\sum_{r\le R}\tau(r)^D\ll_DR\log^{2^{\lceil D\rceil}-1}(2R)$.

Let $F_*(t)=\sum_{n\le t}f_*(n)$ and
$G_0(t)=\sum_{n\le t}G_C(n)$.
The masses of the two functions agree on every complete cylinder with
its initial $L$ suffix digits fixed, so
$\sup_t|F_*(t)-G_0(t)|\ll A(b-1)^{-L}$.
For each initial interval, compare the actual progression count with
$F_*(t)/r$ and the model progression count with $G_0(t)/r$.
Their different reference masses therefore cost at most
$A(b-1)^{-L}/r$; arbitrary intervals are differences of two
initial intervals.
Once the initial and terminal words are fixed, the condition $r\mid n$
selects a unique residue class of the middle variable $m$.
The error for the model, for each modulus and arbitrary truncation,
is therefore at most
$\rho\#\mathcal D_L\#\mathcal D_L^*\ll AB_0^2/q$.
Subtracting the model from the actual function gives
\begin{align*}
 \sum_{r\le x^\Gamma}\tau(r)^D\max_J
 \left|\sum_{m\in J}W_C(rm)\right|
 \ll{}&A\log^{-K-2}q
   +Aq^{\Gamma-1}B_0^2\log^{T_D}q\\
 &+A(b-1)^{-L}\log^{T_D}q,\qquad T_D=2^{\lceil D\rceil}.
\end{align*}
Taking $C\alpha>K+T_D+3$ proves the result.
The rows with $(r,b)>1$ vanish identically.
\end{proof}

\subsection{The full major arcs and the physical model}

\begin{lemma}[Bilinear comparison of the full major arcs with the model]\label{dec:bridge}
Given $P>0$ and $\theta_0>0$, one can choose a fixed $C$ and then
the major-arc width so that, for $MR\asymp_vq$,
$\min(M,R)\ge q^{\theta_0}$, and arbitrary bounded coefficients,
\begin{equation}\label{dec:bridgeeq}
 \sum_{m,r}\alpha_m\beta_r1_{vq\le mr<(v+1)q}1_{(mr,b)=1}
 \{P_{\mathfrak M}f_{\rm all}(mr)-G_{\rm all}(mr)\}
 \ll_{v,P}A\log^{-P}q.
\end{equation}
Here $f_{\rm all}=1_{\mathcal D_k}$ is viewed as a periodic function,
and $G_{\rm all}$ is the periodic model obtained from
\eqref{dec:modeldef} by removing the unit terminal-digit restriction.
\end{lemma}
\begin{proof}
Take the major-arc denominator bound to be
$R_0=\exp(\eta\sqrt{\log q})$, with $\eta>0$ fixed and sufficiently
small, and the offset bound to be a sufficiently large fixed
logarithmic power $H$.
For a denominator containing a prime factor that does not divide $b$,
first discard $O_b(\log r)$ digit factors to remove its smooth part.
The remaining denominator is coprime to $b$ and greater than $1$,
so the orbit-decay estimate above still applies.
Choosing $\eta$ sufficiently small allows for all $O(R_0^2H)$
major-arc grid points. The total Fourier pairing over these
frequencies is $O(Ae^{-c\sqrt{\log q}})$.

Absorb the unit restriction into the coefficients and write
\[
 \mathcal B(\xi)=\sum_{m\in I_M,\,n\in I_R}\alpha_m\beta_n
                 1_{vq\le mn<(v+1)q}e(-mn\xi).
\]
Write $\mathcal B_0$ for the same sum without the product cutoff.
For a reduced denominator $r$ supported on primes dividing $b$ and
an integer offset $z$, the bilinear Cauchy estimate gives
\begin{equation}\label{dec:classicalbilinear}
 |\mathcal B(c/r+z/q)|
 \ll\frac{q\log^2q}{\sqrt{r(1+|z|)}}+q^{1-\theta_0/2+o(1)}.
\end{equation}
For a reduced fraction $a/s$ with $|\xi-a/s|\le s^{-2}$, Cauchy's inequality,
geometric sums, and grouping by separation give
$|\mathcal B_0(\xi)|^2\ll MR(MR/s+M+R+s)\log(2s)$.
For $z=0$, take $s=r$.
For $z\ne0$, apply Dirichlet approximation again with denominator
bound $\lfloor2q/(r|z|)\rfloor$.
The old fraction $c/r$ is excluded, and the lower bound $1/(rs)$
for the difference of two distinct fractions forces
$s\asymp q/(r|z|)$.
Substitution into the preceding estimate, followed by Perron's
formula to retain the product truncation, proves
\eqref{dec:classicalbilinear}.

Choose radix powers $U\asymp r$, $V\asymp Z$ with $UV\le q$.
Keeping the two end blocks bounds $F_q(c/r+z/q)$ by
$F_U(c/r+z/q)F_V(z/V)$, because $r$ divides $q/V$ for large $q$.
Variation bounds the sum over $c$ of the first factor's local suprema
by $O(r^{\gamma_0})$. The shifted first moment sums the second factor
over $z$ at cost $O(Z^{\gamma_0})$. Thus for each dyadic $Z\ge1$,
\begin{equation}\label{dec:packlone}
 \sum_{(c,r)=1}\sum_{|z|\asymp Z}
 |\widehat f_{\rm all}(c/r+z/q)|\ll A(rZ)^{\gamma_0}.
\end{equation}
The case $Z=1$ also includes $z=0$.
The chosen digit blocks are disjoint because $r,Z=q^{o(1)}$.

To control the physical boundaries of the model, let $Q_0=q/B_0$
and $T=\lfloor\log^{9C/8}q\rfloor$, and use the probability kernel
on $\mathbb Z/Q_0\mathbb Z$ given by
\[
 K_T(u)=\frac1{Q_0T}\left|\sum_{j<T}e_{Q_0}(ju)\right|^2,\qquad
 G_T(n)=\sum_uK_T(u)G_{\rm all}(n-B_0u).
\]
These translations preserve the last $L$ digits.
The Fourier support is $a=cq/B_0+z$ with $|z|<T$, and the
Fourier coefficients are multiplied by $(1-|z|/T)_+$.
The Fej\'er tail bound is
$\sum_{|u|\ge h}K_T(u)\ll\min(1,Q_0/(Th))$.
Thus the smoothing error at physical distance $D$ from an initial-word
boundary is at most $O(\rho\min(1,q/(TD)))$.
The absolute coefficient of any multiplicative test is at most
$\tau_2(n)$, and the Dirichlet hyperbola decomposition gives
\[
 \sum_{n\in J}\tau_2(n)\ll_v |J|\log q+\sqrt q
\]
for every interval $J$ in the range $n\asymp_vq$.
Decomposing by distance from the at most $B_0$ initial-word boundaries
gives
\begin{equation}\label{dec:smoothing}
 \sum_n\tau_2(n)|G_{\rm all}(n)-G_T(n)|
 \ll A B_0^{3-2\alpha}T^{-1}\log^2q
      +Aq^{-1/2}B_0^{3-2\alpha}.
\end{equation}

The function $G_T$ also satisfies \eqref{dec:packlone}. Its support
in the packet requires $r\mid B_0$, since distinct rational centres
of denominators $r,B_0$ are separated by at least $1/(rB_0)$, greater
than the sum of their polylogarithmic widths divided by $q$. For
$r\mid B_0$, the normalized magnitude of $G_{\rm all}$'s
Fourier transform is
\[
 F_{B_0}(c/r+z/q)F_{B_0}(z/B_0)
       \frac{|D_Y(z/(B_0Y))|}{Y},\qquad
 D_Y(t)=\sum_{m<Y}e(mt).
\]
The last factor is at most $\min(1,O(B_0/(1+|z|)))$, and the Fej\'er multiplier defining $G_T$ has modulus at most one.
For $Z\le B_0$, use the first moment.
For $Z>B_0$, periodic repetition of the second factor gives
$O(ZB_0^{\gamma_0-1})$; multiplication by $O(B_0/Z)$ still
leaves at most $O(Z^{\gamma_0})$.

Put $S=B_0^{1/10}$.
For $b=10$ or $11$, every smooth denominator $r\le S$ divides $B_0$.
When $r,|z|\le S$, the actual digital measure and the model differ only
in the middle block. The normalized measures both have mass one, and the phase
variation across that block is $O(|z|/B_0)$. Hence
\[
 |\widehat f_{\rm all}(c/r+z/q)-\widehat G_T(c/r+z/q)|
 \ll A|z|(B_0^{-1}+T^{-1}).
\]
The total pairing error from the small packets is at most
\[
 O(AS^4(B_0^{-1}+T^{-1})).
\]
For the remaining packets, multiply \eqref{dec:classicalbilinear}
by \eqref{dec:packlone}. The main term in each packet is
\[
 A\log^2q(rZ)^{-(1/2-\gamma_0)}.
\]
The sums over smooth denominators and dyadic offset scales converge.
The total from the second term is $Aq^{-\theta_0/2+o(1)}$.

The logarithmic exponents in the three error terms without a power
saving in $q$ are, respectively,
\[
 -C(2\alpha-15/8)+2,\qquad -3C/5,\qquad
 -C(23/154)/10+2+o(1).
\]
It is enough to take $C\ge100(P+10)$, and then to take the
major-arc offset exponent at least $9C/8+1$, increasing it further
to obtain the precision required in \eqref{dec:minorarcII}.
Once $C$ is fixed, all these logarithmic scales are eventually below
$R_0$. Adding \eqref{dec:smoothing} proves the result.
The physical unit conditions are absorbed separately into the two
coefficient sequences; the two periodic functions being compared
remain precisely those in \eqref{dec:bridgeeq}.
The argument works for every fixed $C$ above the stated threshold,
with the major-arc widths chosen afterwards. This quantifier permits
the same $C$ to be used in Type I and Type II.
\end{proof}

Combining \eqref{dec:minorarcII} with Lemma~\ref{dec:bridge} gives
an estimate for the same $W_C$:
given $v,T$, after a single choice of sufficiently large $C=C(v,T)$,
for arbitrary bounded coefficients, $MR\asymp_vq$, and
$q^{.249999}\le R\le q^{.4244}$, we obtain \eqref{dec:fullII}
with $K=T$.
An additional common product truncation is also permitted.

\subsection{Extension to divisor-bounded coefficients}

First, for fixed $J>0$, we record the growth bound
\begin{equation}\label{dec:divisormoment}
 \sum_n(f_*(n)+G_C(n))\tau(n)^J
 \ll_{v,J,C}A\log^{C_J}q,\qquad C_J=2^{\lceil5J\rceil}+1.
\end{equation}
A direct proof is as follows.
Remove the prime factors greater than $x^{1/4}$.
There are at most four, counted with multiplicity, and their
contribution to the divisor function is at most $16$.
The remaining prime factors can be grouped greedily into at most five
blocks, each with product at most $\sqrt x$. Thus
\[
 \tau(n)^J\le16^J\sum_{r\mid n, r\le\sqrt x}\tau(r)^{5J}.
\]
For $f_*$, use the weighted arithmetic-progression bound from the
proof of Lemma~\ref{dec:typeI}.
For $G_C$, use the bound $O(A/r+AB_0^2/q)$ for each modulus
from the Chinese remainder theorem.
This proves \eqref{dec:divisormoment}.
Its logarithmic exponent $C_J$ is independent of the model parameter $C$.

Suppose that $|\alpha_m|,|\beta_r|\le\tau(m)^D,\tau(r)^D$.
For a target precision $K$, take
\[
 J=3D+1,\quad L_*=C_J+K+3,\quad K_*=K+2L_*+3.
\]
Truncate both coefficient sequences at $\log^{L_*}q$.
For the small coefficients, use \eqref{dec:fullII} with precision $K_*$.
For the large coefficients, sum by the actual product $t=mr$.
There are at most $\tau(t)$ factorizations, the product of the two
weights is at most $\tau(t)^{2D}$, and the exceptional indicator is
at most $\tau(t)^D/\log^{L_*}q$.
By \eqref{dec:divisormoment}, the total is $O(A\log^{-K-3}q)$.
There are only $O_v(\log q)$ nonempty dyadic boxes.
Thus \eqref{dec:fullII} extends to coefficients bounded by any fixed
power of the divisor function. We can now conclude the common-model proposition.

\begin{proof}[Proof of Proposition~\ref{dec:commonmodel}]
Choose $J,L_*,K_*$ as above. The Type I estimate requires
$C\alpha>K+2^{\lceil D\rceil}+3$.
The bounded-coefficient Type II estimate with precision $K_*$
requires only a further lower bound for $C$.
Its extension to divisor-bounded coefficients uses a logarithmic
moment exponent $C_J$ independent of $C$, so the two choices are
compatible. Choose the common major arcs after fixing $C$.
\end{proof}

\subsection{The asymptotic sieve}
\label{dec:sieve}

We now state precisely the asymptotic sieve theorem used
\cite[Theorem 2.2(a)]{FordMaynard2024}.
Let $w_n$ be supported on $(x/2,x]$, and suppose that the parameters
$(\gamma,\theta,\nu)$ satisfy
\[
 \tfrac12\le\gamma\le1-\theta-\nu\quad\text{or}\quad
 1-\theta\le\gamma<1,\qquad
 0\le\theta<\theta+\nu\le\tfrac12.
\]
For a sufficiently large fixed $B_*$, assume
\begin{align}
 &\sum_{m\le x^\gamma}\tau(m)^{B_*}\max_I
 \left|\sum_{\substack{x/2<mn\le x\\n\in I}}w_{mn}\right|
 \le x\log^{-B_*}x,\label{dec:FM_I}\\
 &\left|\sum_{\substack{(x/2)^\theta<m\le x^{\theta+\nu}\\x/2<mn\le x}}
 \xi_m\kappa_nw_{mn}\right|\le x\log^{-B_*}x,\quad
 |\xi_m|\le\tau(m)^{B_*}, |\kappa_n|\le\tau(n)^{B_*},\label{dec:FM_II}\\
 &\sum_n|w_n|\tau(n)\le x\log^\varpi x,\qquad
 w_n\ge-x^{\nu/10}.\label{dec:FM_growth}
\end{align}
Let $M=\lfloor(1-\gamma)^{-1}\rfloor$.
Suppose that for every integer $n\ge M+1$ there is an
$a\in\mathbb N$ with $a/n\in[\theta,\theta+\nu]$, and that
there is a positive integer $h$ such that
\[
 h(1-\gamma)\in[\theta,\theta+\nu]\cup[1-\theta-\nu,1-\theta],
\]
Then, for given $A_{\rm out}>1$ and $\varpi\ge1$, taking $B_*$
sufficiently large gives $\left|\sum_pw_p\right|\ll x\log^{-A_{\rm out}}x$.

We use the exact parameters
\begin{equation}\label{dec:FMparameters}
 \begin{aligned}
 \theta&=499999/2000000,&\qquad \gamma&=1500001/2000000,\\
 \nu&=348001/2000000,& \theta+\nu&=53/125.
 \end{aligned}
\end{equation}
First take the intermediate Type II interval
$[.24999925,.4242]$ inside \eqref{dec:rawwindow}.
By \cite[Proposition~4.11]{FordMaynard2024}, Type II supplies
maximal Type I on both factor ranges with a loss of one logarithm. The complementary range is
$[x^{.5758},(x/2)^{.75000075}]$.
It overlaps the range $[1,x^{.629}]$ of Lemma~\ref{dec:typeI},
and its upper endpoint strictly exceeds the required
$x^{.7500005}$, so \eqref{dec:FM_I} has no gap.
The final range in \eqref{dec:FM_II} also lies strictly inside the
input range.

The parameters satisfy $\gamma=1-\theta$ and $M=4$.
For $n=5$, take $a=2$.
For $n\ge6$, we have $n\nu>1$, so take $a=\lceil n\theta\rceil$.
For the other combinatorial condition, take $h=3$, since
$3(1-\gamma)=.7499985\in(.576,.7500005)$.

Fix the parameters in the following order, so that all estimates
apply to the same model.
First take $A_{\rm out}=3$ and $\varpi=2$, and let the sieve theorem
determine $B_*$. Choose integers $D\ge B_*+4$ and $K\ge D+4$,
then determine $J,L_*,K_*$ by the extension to divisor-bounded
coefficients above. Finally make a single choice of
\[
 C\ge100(K_*+10),\qquad C\alpha>K+2^{\lceil D\rceil}+3,
\]
and fix a common major-arc width. Put
\[
 a_n=(q/A)f_*(n),\qquad b_n=(q/A)G_C(n),\qquad w_n=a_n-b_n.
\]
These sequences are supported on $(x/2,x]$.
When $v=1$, the unit condition makes them zero at the left endpoint.
Using $\tau(n)\le2\sum_{r\mid n,r\le\sqrt x}1$, the digital
arithmetic-progression estimate already proved, and the Chinese
remainder theorem for the model, we obtain
\[
 \sum_n(f_*+G_C)\tau(n)\ll_{v,C}A\log q.
\]
Thus the total growth condition in \eqref{dec:FM_growth} holds.
For the pointwise negative part, use the actual model value:
\[
 b_n\le(q/A)\rho=B_0^{2-2\alpha}=(\log q)^{O_C(1)}\le x^{\nu/10}
\]
for sufficiently large $q$.
Proposition~\ref{dec:commonmodel} applies to this fixed $C$.
After multiplication by $q/A$, its Type I and Type II bounds become
$O_v(q\log^{-K}q)$ with arbitrarily prescribed $K$.
The strict margins between the $q$-ranges and the displayed $x$-ranges
absorb the fixed factor $x/q=v+1$.
Taking the prescribed logarithmic precision gives
\eqref{dec:FM_I}--\eqref{dec:FM_II}. Hence
\begin{equation}\label{dec:primecomparison}
 \left|\sum_p(f_*(p)-G_C(p))\right|\ll_vA\log^{-3}q.
\end{equation}

\subsection{The leading constant and arbitrary real endpoints}
\label{dec:endpoints}

Each pair consisting of an admissible initial word and a unit terminal
word in the model corresponds to a physical interval of length $q/B_0$
and a unit residue class modulo $B_0$.
The Siegel--Walfisz theorem applies uniformly to these moduli, which are
bounded by a fixed power of $\log q$, and to all endpoints.
Since $v$ is fixed, $\log t=\log q+O_v(1)$, and therefore
\begin{align}
 \sum_pG_C(p)
 &\sim\rho(b-1)^Lh(b-1)^{L-1}
       \frac{q}{B_0\varphi(B_0)\log q}\notag\\
 &=\kappa_{b,d}\frac{(b-1)^k}{k\log b},\qquad
 \kappa_{b,d}=\frac{b\,h}{(b-1)\varphi(b)}.
 \label{dec:prefixasymptotic}
\end{align}
Once $C$ is fixed, the numbers of intervals and residue classes are
both bounded by fixed powers of $\log q$, and are absorbed by choosing
sufficiently strong logarithmic savings in the prime number theorem
for arithmetic progressions.
A prime that is not a unit must divide $b$, and contributes nothing
to a sufficiently large fixed-prefix interval.
Equation \eqref{dec:primecomparison} therefore proves
\eqref{dec:prefixasymptotic} for every fixed admissible positive prefix.

The fixed-prefix estimate has the same constant $\kappa_{b,d}$
for every admissible positive prefix $v$. Lemma~\ref{prelim:endpoints}
therefore gives the asserted asymptotic for ordinary expansions at every
real endpoint, completing the proof of
Theorem~\ref{dec:digital-criterion}. The allowed unit last digits give
\[
 \kappa_{10,d}=
 \begin{cases}
 5/6,&d\in\{1,3,7,9\},\\
 10/9,&d\in\{0,2,4,5,6,8\},
 \end{cases}
 \qquad
 \kappa_{11,d}=
 \begin{cases}
 11/10,&d=0,\\
 99/100,&d\in\{1,\ldots,10\}.
 \end{cases}
\]
In particular, Fourier reflection of the deleted digit preserves the
amplitude estimates but need not preserve the arithmetic constant:
in base eleven it interchanges exclusion of $0$ and exclusion of a
nonzero unit. Proposition~\ref{dec:verified-inputs} verifies the digital
hypotheses for every required digit, and hence proves
Theorem~\ref{main:exact}.

\section{Completion of the main results}
\label{sec:completion}

\begin{proof}[Proof of Theorems~\ref{main:exact} and \ref{main:order}]
Theorem~\ref{dec:digital-criterion} gives the exact prime asymptotic
from the three shifted moment bounds and the one-digit correlation
bound. Appendix~\ref{app:exact-inputs} verifies these hypotheses for
every omitted digit in bases ten and eleven. The local density and
fixed-prefix completion in Section~\ref{sec:preliminaries} give
the constant $\kappa_{b,d}$ and the assertion for arbitrary real
endpoints. This proves Theorem~\ref{main:exact}.

Theorems~\ref{octal:prime-count} and \ref{nine:theorem} prove the
two-sided estimates in terms of $x^{c_b}/\log x$.
Lemma~\ref{prelim:count} gives $A_{b,d}(x)\asymp x^{c_b}$.
\end{proof}

\section*{AI Disclosure}
The author selected the research topic; artificial intelligence completed all mathematical derivations, proofs, and writing.

\appendix
\addtocontents{toc}{\protect\setcounter{tocdepth}{1}}
\section{Continuous bounds from finite grids}\label{app:continuous}

The finite certificates must bound a transfer operator or a correlation
profile throughout its continuous state space. We establish the
interpolation estimates that control the values between grid points and
the curvature bounds needed to apply them. For transfer operators, the
resulting supersolution inequality combines with
Lemma~\ref{prelim:supersolution} to give a moment bound for every word length.

\subsection{Interpolation and transfer supersolutions}

\begin{lemma}[Semiconvex grid interpolation]\label{prelim:grid}
If a continuous function $u$ satisfies $u''\geqslant-C$ distributionally
on an interval $[a,a+\Delta]$, then
\[
 u(t)\leqslant
 \frac{a+\Delta-t}{\Delta}u(a)+
 \frac{t-a}{\Delta}u(a+\Delta)
 +\frac C2(t-a)(a+\Delta-t).
\]
In particular, its supremum is at most the larger endpoint value plus
$C\Delta^2/8$. On a rectangular grid, if the second derivative in coordinate
$j$ is bounded below by $-C_j$, the supremum is at most the largest vertex
value plus $\sum_jC_j\Delta_j^2/8$.
\end{lemma}
\begin{proof}
The function $u(t)+Ct^2/2$ is convex and hence lies below its endpoint
chord. Subtract $Ct^2/2$ and use
$(t-a)(a+\Delta-t)\leqslant\Delta^2/4$. Applying this argument successively
in each coordinate proves the rectangular version.
\end{proof}

All transfer supersolutions are periodic linear interpolants of positive
rational node values. A Lipschitz alternative is sometimes convenient:
if the residual $\Lambda h-T_ph$ has Lipschitz constant $L$ on each
cell of length $1/N$, a nodal margin $L/(2N)$ suffices.
For the sharper second-derivative criterion, Lemma~\ref{prelim:grid} gives
\begin{equation}\label{prelim:semiconvex}
 \Lambda h(j/N)-(T_ph)(j/N)\geqslant C/(8N^2)
 \quad(0\leqslant j<N)
 \quad\Longrightarrow\quad T_ph\leqslant\Lambda h,
\end{equation}
provided $(T_ph)''\geqslant-C$ on each open cell.

\begin{lemma}[A common transfer certificate]\label{prelim:operator-cert}
Let $p\geqslant1$ and
\[
 (T_ph)(x)=\sum_{i=0}^{b-1}|Z_i(x)|^p h((i+x)/b),
 \qquad |Z_i|\leqslant1,\quad |Z_i'|\leqslant a,\quad |Z_i''|\leqslant c.
\]
Suppose that $h$ is positive, periodic, linear on the $1/N$ grid,
$\|h\|_\infty\leqslant1$, and $b\mid N$. If
\[
 \Lambda h(j/N)-(T_ph)(j/N)
 \geqslant\frac{bpc+2pa\operatorname{Lip}(h)}{8N^2}
 \qquad(0\leqslant j<N),
\]
then $T_ph\leqslant\Lambda h$ everywhere.
\end{lemma}
\begin{proof}
Convexity of $z\mapsto|z|^p$ gives, distributionally,
\[
 (|Z_i|^p)''\geqslant-pc,\qquad
 \bigl|(|Z_i|^p)'\bigr|\leqslant pa.
\]
For completeness, the Hessian term in the chain rule is nonnegative,
while the remaining term is bounded below by
$-p|Z_i|^{p-1}|Z_i''|\geqslant-pc$. At zeros, use the convex
regularization $(|Z_i|^2+\epsilon)^{p/2}$ and pass to the limit.
Since $b\mid N$, no inverse branch crosses a breakpoint of $h$ inside
an input cell. The product rule therefore gives
\[
 (T_ph)''\geqslant-bpc-2pa\operatorname{Lip}(h)
\]
on that cell. Apply \eqref{prelim:semiconvex}.
\end{proof}

\subsection{Curvature of affine-correlation profiles}

The affine-correlation certificates require the same interpolation principle
in two variables. The following bounds remain valid at zeros of the
one-digit amplitudes.

\begin{lemma}[Curvature of Hellinger profiles]\label{prelim:hellinger-curvature}
Let $A(x)=(A_0(x),\ldots,A_{b-1}(x))$ be a $C^2$ complex vector with
\[
 \|A(x)\|_2=1,\qquad \|A'(x)\|_2\leqslant a,\qquad
 \|A''(x)\|_2\leqslant c.
\]
Put $P_i(x)=|A_i(x)|^2$,
\[
 H(x,E)=\sum_i|A_i(x)||A_i(x+E)|,\qquad
 \rho(E)=\sup_{x\in[0,1]}H(x,E).
\]
For arbitrary fixed shifts $s_i$, define
$G(x,t)=\sum_iP_i(x)\rho(t+s_i)^2$. Then
$0\leqslant H,\rho,G\leqslant1$, and
\[
 H_{xx}\geqslant-2c-2a^2,\qquad H_{EE}\geqslant-c,\qquad
 (\rho^2)''\geqslant-2c,\qquad
 G_{xx}\geqslant-2c,\quad G_{tt}\geqslant-2c
\]
in the sense of distributions.
\end{lemma}
\begin{proof}
Cauchy--Schwarz gives $H\leqslant1$, and the other range bounds follow.
To justify differentiation at zeros, set
$a_{i,\epsilon}=(|A_i|^2+\epsilon^2)^{1/2}$. Direct differentiation gives
\[
 |a_{i,\epsilon}'|\leqslant|A_i'|,
 \qquad a_{i,\epsilon}''\geqslant-|A_i''|,
 \qquad \|(a_{i,\epsilon})_i\|_2=(1+b\epsilon^2)^{1/2}.
\]
The product rule and Cauchy--Schwarz therefore bound the second
$x$ derivative of the regularized $H$ below by
$-2c(1+b\epsilon^2)^{1/2}-2a^2$, and its second $E$ derivative below by
$-c(1+b\epsilon^2)^{1/2}$. Let $\epsilon\downarrow0$.

For fixed $x$, squaring the nonnegative function $H(x,E)\leqslant1$
gives $(H(x,\cdot)^2)''\geqslant-2c$. Suprema preserve a common
semiconvexity bound, because adding $cE^2$ makes every such function
convex. Since $\rho^2=\sup_x H(x,\cdot)^2$, this proves the bound for
$\rho^2$, and averaging with the probabilities $P_i(x)$ gives $G_{tt}$.
Finally,
\[
 G_{xx}
 =2\sum_i\rho(t+s_i)^2
       \bigl(|A_i'|^2+\operatorname{Re}(A_i''\overline{A_i})\bigr)
 \geqslant-2\|A''\|_2\|A\|_2\geqslant-2c.
\]
\end{proof}

\subsection{Rational enclosures}

Trigonometric enclosures use Machin's identity
\[
 \pi=16\arctan(1/5)-4\arctan(1/239),
\]
alternating-series bounds, and Taylor bounds for cosine with directed
rounding. A rational-power bound $(a/D)^{u/v}\leqslant w/D$ is certified
by the integer comparison
\begin{equation}\label{prelim:integer-power}
 w^vD^{u-v}\geqslant a^u\qquad(u\geqslant v).
\end{equation}
The parameters and strict margins for the moment and affine-correlation
bounds in bases ten and eleven are given in Appendix~\ref{app:exact-inputs}.
The corresponding moment certificates in bases eight and nine are given in
Appendix~\ref{app:small-certificates}.


\section{Unit-circle zeros and one-sided shifts}\label{app:shifts}

The shift comparison used in the carry decomposition depends on the
arithmetic location of the zeros of the digit polynomial. Their orders
are coprime to the base, so an unshifted digital argument stays a controlled
distance from every zero. We first classify these zeros and then compare
the full digital products, retaining the multiplicity of each zero.

\begin{lemma}[Zeros of a one-digit-deleted polynomial]\label{prelim:zeros}
For $b\geqslant3$ and $0\leqslant d<b$, the unit-circle zeros of
$P_{b,d}(z)=\sum_{j=0}^{b-1}z^j-z^d$ are exactly the points $z\ne1$
satisfying
\begin{equation}\label{prelim:zero-classification}
 \bigl(z^{b-1}=1,\ z^d=1\bigr)
 \quad\hbox{or}\quad
 \bigl(z^{b+1}=1,\ z^{d+1}=-1\bigr).
\end{equation}
Their orders are coprime to $b$. Every such zero is simple, except that
$z=-1$ is a double zero when $b\equiv1\pmod4$ and $d=(b-1)/2$.
\end{lemma}
\begin{proof}
Since $P_{b,d}(1)=b-1$, consider $z=e^{i\theta}\ne1$.
If $P_{b,d}(z)=0$, the geometric sum has modulus one, so
\[
 \left|\frac{\sin(b\theta/2)}{\sin(\theta/2)}\right|=1.
\]
Thus $\cos(b\theta)=\cos\theta$, which gives $z^{b-1}=1$ or
$z^{b+1}=1$. The geometric sum equals $1$ in the first case and
$-z^{-1}$ in the second; this proves both necessity and sufficiency
of \eqref{prelim:zero-classification}.

At a zero of the first type and of the second type, respectively,
\[
 zP_{b,d}'(z)=-\frac{(b-1)z}{1-z}-d,\qquad
 zP_{b,d}'(z)=z^{-1}\frac{(d-b)-(d+1)z}{1-z}.
\]
Either derivative can vanish on the unit circle only when
$d=(b-1)/2$ and $z=-1$. This point is a zero precisely when $d$ is
even, equivalently $b\equiv1\pmod4$. Writing $b=2d+1$ in that case,
\[
 P_{b,d}''(-1)=2d^2-d(d-1)=d(d+1)>0,
\]
so its multiplicity is exactly two. Finally, the orders in
\eqref{prelim:zero-classification} divide $b-1$ or $b+1$.
\end{proof}

\begin{lemma}[One-sided shift comparison]\label{prelim:shift}
Fix $b,d$ and $K>0$. Let $q$ and $Q=b^L\leqslant q$ be powers of $b$,
let $A\in\mathbb Z$, and let $|\delta|\leqslant1$. If
$F_Q((A+\delta)/Q)\geqslant q^{-K}$, then
\begin{equation}\label{prelim:shift-bound}
 F_Q((A+\delta)/Q)
 \leqslant \exp\!\bigl(O_{b,d,K}(\sqrt{\log q})\bigr)F_Q(A/Q).
\end{equation}
\end{lemma}
\begin{proof}
For $1\leqslant j\leqslant L$, write
$u_j=f_{b,d}(A/b^j)$ and $v_j=f_{b,d}((A+\delta)/b^j)$.
We first prove the uniform single-factor estimates
\begin{equation}\label{prelim:factor-comparison}
 u_j>0,\qquad v_j/u_j\ll_{b,d}1,\qquad |v_j-u_j|\ll_{b,d}b^{-j}.
\end{equation}
The last bound follows from the Lipschitz property of $f_{b,d}$.
For the first two, factor the polynomial as
\[
 P_{b,d}(z)=R(z)\prod_{\zeta\in\mathcal Z}(z-\zeta)^{m_\zeta},
\]
where $\mathcal Z$ is the set of unit-circle zeros and $R$ has none
on that circle. Thus $|R|$ is bounded above and away from zero there.
Write $\zeta=\mathrm e(a/r)$ in lowest terms. Lemma~\ref{prelim:zeros}
gives $(r,b)=1$, so
\[
 \|A/b^j-a/r\|\geqslant (rb^j)^{-1}.
\]
Consequently the denominator in
\[
 \frac{|\mathrm e((A+\delta)/b^j)-\zeta|}
      {|\mathrm e(A/b^j)-\zeta|}
 \leqslant
 1+\frac{2\pi b^{-j}}{4/(rb^j)}
\]
is nonzero and the ratio is uniformly bounded. Taking the product,
with the actual multiplicities $m_\zeta$, proves
\eqref{prelim:factor-comparison}.

At positions with $u_j>b^{-j/2}$, the ratio is
$1+O_{b,d}(b^{-j/2})$, whose product is bounded. At the other positions,
$v_j\ll_{b,d}b^{-j/2}$, and hence $v_j\leqslant b^{-j/4}$ once $j$
exceeds a fixed constant. Since $\prod_jv_j\geqslant q^{-K}$ and
$v_j\leqslant1$ at every position,
\[
 \sum_{\substack{j\text{ exceptional}\\j\text{ sufficiently large}}}
 j\leqslant4K\log_bq.
\]
There are therefore $O_{b,d,K}(\sqrt{\log q})$ exceptional positions.
Each costs a bounded ratio by \eqref{prelim:factor-comparison}, proving
\eqref{prelim:shift-bound}.
\end{proof}


\section{Verification of the exact-asymptotic digital inputs}
\label{app:exact-inputs}\label{sec:eleven}

This appendix proves Proposition~\ref{dec:verified-inputs}.
The moment exponents at orders $1$ and $235/154$ in base ten come
from Maynard \cite[Lemmas~10.3--10.4]{Maynard2019}. Karwatowski
\cite[Sections~3 and~5]{Karwatowski2022} gives the corresponding
eigenvalue criteria and reports their verification in base eleven. We give rational certificates for both bases, including
the quantitative $3/2$-moments and the affine contraction used here.
Karwatowski also uses a $3/2$-moment for the base-nine endpoint digits
\cite[Section~3, p.~5]{KarwatowskiBaseNine}. The continuous
certification principles are proved in Appendix~\ref{app:continuous};
the exact rational data and their verification are specified in
Appendix~\ref{app:reproduce}.

\subsection{Uniform shifted moments}

To make the passage from a finite certificate to arbitrary word length explicit,
we record the common certification principle. For $p\ge1$, define
\[
 (T_ph)(x)=\sum_{i=0}^{b-1}f((i+x)/b)^p h((i+x)/b).
\]
Expansion one digit at a time gives the exact identity
\[
 T_p^k1(x)=\sum_{a<b^k}F_{b^k}((a+x)/b^k)^p.
\]
If a positive periodic piecewise linear function $h$ satisfies
$T_ph\le\Lambda h$, iteration of the positive operator gives
\begin{equation}\label{dec:transferiterate}
 T_p^k1(x)\le \frac{\max h}{\min h}\Lambda^k.
\end{equation}
The passage from rational nodal inequalities to a continuous supersolution
uses Lemmas~\ref{prelim:grid} and \ref{prelim:operator-cert}.
Their distributional formulation includes zeros of the digit polynomial.
Thus a finite certificate proves a uniform statement for every word length;
the prefactor in \eqref{dec:transferiterate} is independent of $k$.

\subsubsection{Base ten}

\begin{lemma}\label{dec:decimalcert}
For $b=10$, estimates \eqref{dec:moment1}--\eqref{dec:momentnew} hold.
\end{lemma}
\begin{proof}
Put
\[
 Z_i(x)=\frac19\sum_{\substack{0\le j<10\\j\ne d}}
                         e((j-9/2)(i+x)/10).
\]
The finite sums and the inequality $2\pi<44/7$ give
\[
 |Z_i|\le1,\quad |Z_i'|\le110/63,\quad |Z_i''|\le2662/735.
\]
Let $N$ be divisible by $10$, and let $h$ be the periodic linear
interpolant of the certified nodal vector, normalized by
$\|h\|_\infty\le1$, with Lipschitz constant $L_h$.
Lemma~\ref{prelim:operator-cert}, with
$a=110/63$ and $c=2662/735$, shows that it suffices to choose
\[
 c_0\ge10p(2662/735),\qquad c_1\ge2p(110/63),
\]
and verify, at every node $j/N$, that
\begin{equation}\label{dec:nodecert}
 \Lambda h(j/N)-T_ph(j/N)>\frac{c_0+c_1L_h}{8N^2}.
\end{equation}
The rational parameters used are
\[
\begin{array}{c|c|c|c|c}
 p&\text{exponent}&\Lambda&N&(c_0,c_1)\\ \hline
 1&27/77&1121/500&1000&(37,4)\\
 235/154&59/433&1368543/1000000&1000&(56,6)\\
 3/2, d\notin\{4,5\}&287/2000&13915537/10000000&1000&(55,6)\\
 3/2, d\in\{4,5\}&14387/100000&6963699/5000000&2000&(55,6)
\end{array}
\]
The rational nodal vectors and their exact verification are included in
the supplementary material described in Appendix~\ref{app:reproduce}.
Trigonometric values are enclosed using a rational interval for $\pi$ from
Machin's identity and Taylor polynomials with rigorous remainder bounds.
The identity
\[
 f(t)^2=\frac{9+2\sum_{r=1}^9c_r\cos(2\pi rt)}{81},\qquad
 c_r=\#\{j:j,j+r\ne d, 0\le j<j+r\le9\},
\]
allows each nodal value to be rounded upwards to a rational with denominator
$D=2^{40}$. If $p=u/v$ and $f\le a/D$, the integer condition
$w^vD^{u-v}\ge a^u$ ensures that $f^p\le w/D$.
After denominators are cleared, every margin in \eqref{dec:nodecert} is
a positive integer. The inequalities $\Lambda\le10^{\text{exponent}}$
are also checked by comparing integer powers. The five representative digits
$d=0,\ldots,4$, together with $f_{9-d}=f_d$, cover all digits.
The conclusion follows from \eqref{dec:transferiterate}.
\end{proof}

\subsubsection{Base eleven}

Put $D_d=\{0,\ldots,10\}\setminus\{d\}$ and $f_d=f_{11,d}$.
The reflection
$f_{10-d}=f_d$ reduces the verification to $d=0,\ldots,5$.

\begin{lemma}[Uniform moments in base eleven]\label{eleven:moments}
For every $d\in\{0,\ldots,10\}$, every $q=11^k$, and every real $x$,
\begin{align}
 \sum_{a<q}F_q((a+x)/q)&\ll q^{27/77},\label{eleven:l1}\\
 \sum_{a<q}F_q((a+x)/q)^{235/154}&\ll q^{59/433},\label{eleven:oldmoment}\\
 \sum_{a<q}F_q((a+x)/q)^{3/2}&\ll q^{287/2000}.\label{eleven:newmoment}
\end{align}
The implied constants are independent of $k,x$.
\end{lemma}
\begin{proof}
Each certificate uses $N=1100$ nodes and common denominator $D=2^{40}$,
specifying
\[
 h(j/N)=v_j/D,\qquad 0<v_j\leqslant D,
\]
with periodic linear interpolation between nodes. The positive vectors
$v$ and the integer arrays of nodal weights are given in the
supplementary file \path{supplement/base11/moments.json};
see Appendix~\ref{app:reproduce}. The three sets of parameters are
\begin{center}
\begin{tabular}{ccccc}
\toprule
$p$ & $g$ & $\Lambda$ & $c_0$ & $c_1$\\
\midrule
$1$ & $27/77$ & $23/10$ & $41$ & $4$\\
$235/154$ & $59/433$ & $137/100$ & $62$ & $6$\\
$3/2$ & $287/2000$ & $7/5$ & $62$ & $6$\\
\bottomrule
\end{tabular}
\end{center}
In each row, $\Lambda\leqslant11^g$ is verified by an integer-power
comparison.

To show that these data control the continuous state space, let
\[
 Z_i(x)=\frac1{10}\sum_{a\in D_d}
          \mathrm e((a-5)(i+x)/11).
\]
Using $\sum_{a=0}^{10}|a-5|=30$ and
$\sum_{a=0}^{10}(a-5)^2=110$, we obtain
\[
 |Z_i|\leqslant1,\qquad
 |Z_i'|\leqslant\frac{2\pi\cdot30}{110}\leqslant\frac{110}{63},
 \qquad
 |Z_i''|\leqslant\frac{4\pi^2\cdot110}{1210}
       \leqslant\frac{2662}{735}.
\]
Let $L_h=\operatorname{Lip}(h)$. The constants in the table satisfy
$c_0\geqslant11p(2662/735)$ and $c_1\geqslant2p(110/63)$.
Since $11\mid N$, Lemma~\ref{prelim:operator-cert} shows that it is enough
to check at every node that
\begin{equation}\label{eleven:node}
 \Lambda h(j/N)-(T_ph)(j/N)
 >\frac{c_0+c_1L_h}{8N^2}.
\end{equation}

Upper bounds for the nodal weights are computed directly from
\[
 f_d(t)^2=\frac{10+2\sum_{r=1}^{10}c_r\cos(2\pi rt)}{100},
 \qquad c_r=\#\{a\in D_d:a+r\in D_d\}.
\]
The program \path{exact_trig.py} uses fixed-point denominator $2^{90}$,
Machin's formula, the degree-$40$ Taylor polynomial, and remainder bound
$(22/7)^{42}/42!$ to produce directed cosine intervals. An upward-rounded
square root bounds the amplitude, and \eqref{prelim:integer-power}
bounds its rational $p$th power. These weights also have denominator $D$.
If $m=\max_j|v_{j+1}-v_j|$ and $\Lambda=u/\ell$ with coprime positive integers
$u,\ell$, clearing denominators in \eqref{eleven:node} gives the integer
inequality
\[
 \begin{split}
 \mathcal M_j={}&8N^2\left[
  11Du\,v_j-\ell\sum_{i=0}^{10}w_{iN+j}
       \bigl((11-r)v_{iN/11+\lfloor j/11\rfloor}
               +r v_{iN/11+\lfloor j/11\rfloor+1}\bigr)\right]\\
 &-11D\ell(c_0D+c_1Nm)>0,\qquad r=j\bmod11.
 \end{split}
\]
The subscripts of $v_j$ are taken modulo $N$. The index $iN+j$ of
the directed weight $w_{iN+j}$ runs over the full $11N$-point grid. For each moment, the smallest integer
margin over the six digit representatives and all nodes is
\begin{center}
\begin{tabular}{cl}
\toprule
$p$ & $\min_{d,j}\mathcal M_j$\\
\midrule
$1$ & $39943001465673152637111041290240$\\
$235/154$ & $191058315489861615435892981836800$\\
$3/2$ & $12282688731347493451960444195840$\\
\bottomrule
\end{tabular}
\end{center}
The program \path{supplement/base11/verify_moments.py} computes the trigonometric bounds
from the digit set and checks every weight, nodal interpolation, and
margin above using integer and rational arithmetic. Thus
\eqref{prelim:semiconvex} gives $T_ph\leqslant\Lambda h$, and
Lemma~\ref{prelim:supersolution} proves the bounds for every $k,x$.
\end{proof}

\subsection{One-digit affine contraction}

We use the quantities $P_i$, $H$, $\rho$ and $K_2$ defined in
Section~\ref{dec:inputs}. The path estimate itself was proved in
Lemma~\ref{dec:hellinger}; it remains to bound its one-digit coefficient.

\subsubsection{Base ten}

\begin{lemma}[One-digit affine contraction in base ten]
\label{dec:decimalhellinger}
For base ten, the following upper bounds for $K_2$ and admissible values
of $\delta$ satisfy $K_2\le10^{-3\delta}$:
\[
\begin{array}{c|ccccc}
 \text{representative }d&0&1&2&3&4\\ \hline
 \text{upper bound for }K_2&19/20&967/1000&239/250&481/500&241/250\\
 \text{admissible }\delta&7/1000&3/625&3/500&11/2000&53/10000
\end{array}
\]
The reflected digit $9-d$ has the same bounds. In particular,
\eqref{dec:groups} is valid for every decimal digit.
\end{lemma}
\begin{proof}
Use the centered amplitudes
\[
 A_i(x)=90^{-1/2}\sum_{\substack{0\le j<10\\j\ne d}}
             e((j-9/2)(i+x)/10),\qquad 0\le i<10.
\]
Discrete Parseval gives $\|A\|_2=1$ and
\[
 \|A'\|_2^2\le(2\pi/10)^2\frac{55}{6}<4,\qquad
 \|A''\|_2^2\le(2\pi/10)^4\frac{3223}{24}<25.
\]
Lemma~\ref{prelim:hellinger-curvature} therefore gives
$H_{xx}\ge-18$, $H_{EE}\ge-5$ and $(\rho^2)''\ge-10$.
For $G_a(x,t)=\sum_iP_i(x)\rho(ai+t)^2$, both coordinate second
derivatives are bounded below by $-10$.

Take $M=200$. At each $E=k/M$, maximize the directed upper bounds
for $H(j/M,E)$ and add $18/(8M^2)$. This bounds the continuous
supremum defining $\rho(E)$ by Lemma~\ref{prelim:grid}.
Evaluate $G_a$ on the product grid in $(x,t)$ and add $28/(8M^2)$;
this exceeds the required total curvature correction $20/(8M^2)$.
The same rational trigonometric enclosures used in
Lemma~\ref{dec:decimalcert}, with upward-rounded probabilities and
square roots of denominator $2^{28}$, give the displayed bounds.
The computation covers $0\le x\le1$, $0\le t\le10$ and
$a=1,\ldots,5$; evenness and ten-periodicity of $\rho$ cover the
remaining nonzero residue classes.
All these finite operations and their directed roundings are specified
in the supplementary verification program.
Finally, if an entry is $K_2\le u/v$ and $\delta=r/s$, the comparison
with $10^{-3\delta}$ is checked by the integer inequality
$u^s10^{3r}\le v^s$. Each of the five pairs in the table satisfies it.
\end{proof}

\subsubsection{Base eleven}

\begin{lemma}[Affine contraction in base eleven]\label{eleven:hellinger}
For every one of the eleven forbidden digits,
\begin{equation}\label{eleven:K2}
 K_2\leqslant\frac{24}{25}\leqslant11^{-3/250}.
\end{equation}
Consequently, in the notation of Lemma~\ref{dec:hellinger},
\[
 \mathcal H_L\leqslant(B/11^m)^{-1/250},\qquad B=11^L,
\]
uniformly in the initial states, intercepts, and unit slopes satisfying
$\lambda\equiv1\pmod {11}$.
\end{lemma}
\begin{proof}
Let
$A_i(x)=110^{-1/2}\sum_{a\in D_d}\mathrm e((a-5)(i+x)/11)$.
Discrete Parseval and the centered exponents give
\[
 \|A\|_2=1,\qquad
 \|A'\|_2^2\leqslant(2\pi/11)^2\frac{110}{10}<4,
 \qquad
 \|A''\|_2^2\leqslant(2\pi/11)^4\frac{1958}{10}<25.
\]
Here $1958=\sum_{a=0}^{10}(a-5)^4$. Apply
Lemma~\ref{prelim:hellinger-curvature} with $a=2$, $c=5$. It gives
\[
 H_{xx}\geqslant-18,\qquad H_{EE}\geqslant-5.
\]
For $G_a(x,t)=\sum_iP_i(x)\rho(ai+t)^2$, the same lemma gives
\[
 (G_a)_{xx}\geqslant-10,\qquad (G_a)_{tt}\geqslant-10.
\]

Use a grid with $M=80$ and upward-rounded probabilities and amplitudes
with denominator $D_0=2^{28}$. First maximize over $x=j/M$ and add
$18/(8M^2)$ to bound $\rho(E)$ at each $E=k/M$. Next evaluate $G_a$ on the
full product grid in $(x,t)$ and add $20/(8M^2)$, as justified by
Lemma~\ref{prelim:grid}. Since $\rho$ is even
and $11$-periodic, it suffices to take $a=1,\ldots,5$; this reduction covers
every slope class. The program \path{supplement/base11/verify_hellinger.py}
computes these directed trigonometric intervals, square roots, and finite
suprema. The resulting upper bounds,
rounded further upward to denominator $10^4$, are
\begin{center}
\begin{tabular}{ccccccc}
\toprule
Digit representative $d$ & $0$ & $1$ & $2$ & $3$ & $4$ & $5$\\
\midrule
Upper bound for $K_2$ & $.9284$ & $.9378$ & $.9421$ & $.9476$ & $.9465$ & $.9373$\\
\bottomrule
\end{tabular}
\end{center}
Every entry is less than $24/25$. The final power comparison is exactly
$24^{250}11^3\leqslant25^{250}$. The two-ended path H\"older inequality is
proved in Lemma~\ref{dec:hellinger}; substituting $t=2$ gives the asserted
correlation contraction.
\end{proof}

\subsection{Scalar comparisons and the shift input}

For base ten, the inequalities
\[
 9^{500}>10^{477},\qquad 9^{100}<10^{99},\qquad
 9^{21}>10^{20},\qquad 730^{5000}<10^3\,729^{5000}
\]
give $1/100<1-\alpha<23/500$, $\alpha>20/21$, and
$\log_{10}(1+9^{-3})<3/5000$.
For base eleven, the corresponding integer comparisons are
\[
 10^{500}>11^{477},\qquad 10^{100}<11^{99},\qquad
 10^{21}>11^{20},\qquad 1001^{5000}<11^3\,1000^{5000}.
\]
Thus in this base as well,
\[
 1/100<1-\alpha<23/500,\qquad \alpha/2>10/21,\qquad
 \log_{11}(1+10^{-3})<3/5000.
\]
The even-moment identities \eqref{prelim:even} consequently give
\eqref{dec:evenmoments} in both bases.
The comparison needed for maximal Type I follows from
$\alpha>477/500$ and
\[
 \frac{(50/77)\alpha}{\alpha+(50/77)(1-\alpha)}
 >\frac{23850}{37879}>\frac{629}{1000}.
\]
The programs \path{supplement/base10/verify_base10.py} and
\path{supplement/base11/verify_base11.py} check these exact comparisons.

The one-sided shift estimate is a structural property of the deleted-digit
polynomial, proved in Appendix~\ref{app:shifts}. Its base-eleven form is
recorded here to specify this input explicitly.

\begin{corollary}[Base-eleven zeros and shift comparison]
\label{eleven:zeros}\label{eleven:shift}
The unit-circle zeros of $P_{11,d}$ are precisely
\[
 z^{10}=1,\ z^d=1,\ z\ne1;
 \qquad
 z^{12}=1,\ z^{d+1}=-1,\ z\ne1.
\]
They are all simple and have order coprime to $11$. For every fixed $K>0$,
if $Q\leqslant q$ are powers of $11$, $A\in\mathbb Z$,
$|\delta|\leqslant1$, and $F_Q((A+\delta)/Q)\geqslant q^{-K}$, then
\begin{equation}\label{eleven:distortion}
 F_Q((A+\delta)/Q)
 \leqslant\exp\!\bigl(O_{d,K}(\sqrt{\log q})\bigr)F_Q(A/Q).
\end{equation}
\end{corollary}
\begin{proof}
Apply Lemmas~\ref{prelim:zeros} and \ref{prelim:shift} with $b=11$.
The exceptional double-zero case cannot occur, since $11\equiv3\pmod4$.
\end{proof}

\begin{proof}[Proof of Proposition~\ref{dec:verified-inputs}]
Lemmas~\ref{dec:decimalcert} and \ref{eleven:moments} prove the three
uniform shifted moment estimates. Lemma~\ref{dec:decimalhellinger}
gives the required value of $\delta$ for the decimal central digits
and a value at least $1/250$ for every other decimal digit.
Lemma~\ref{eleven:hellinger} gives $\delta=1/250$ for every base-eleven
digit. These are precisely the coupled inputs in \eqref{dec:groups}.
The scalar comparisons above establish the additional elementary
inequalities used in the common argument.
\end{proof}

\section{The base-nine Fourier certificates}
\label{app:small-certificates}

The certificates provide positive supersolutions for the two transfer
operators. Interpolation extends the nodal inequalities to every
point of the circle. Appendix~\ref{app:reproduce} gives the exact
verification commands.

\subsection{Continuous Fourier moments in base nine}
\label{app:nine-moments}

Use the notation and parameters of Proposition~\ref{nine:fourier}.

\begin{proof}[Proof of Proposition~\ref{nine:fourier}]
Let
$f_d(t)=8^{-1}|\sum_{a\ne d,\,0\le a\le8}\mathrm e(at)|$,
and define a positive operator on continuous functions on the circle by
\[
 (\mathcal T_{d,s}h)(x)=\frac19\sum_{j=0}^8
 f_d\left(\frac{x+j}{9}\right)^s h\left(\frac{x+j}{9}\right).
\]
Expanding the inverse branches successively gives the exact identity
\begin{equation}\label{nine:transfer}
 \mathcal T_{d,s}^k1(x)=9^{-k}
 \sum_{a<9^k}F_{9^k}\left(\frac{a+x}{9^k}\right)^s.
\end{equation}
If $h>0$ and $\mathcal T_{d,s}h\le\lambda h$, then
$\mathcal T_{d,s}^k1\le\lambda^kh/\min h$.
Integration over $x\in[0,1]$ also gives the continuous moment.
Thus one fixed continuous supersolution controls every word length and every shift.

The base-nine certificate data in the mathematical supplement specify
ten positive rational vectors, one for each $d=0,\ldots,4$ and
$s=1,3/2$. Each defines a periodic piecewise linear function through
$8192$ nodes with denominator $10^9$.
The arrays are recorded in the file \texttt{moment\_certificates.json}
in \texttt{supplement/base9}.
The common target constants are
\begin{equation}\label{nine:lambdas}
 \lambda_1=\frac{2477}{10000},\qquad
 \lambda_{3/2}=\frac{314229}{2000000}.
\end{equation}
The following rational estimates certify the supersolution inequality
on the whole circle.

The number of ordered pairs of allowed digits with difference $r$ is
$c_r(d)=9-r-\mathbf1_{d+r\le8}-\mathbf1_{d-r\ge0}\ge0$.
Consequently,
\begin{equation}\label{nine:square-factor}
 f_d(t)^2=\frac{8+2\sum_{r=1}^8c_r(d)\cos(2\pi rt)}{64}.
\end{equation}
Take the arithmetic scale $S=10^{24}$.
Machin's identity
$\pi=16\arctan(1/5)-4\arctan(1/239)$,
with $50$ and $15$ terms of the respective alternating series, gives
\[
 \frac{3141592653589793238462643}{S}<\pi<
 \frac{3141592653589793238462644}{S}.
\]
Use periodicity and evenness to reduce the angles to $[0,\pi]$.
Evaluate the degree-$48$ Taylor polynomial for cosine with outward integer
rounding at each step, and bound the remainder by $4^{50}/50!$.
The cosine coefficients in~\eqref{nine:square-factor} are nonnegative,
so this gives an upper rational enclosure $v/S\in[0,1]$ for $f_d(t)^2$.
The corresponding first- and $3/2$-moment factors are bounded above by
\[
 S^{-1}\left\lceil\sqrt{vS}\right\rceil,\qquad
 S^{-1}\left\lceil(v^3S)^{1/4}\right\rceil
\]
respectively. All roots are computed with integer square roots and checked
again by integer powers.

Let the verification grid have $n$ nodes, with node values $h_i/G$, and put
\[
 M_h=\max_i h_i/G,\qquad
 L_h=\frac nG\max_i|h_{i+1}-h_i|.
\]
At $x=i/n$, the interpolated value on the inverse branch
$t=(i+jn)/(9n)$ is exactly
$((9-r)h_v+rh_{v+1})/(9G)$, where $i+jn=9v+r$.
Multiply by the root enclosures above and divide by the transfer average's
factor $9$. This produces an upper bound for
$\mathcal T_{d,s}h(i/n)$ with denominator $81SG$.

Centering the digit phases at $4$ gives
$\operatorname{Lip}(f_d)\le5\pi<110/7$.
Hence the circular Lipschitz constant of the residual is at most
\begin{equation}\label{nine:lip}
 L_*=\frac{s(110/7)M_h+L_h}{9}+\lambda_sL_h.
\end{equation}
Using exact fractions, the verifier checks
\begin{equation}\label{nine:residual}
 \max_i\{\mathcal T_{d,s}h(i/n)-\lambda_sh(i/n)\}
       +\frac{L_*}{2n}<0.
\end{equation}
Nine certificates use $n=8192$.
For $d=3,s=3/2$, each segment of the same linear interpolant is subdivided
exactly into sixteen equal parts, giving $n=131072$.
All ten checks give strict inequalities.
Every real point is within $1/(2n)$ of a grid point, so
\eqref{nine:residual} proves the supersolution inequality on the whole circle.
Every computation uses only fixed integer arrays and integer or rational arithmetic.

Finally, the verifier checks the integer inequalities
\[
 2477^{200}9^{127}<10000^{200},\qquad
 314229^{40000}9^{33693}<2000000^{40000}.
\]
Together with~\eqref{nine:transfer}, these give the first-moment discrete
exponent $73/200$ and the $3/2$-moment discrete exponent $6307/40000$.
The reflection identity $f_{8-d}=f_d$ covers the other four digits.
Applying H\"older's inequality on the same grid,
\[
 \sum F^{4/3}\le\left(\sum F\right)^{1/3}
                     \left(\sum F^{3/2}\right)^{2/3},
\]
gives $g_{4/3}=((73/200)+2(6307/40000))/3=13607/60000$.
\end{proof}

\section{Computational certificates}\label{app:reproduce}

\subsection{Exact data and acceptance criteria}

The electronic supplementary material contains the fixed rational data and
verification programs for the finite inequalities used in the proof.
The data consist of positive vectors, nodal upper bounds, and nonnegative
coefficients for dual and Farkas identities. Each acceptance criterion uses
only integer or rational arithmetic. Appendix~\ref{app:continuous} establishes the continuous certification
criteria. Their applications to the digital inputs are proved in
Appendices~\ref{app:exact-inputs} and~\ref{app:small-certificates}; the
exponent systems and sieve integrals are derived in the corresponding
analytic sections.

\begin{center}
\begin{tabular}{>{\raggedright\arraybackslash}p{.12\textwidth}
 >{\raggedright\arraybackslash}p{.40\textwidth}
 >{\raggedright\arraybackslash}p{.38\textwidth}}
\toprule
Base & Certificate objects & Result verified\\
\midrule
$8$ & Two shifted moment certificates and rational exponent inequalities
 & Odd-divisor and grouped comparisons for all eight digits\\
$9$ & Ten continuous moments; directed upper sums on full prime-factor simplices
 & Total sieve loss at most $.923405344688153132$\\
$10$ & Fifteen continuous moments; five correlation profiles; $188$ rational systems
 & Positive continuous residuals and strictly negative exponent bounds\\
$11$ & Eighteen continuous moments; six correlation profiles; scalar parameter comparisons
 & The moments in Lemma~\ref{eleven:moments} and $K_2\leqslant24/25$\\
\bottomrule
\end{tabular}
\end{center}

\paragraph{Continuous positive operators.}
A certificate specifies the integer numerators and a common denominator
of a positive periodic linear interpolant $h$.
Starting from the digit set, the verification program computes directed
trigonometric and rational-power enclosures. It then checks the positivity
of $h$, each nodal residual, and the continuous interpolation correction
from Appendix~\ref{app:continuous}. The comparison of the rational
supersolution constant with the claimed spectral growth is checked by
integer powers. These checks establish $T_ph\leqslant\Lambda h$
on the whole circle, and Lemma~\ref{prelim:supersolution} gives the
moment bound for every word length.

\paragraph{Rational systems of exponents.}
The forward and reverse certificates in \texttt{supplement/base10}
contain the exact systems and
their nonnegative dual or Farkas coefficients.
The verifier reconstructs every matrix from the parameters, constraints,
and branches in Sections~\ref{dec:forward}, \ref{dec:reverse},
and~\ref{dec:ranktwo} and checks
equality with the encoded system.
The six denominator coordinates \texttt{(a,b,c,d,e,f)} in the data
correspond, in that order, to $(a,c,h,D,E,G)$ in
Section~\ref{dec:ranktwo}; the coordinate order is part of the
certificate specification.
For a feasible system $Az\leqslant b$, a certificate for the target
$c^tz+c_0$ consists of $y\geqslant0$ satisfying
\[
 y^tA=c^t,\qquad y^tb+c_0<0.
\]
It follows that $c^tz+c_0\leqslant y^tb+c_0<0$ throughout the
real feasible polyhedron. An empty branch is certified by
$y\geqslant0$, $y^tA=0$, and $y^tb<0$.
Every matrix identity and sign comparison is exact.
The same $188$ systems apply in base $11$, whose digital inputs are
verified in Proposition~\ref{eleven:inputs}.

\paragraph{Sieve integrals.}
In base $9$, the programs enumerate the ordered simplex cells of
Proposition~\ref{nine:certificate}. Their volumes are given by the
Eulerian slice formula, with the factorial factors for equal grid
coordinates. A cell can be omitted if its upper multiplicity is zero
or if one subproduct is evaluable throughout the cell. Each remaining
cell retains the full marked
multiplicity, and every arithmetic rounding is upward.
These are the finite upper sums appearing in the sieve-loss bounds.

\subsection{Verification of the certificates}

The supplement contains the rational data and verification programs
for each base. From the directory containing this manuscript, run
\begin{verbatim}
python3 -B -S supplement/verify.py
\end{verbatim}
The programs use the Python standard library. The base-nine sieve
calculation also requires a C++17 compiler with support for
\texttt{unsigned \_\_int128}. The supplement's README lists the checks
and explains how to run them separately.

\clearpage
\phantomsection
\addcontentsline{toc}{section}{References}
\begin{thebibliography}{99}
\raggedright

\bibitem{ChenGuptaLi2026}
B. Chen, V. Gupta, and Y. C. Li,
\emph{Large value estimates for Dirichlet polynomials with characters and
zero density of Dirichlet $L$-functions},
\href{https://arxiv.org/abs/2507.08296v2}{arXiv:2507.08296v2}, 2026.

\bibitem{FordMaynard2024}
K. Ford and J. Maynard,
\textit{On the theory of prime-producing sieves},
preprint (2024),
\href{https://arxiv.org/abs/2407.14368v1}{arXiv:2407.14368v1}.

\bibitem{Harman2007}
G. Harman,
\textit{Prime-Detecting Sieves},
London Mathematical Society Monographs, vol.~33,
Princeton University Press, Princeton, NJ, 2007.

\bibitem{IwaniecKowalski2004}
H. Iwaniec and E. Kowalski,
\textit{Analytic Number Theory},
American Mathematical Society Colloquium Publications, vol.~53,
American Mathematical Society, Providence, RI, 2004.
\href{https://doi.org/10.1090/coll/053}{doi:10.1090/coll/053}.

\bibitem{Karwatowski2022}
F. Karwatowski,
\textit{Primes with one excluded digit},
Acta Arith. \textbf{202} (2022), 105--121.
\href{https://doi.org/10.4064/aa191002-26-8}{doi:10.4064/aa191002-26-8}.

\bibitem{KarwatowskiBaseNine}
F. Karwatowski,
\textit{Digits of primes in base $b=9$},
preprint, 17~pp.,
\href{https://www.math.hhu.de/fileadmin/redaktion/Fakultaeten/Mathematisch-Naturwissenschaftliche_Fakultaet/Mathematik/20_Institut-Lehrstuehle/5_Algebra_und_Zahlentheorie/Karwatowski/Digits_of_primes_in_base_b_9.pdf}{available from the author's webpage}.

\bibitem{Koukoulopoulos2019}
D. Koukoulopoulos,
\textit{The Distribution of Prime Numbers},
Graduate Studies in Mathematics, vol.~203,
American Mathematical Society, Providence, RI, 2019.
\href{https://dms.umontreal.ca/~koukoulo/documents/publications/primes.pdf}{Author's preliminary version}.

\bibitem{Maynard2019}
J. Maynard,
\textit{Primes with restricted digits},
Invent. Math. \textbf{217} (2019), 127--218.
\href{https://doi.org/10.1007/s00222-019-00865-6}{doi:10.1007/s00222-019-00865-6}.

\end{thebibliography}

\end{document}

